% Section: the noise response and the margin of the compensator (v3 §3).
% Author file; sub-files are \input below, including 04k-kernels.tex, 04x-x0band.tex.
\section{The noise response and the margin of the compensator}\label{sec:noise}
% v3: §3

This section proves that the predictable compensator of the jittered composites leaves a positive margin at every
scale. Write $L:=\log x$.
\begin{itemize}
\item Theorem~\ref{thm:g2} shows that the $H(x)$-window average $M$ of the margin satisfies
  $\log x\cdot M(x)\ge 0.97-o(1)-j_{X_0}(x)$ for $x\ge X_0^{1.01}$, and the same with $0.82$ in place of $0.97$ for
  $2X_0\le x<X_0^{1.01}$.
% Logged correction 47 (CORRECTIONS-vs-v3.md): piecewise base value, after Proposition g1.
\item Proposition~\ref{prop:pointwise} extends this to the pointwise margin density $m$.
\item Corollary~\ref{cor:r1} deduces the drift inequality R1$'$ on every window, with $c'=0.15$. This is the input of
  Theorem~\ref{thm:queue}.
\end{itemize}
The jump of the target at $X_0$ is handled exactly by Lemma~\ref{lem:xb}, using the short-interval Saias bound of
Lemma~\ref{lem:si}. The noise enters only through Saias-weighted dyadic sums \eqref{eq:sw-sums}, which depend on $X_0$
alone. This is what removes the divergence of the sup-norm noise response.

The section is organised as follows.
\begin{itemize}
\item \S\ref{sec:noise-notation}: notation.
\item \S\ref{sec:noise-primes}: prime inputs (Chebyshev--Mertens and smooth elements).
\item \S\ref{sec:noise-identity}: the margin identity and its push-through form.
\item \S\ref{sec:noise-saias}: the sharp coefficients (Lemma~M) and the Beurling--Saias bound.
\item \S\ref{sec:kernels}: the operator kernels.
\item \S\ref{sec:noise-perscale}: the noise, per scale.
\item \S\ref{sec:noise-lp}: a local prime number theorem.
\item \S\ref{sec:noise-bridge}: the bridge from the operator to the sharp reference, the $X_0$ jump, and positivity of
  the sharp reference.
\item \S\ref{sec:noise-margin}: the margin and R1$'$.
\end{itemize}

% 04a-setup.tex — v3 §3.0 (notation) and §3.1 (prime inputs).
\subsection{Notation and standing conventions}\label{sec:noise-notation}
% v3: §3.0

Throughout, $\NN$ is the realised multiplicative semigroup and $\PP$ the set of realised primes. Values are real
numbers, and elements with equal value are counted as distinct. The variable $x$ denotes the current scale,
$L_u:=\log u$, and dyadic scales are $u=X_0 2^j$. We write $\Lambda:=\log X_0\ge 10^8$ (\S\ref{sec:params}).

\paragraph{Jitter and slots} (\S\ref{sec:construction}).
\begin{itemize}
\item A prime $q$ decided at slot $s_q\ge X_0$ has its value uniformly distributed on the box $J_q$ of width $H(s_q)$,
  centred at $s_q$. Hence jittered values are $\ge \Xm:=X_0-H(X_0)/2$, and $\NN=\N$ below $\Xm$.
\item $H$ is the jitter width of \S\ref{sec:queue}: $H=\max(H_{\mathrm{std}},M_1^2)$ with
  $H_{\mathrm{std}}(u):=\exp(8(\log\log u)^2)$. It exceeds $H_{\mathrm{std}}$ only on the range
  $\log\log u\in[\approx 44,10^4]$ of Theorem~\ref{thm:caplw}.
\item $\sigma:=\sum_q\delta_{s_q}$ is the slot measure of the decided primes, and
  $\rhot:=J\sigma$, i.e.\ $\rhot(y)=\sum_q \mathrm{box}_{H(s_q)}(y-s_q)$, is the jitter-smoothed slot density.
\end{itemize}

\paragraph{Regularity of $H$.}
\begin{equation}\label{eq:RH}
  \frac{H(vw)}{H(v)}\le w^{\eps}\qquad(v\ge X_0,\ w\ge 1),\qquad
  \eps\le \frac{8(2\log\log X_0+0.65)}{\log X_0}\le 0.03,
\end{equation}
and $H(v)/v$ is decreasing. For $H_{\mathrm{std}}$, put $s:=\log w/\log v$; then
$\log H(vw)-\log H(v)=8\log(1+s)\,(2\log\log v+\log(1+s))$, and \eqref{eq:RH} follows from $\log(1+s)\le s$ and
$\log^2(1+s)/s\le 0.648$. On the raised range $d\log H/d\log u\le 4\cdot 10^{-6}$ (Theorem~\ref{thm:caplw}(c)), so the
same $\eps$ works there.

\paragraph{Decision order} (\S\ref{sec:construction}). $t(p)$ is the decision slot of $p$, so $|t(p)-p|\le H(p)$; for
rational primes $p<X_0$, $t(p)=p$. For $w\in\NN$ put $\tau_w:=\max_{p\mid w}t(p)$.

\paragraph{Operators on functions of $y\ge \Xm$.}
\begin{itemize}
\item The output-scale average: $(Bg)(y):=H(y)^{-1}\int_{W_y}g$, with $W_y:=[y-H(y)/2,\,y+H(y)/2]$.
\item The pointwise sharp latest-prime operator:
  \[
    (\Kpt g)(y):=\sum_{w\in\NN,\ w\ge 2}\frac1w\,\mathbf 1[\tau_w<y/w]\,\mathbf 1[y/w\ge \Xm]\,g(y/w).
  \]
  The cutoff $\Xm$ is the bottom of the support of $\rhot$; every function to which $\Kpt$ is applied ($T'$, $F$,
  $\rhot$ below) vanishes below $\Xm$.
\end{itemize}

\paragraph{The target} (\S\ref{sec:construction}).
\begin{itemize}
\item $T$ is the target \eqref{eq:target}. Its continuous source is $T'_{\mathrm{cont}}=1-\rhodet-\bar\rho_{\Abs}+R'$,
  with $\rhodet(u)=\Psi([u,(1+\kappa)u],X_0)/(\kappa u)$ \eqref{eq:rhodet} and $\kappa=X_0^{-5/12+2\eps}$. On
  $\Good(<u)$ the cap is inactive: $\bar\rho_{\Abs}=\rho_{\Abs}$ (Remark~\ref{rem:stopped}).
\item $R=\sum_k R_k$ is the designed oscillation. It vanishes below $x_1\ge\exp(X_0^2)$, has frequency
  $\gamma_k\in[\Gamma_k,2\Gamma_k]$, and $x_{k+1}\ge\exp(x_k)$.
\item The amplitude condition \eqref{eq:ampl} reads $\alpha_k\gamma_kx_k^{\theta-1}=\kappa_k'/\log x_k$, with
  $\kappa_k'=\kappa_0/\log^2x_k$ and $\kappa_0\le\frac14$. Hence
  \begin{equation}\label{eq:Rkprime}
    |R_k'(v)|\le\kappa_k\,(v/x_k)^{\theta-1},\qquad \kappa_k:=2\kappa_k'/\log x_k.
  \end{equation}
\item $\Abs$ is the absorbed set \eqref{eq:absorbed}. $K_{\Abs}$ denotes the part of $\Kpt$ whose latest prime is
  absorbed, and $\sigma_{\Abs}:=K_{\Abs}\rhot-\rho_{\Abs}$. By Lemma~\ref{lem:rhoA}(c), $|\sigma_{\Abs}|\le 2L^{-4}$, and
  $\sigma_{\Abs}$ vanishes unless $\log\log u>2\log\Lambda+5.97$ ($=42.8$ at $\Lambda=10^8$), since $\Abs$ is empty below
  that height.
\item \emph{The effective source} is
  \[
    T':=T'_{\mathrm{cont}}+K_{\Abs}\rhot=1-\rhodet+R'+\sigma_{\Abs}.
  \]
  The term $\sigma_{\Abs}$ is always kept in $T'$ (\S\ref{sec:noise-identity}).
\item \emph{The jump at $X_0$.} $T'=0$ below $X_0$ (there $T=\lfloor\cdot\rfloor$), and
  $T'(X_0^+)=\jz:=1-\rhodet(X_0)\le 1.0001/\log X_0$ by \eqref{eq:huxley}. Put
  \[
    J:=\jz\mathbf 1_{[X_0,\infty)},\qquad \Tc:=T'-J,
  \]
  so that $\Tc$ is continuous at $X_0$. Every oscillation bound below refers to $\Tc$. The jump is carried exactly, in
  chain form, by Lemma~\ref{lem:xb}.
\item $\tilde T'(y):=\#\{\text{continuous slots in }W_y\}/H(y)+(BK_{\Abs}\rhot)(y)$. Then
  \begin{equation}\label{eq:Ttilde}
    \tilde T'=BT'+\vartheta/H,\qquad|\vartheta|\le 6,
  \end{equation}
  where $\vartheta=\vartheta(y)$ collects the granularity: slots against $\int T'_{\mathrm{cont}}$, the two slot
  boundaries, and $\rhot$ against $\sigma(W_y)/H(y)$.
\end{itemize}

\paragraph{Sharp coefficients.} $c_n(y):=\mu(n)\chi^t_y(n)$ (Lemma~\ref{lem:lemmaM}), and
$\rhosh(y):=\sum_{n\in\NN}c_n(y)T'(y/n)/n$ is the pointwise sharp reference.

\paragraph{Constants.}
\begin{itemize}
\item The margin target is $c':=0.15$, the drift constant of R1$'$ (\S\ref{sec:queue}).
\item $M_1(u)$, $L_0(u)$, $y(u)$ and $B(u)$ are the per-scale parameters \eqref{eq:params}, with the tolerance
  $M_1\le H/L^{14}$. By \eqref{eq:queueB}, $B(u)\le 28M_1(u)L_u^2$ for $\log u\ge 10^8$.
\item $K_M=1.01$ and $K'=2.20$ (Lemma~\ref{lem:cm}); $C_0=0.70$ (Lemma~\ref{lem:smooth-v}).
\item $B_1:=1$, so that $\sum_{n\in\NN,\,n\le y}1/n\le\log y+B_1$ for every $y\ge 1$. For $\N$ the harmonic sum is
  $\le\log y+1$; for $y\ge\Xm$ the $\N$-excess is $\le\gamma+1/y$, and by \eqref{eq:HBS} the system $\NN$ adds
  $\le\int_{\Xm}^y\eps_1(v)\,dv/v+\eps_1(y)\le 2\eta_1\le 0.002$.
\item $C_{\mathrm{BS}}\le 38.3$ and $C_S\le 77$ (Theorem~\ref{thm:bs}, Corollary~\ref{cor:saias});
  $C_N:=C_S(2.04+4.71/\log X_0)\le 158$ (Corollary~\ref{cor:kr-weights}).
\end{itemize}

\paragraph{Regime and $o(1)$.} $X_0\ge\max(e^{10^8},X_{\mathrm{Hux}})$ is fixed; it is never sent to infinity. In this
section, $o(1)$ denotes a term bounded, at every $\Lambda\ge10^8$ and uniformly in $x$, by
$\exp(-10^{-3}X_0^{0.005}/\Lambda)+10^5\Lambda^2/X_0^2\le 10^{-3}$. These are the terms of
Proposition~\ref{prop:g1}(B), Steps~1--4: errors of the local prime number theorem at $\log x\ge X_0^{1/100}$, and
$\kappa_0$-terms with $\log x_k\ge X_0^2$. So $0.97-o(1)\ge0.969$ and $0.96-o(1)\ge0.959$.

\subsection{Prime inputs}\label{sec:noise-primes}
% v3: §3.1

\paragraph{The hypothesis \textup{(H$_{\mathrm{BS}}$)}.} $N(v)=\lfloor v\rfloor$ for $1\le v<\Xm$, and
$|N(v)-\lfloor v\rfloor|\le\eps_1(v)\,v$ for $v\ge\Xm$. Put
\begin{equation}\label{eq:HBS}
  \eta_1:=\sup_x\sum_{k\ge0}\sup\{\eps_1(v):v\in[x2^{-k-1},x2^{-k}]\cap[\Xm,\infty)\},\qquad
  \eta_2:=\sup_{v\ge\Xm}\eps_1(v)\log v,
\end{equation}
and require $\eta_1,\eta_2\le 10^{-3}$. In the construction, write
$N-\lfloor\cdot\rfloor=(N-T)+(T-\lfloor\cdot\rfloor)$.
\begin{itemize}
\item On $\Good(<z)$ (Definition~\ref{def:good}), $|N-T|\le B(u)+H(u)/2+2$ at scales $u<z$, by Theorem~\ref{thm:queue}
  and Corollary~\ref{cor:lag}.
\item $T-\lfloor\cdot\rfloor$ is the B-lin error plus $R$ plus the absorbed part. By \eqref{eq:T1} and
  Lemma~\ref{lem:lemmaP} it is bounded by $2\kappa u\cdot\min(1,C\rho(\log u/\log X_0))+O(u^{\theta}\log^*u)$.
\item On $[\Xm,X_0)$, $N-\lfloor\cdot\rfloor$ counts the jittered primes landing below $X_0$. There are at most
  $\#\{\text{slots in }[X_0,X_0+H]\}\le 2H(X_0)$ of them, so $\eps_1\le 3H(X_0)/X_0$ there.
\item Hence $\eta_1,\eta_2=O(\kappa\log X_0)+O(X_0^{-1/2})\ll 10^{-3}$.
\end{itemize}
No sup-norm power saving is assumed; under B-lin it would fail near $u\approx X_0^2$.

\begin{lemma}[Chebyshev--Mertens]\label{lem:cm}
% v3: Lemma 3.1
Assume \textup{(H$_{\mathrm{BS}}$)} on $[1,z]$, and let $\psi(x):=\sum_{p^k\le x}\log p$, summed over the realised
primes. Then for all $1\le x\le z$:
\begin{enumerate}[label=\textup{(\roman*)}]
\item $\psi(x)-\psi(x/2)\le 0.71x$, and $\psi(x)\le 1.42x$;
\item $S(x):=\sum_{p^k\le x}\log p/p^k$ satisfies $-1.01\le S(x)-\log x\le 0.43$; we put $K_M:=1.01$;
\item $S_p(x):=\sum_{p\le x}\log p/p$ satisfies $-1.77\le S_p(x)-\log x\le 0.43$;
\item for $2\le y\le x\le z$,
  \[
    \Big|\sum_{y<p\le x}\frac1p-\log\frac{\log x}{\log y}\Big|\le\frac{2.20}{\log y}=:\frac{K'}{\log y}.
  \]
\end{enumerate}
The proof at $x$ uses $N$ only on $[1,x]$. In particular \textup{(i)--(iv)} hold surely on $\Good(<z)$.
\end{lemma}

\begin{proof}
Write $N=\lfloor\cdot\rfloor+E_1$, with $E_1$ supported on $[\Xm,\infty)$ and $|E_1(v)|\le\eps_1(v)v$. Since $\log$ is
additive on $\NN$, we have $\log\cdot\,dN=d\psi*dN$, that is, the Chebyshev identity
$\int_{[1,x]}\log v\,dN(v)=\int_{[1,x]}N(x/a)\,d\psi(a)$. Inserting $N=\lfloor\cdot\rfloor+E_1$ on both sides and writing
$\lfloor x/a\rfloor=x/a-\{x/a\}$ gives
\begin{equation}\label{eq:cm-identity}
  x\,S(x)=\sum_{n\le x}\log n+\int\{x/a\}\,d\psi(a)+\int_{[1,x]}\log v\,dE_1(v)-\int E_1(x/a)\,d\psi(a).
\end{equation}

\emph{Two error integrals.} By partial summation,
$\int_{[1,x]}\log v\,dE_1(v)=E_1(x)\log x-\int_{\Xm}^xE_1(v)\,dv/v$, so its absolute value is at most
$x(\eta_2+\sup\eps_1)$. Next, group $a$ into dyadic blocks $[A,2A)$ with $x/A\ge\Xm$. On such a block,
$|E_1(x/a)|\le\eps_1(x/a)\,x/a$, and $\int_{[A,2A)}(x/a)\,d\psi(a)\le(x/A)(\psi(2A)-\psi(A))\le 1.42x$ by (i) at the
smaller scale $2A\le 2x/\Xm$. Summing over the blocks,
\begin{equation}\label{eq:cm-psi}
  \int|E_1(x/a)|\,d\psi(a)\le 1.42\,\eta_1\,x .
\end{equation}
The proof of (i) is by induction on $x$, so this use of (i) at scales $\le 2x/\Xm$ is legitimate.

\emph{(i).} Put $\ell_N(x):=\int_{[1,x]}\log v\,dN$. The Chebyshev identity gives
$\ell_N(x)-2\ell_N(x/2)=\int[N(x/a)-2N(x/2a)]\,d\psi(a)$. In the bracket, $\lfloor v\rfloor-2\lfloor v/2\rfloor\ge
\mathbf 1_{[1,2)}(v)$, so the $\lfloor\cdot\rfloor$ part contributes at least $\psi(x)-\psi(x/2)$. Hence
\[
  \psi(x)-\psi(x/2)\le\big[S_0(x)-2S_0(x/2)\big]+\Big|\int_{[1,x]}\!\log v\,dE_1\Big|+2\Big|\int_{[1,x/2]}\!\log v\,dE_1\Big|
  +\int\big(|E_1(x/a)|+2|E_1(x/2a)|\big)\,d\psi(a),
\]
with $S_0(x):=\sum_{n\le x}\log n$. By Stirling, $S_0(x)-2S_0(x/2)\le x\log2+3\log x+3$. The two log-integrals are at
most $4x(\eta_2+\sup\eps_1)\le 0.008x$, and by \eqref{eq:cm-psi} the last integral is at most
$2.84\,\eta_1x\le 0.0043x$. So
$\psi(x)-\psi(x/2)\le(\log2+0.008+0.0043)x+O(\log x)=0.7054x+O(\log x)\le0.71x$ for $x\ge\Xm$. Below $\Xm$ the system
is $\N$, where the maximum of $(\psi(x)-\psi(x/2))/x$ is $0.680$, attained at $x=5$. Summing over dyadic blocks gives
$\psi(x)\le 0.71x\sum_{k\ge0}2^{-k}=1.42x$.

\emph{(ii).} In \eqref{eq:cm-identity}, $0\le\int\{x/a\}\,d\psi\le\psi(x)\le1.42x$, and
$|S_0(x)-(x\log x-x)|\le\log(2\pi x)$. The two $E_1$-terms are at most $(\eta_2+\sup\eps_1+1.42\eta_1)x\le 0.004x$. So
$S(x)-\log x\in[-1.004,\,0.424]$ for $x\ge\Xm$. Below $\Xm$ it equals the value for $\N$, which lies in $[-0.753,0]$:
this is an exact computation for $x\le10^7$. For $10^7\le x<\Xm$, Rosser and Schoenfeld
\cite[(2.8) and Thm.~6, (3.21)--(3.22)]{RosserSchoenfeld1962} give
$\sum_{p\le x}(\log p)/p=\log x+E\pm1/(2\log x)$, where the lower bound holds for $x>1$ and the upper for $x\ge319$. By
their (2.8), adding the prime powers gives $S(x)-\log x=-\gamma\pm0.04$ there.

\emph{(iii).} $S-S_p=\sum_{p^k\le x,\,k\ge2}\log p/p^k\in[0,0.7554]$: the rational primes contribute at most
$\sum_p\log p/(p(p-1))=0.75537$, and primes $\ge\Xm$ add at most $3/X_0$. Combine with (ii).

\emph{(iv).} Put $E_p:=S_p-\log\in[-1.77,0.43]$. Partial summation on $(y,x]$ gives
\[
  \sum_{y<p\le x}\frac1p-\log\frac{\log x}{\log y}=\frac{E_p(x)}{\log x}-\frac{E_p(y)}{\log y}
  +\int_y^x\frac{E_p(v)\,dv}{v\log^2v}.
\]
If $m\le E_p\le M'$, the first and third terms together lie in $[m/\log y,\,M'/\log y]$, and the second in
$[-M'/\log y,\,-m/\log y]$. So the left side has absolute value at most $(M'-m)/\log y=(0.43+1.77)/\log y$.
\end{proof}

\begin{remark}\label{rem:cm-chain}
Part (iv) is the form used in the pivot chains of \S\ref{sec:queue} and Appendix~\ref{app:rankfull}. It applies once
per chain, so nothing accumulates. A single-function version also exists: the function
$E_z(\sigma):=\sum_{p\le e^\sigma}1/p-\log\sigma-M_z$, with $M_z$ defined from the primes $\le 2z$, satisfies
$|E_z(\sigma)|\le 2.20/\sigma$, but it is not adapted.
\end{remark}

\paragraph{Facts on Dickman's function.} Let $\rho$ be Dickman's function, with $\rho(v):=1$ for $v\le1$, and put
$r(v):=-\rho'(v)/\rho(v)=\rho(v-1)/(v\rho(v))$ for $v>1$.
\begin{enumerate}[label=(R\arabic*)]
\item\label{R1} Exact identities. $v\rho(v)=\int_{v-1}^v\rho$ for $v\ge1$, and $\int_a^\infty\rho(w-1)\,dw/w=\rho(a)$
  for $a\ge1$ (both sides vanish at $\infty$ and have derivative $-\rho(a-1)/a$). For $u\ge2$,
  \[
    \int_{u-1}^u\rho(w-1)(u-w)\frac{dw}{w}=\rho(u-1)-u\rho(u)=\rho(u-1)\big(1-1/r(u)\big),
  \]
  which follows from the first two identities: the left side is $u(\rho(u-1)-\rho(u))-(u-1)\rho(u-1)$.
\item\label{R2} $r$ is increasing on $(1,\infty)$, so $\rho$ is log-concave on $[1,\infty)$. On $[1,2]$,
  $r=1/(v(1-\log v))$, which is increasing. For $v>2$, the first identity of \ref{R1} gives $r(v)=1/I(v)$ with
  $I(v):=\int_0^1\exp(-\int_{v-1}^{v-1+s}r)\,ds$, and
  $I'(v)=-\int_0^1e^{(\cdots)}\big(r(v-1+s)-r(v-1)\big)\,ds$. Let $v^*:=\inf\{v>2:r'(v)\le0\}$. Then $r$ is strictly
  increasing on $[v^*-1,v^*]$, so $I'(v^*)<0$ and $r'(v^*)>0$, a contradiction.
\item\label{R3} $r(2)=1/(2-2\log2)=1.6294$. Log-concavity gives $I(v)\le(1-e^{-r(v-1)})/r(v-1)$, hence
  $r(v)\ge r(v-1)/(1-e^{-r(v-1)})$. Iterating from $r(2)$ gives a computable lower bound $r_{\mathrm{low}}$. Also
  $\rho(w-1)/\rho(u-1)\le e^{-r(u-1)(w-u)}$ for $w\ge u$.
\item\label{R4} For $\delta\in(0,1)$, \ref{R1} and monotonicity give $v\rho(v)\ge\delta\rho(v-1+\delta)$, so
  $\exp(\int_{v-1+\delta}^vr)=\rho(v-1+\delta)/\rho(v)\le v/\delta$. Since $r$ is increasing,
  $\int_{v-1+\delta}^vr\ge(1-\delta)r(v-1+\delta)$, and with $x:=v-1+\delta$,
  \[
    r(x)\le\frac{\log((x+1)/\delta)}{1-\delta}\qquad(0<\delta<1).
  \]
  The minimum over $\delta$ is denoted $r_{\mathrm{up}}(x)$; for example $r_{\mathrm{up}}(100)=7.66$. Taking
  $\delta=1/\log x$ gives $r(x)\le\log x+\log\log x+2$ for $x\ge30$.
\end{enumerate}

\begin{lemma}[smooth elements, with Dickman quality]\label{lem:smooth-v}
% v3: Lemma 3.2
Assume \textup{(H$_{\mathrm{BS}}$)}. Let $z>1$, $L_z:=\log z\ge100$, and $K:=3.2\ge K_M$. For $\ell:=\log Y=uL_z$ with
$1\le u\le L_z$,
\[
  v(\ell,z):=\sum_{n\in\NN\cap(Y,2Y],\ \Pplus(n)\le z}\frac1n\le C_0\,\rho(u)\,e^{a_V},\qquad
  a_V:=\frac{4Ku^2\log(u+2)}{\ell},\qquad C_0:=0.70.
\]
\end{lemma}

\begin{proof}
\emph{Hildebrand's identity.} Since $\log n=\sum_{p^k\mid n}\log p$ in the free semigroup $\NN$, and writing
$n=p^km$,
\[
  \sum_{n\in(Y,2Y],\ \Pplus(n)\le z}\frac{\log n}{n}=\sum_{p\le z,\ k\ge1}\frac{\log p}{p^k}\,v(\ell-k\log p,z).
\]
The dyadic window shifts exactly, so this is an identity. Its left side is $\ge\ell\,v(\ell,z)$.

\emph{Small heights.} For $0\le\ell'\le L_z$, $v(\ell',z)\le\sum_{n\in(e^{\ell'},2e^{\ell'}]}1/n$. For
$e^{\ell'}\ge72$ this is at most $\log2+1/(4e^{\ell'})+3(\eta_2+\sup\eps_1)\le C_0$; on $\N$ the exact value is at most
$0.6966$ there. For $1\le e^{\ell'}<72$ it is the value for $\N$, which is at most $0.8334$ (attained for
$e^{\ell'}\in[1.5,2)$). For $\ell'<0$, $v(\ell',z)\le1$.
% Logged correction 10 (CORRECTIONS-vs-v3.md): v3's base case v <= C_0 for all l' <= L_z fails for 1.5 <= e^{l'} < 72;
% the excess enters the induction step as X(u) and is certified (aux_cert_arb.py, third line).

\emph{Induction.} We show $v(\ell',z)\le G_V(\ell'):=C_0\rho(\ell'/L_z)e^{a_V(\ell')}$ for $\log72\le\ell'\le L_z^2$, by
induction on $\ell'$ at fixed $z$; for $\ell'\le L_z$ this is the bound just proved. As a function of $u'=\ell'/L_z$,
$a_V=4Ku'\log(u'+2)/L_z$ is increasing. Let $\ell\in(L_z,L_z^2]$, so $u\in(1,L_z]$, and put $F(\sigma):=G_V(\ell-\sigma)$ for
$0\le\sigma\le L_z$. Since $\rho$ is non-increasing and $a_V$ is increasing,
\begin{equation}\label{eq:lemV-maxF}
  \max_{0\le\sigma\le L_z}F(\sigma)\le C_0\,\rho(u-1)\,e^{a_V(u)} .
\end{equation}
($F$ is increasing in $\sigma$ when $u\ge2$, but for $1<u<2$ it is only unimodal; \eqref{eq:lemV-maxF} is what is used.)
% Logged correction 10: v3 asserts F is increasing for all u; for 1<u<2 it is unimodal, and max F <= C_0 rho(u-1) e^{a_V(u)}
% gives the same term 2K_M r(u).
\begin{itemize}
\item \emph{Main terms, $p^k\le z$.} Partial summation against $S=\log+E$, where $E(0)=0$ and $|E|\le K_M$
  (Lemma~\ref{lem:cm}(ii)), gives $\int F\,dE=F(L_z)E(L_z)-\int E\,dF\le K_MF(L_z)+K_M\,\TV(F)\le2K_M\max F$, since $F$ is
  unimodal. Hence
  \[
    \sum_{p^k\le z}\frac{\log p}{p^k}F(k\log p)\le\int_0^{L_z}F(\sigma)\,d\sigma+2K_M\,C_0\rho(u-1)e^{a_V(u)} .
  \]
\item \emph{Small heights.} Here $\ell-k\log p\ge\ell-L_z=(u-1)L_z$. If $e^{(u-1)L_z}<72$, some of these terms have
  $e^{\ell-k\log p}<72$, where $v$ may exceed $G_V$ by up to $0.8334-C_0\le0.1334$. Their total excess is at most
  \[
    0.1334\sum_{p^k\in(Y/72,z]}\frac{\log p}{p^k}\le0.1334\,(\log72+1.44)\le0.7627,
  \]
  by Lemma~\ref{lem:cm}(ii). This can happen only for $1<u<1+\log72/L_z\le1.043$.
\item \emph{Prime powers $p^k>z$ with $p\le z$.} These contribute at most $22e^{-0.4L_z}$ relative to the normaliser below.
  This uses $\rho(u-t)/\rho(u)\le e^{t\,r(u)}$, $v\le G_V/C_0$ for $\ell'<\log72$, and
  $1-r(u)/L_z\ge1-r_{\mathrm{up}}(100)/100\ge0.92$.
\end{itemize}

\emph{The one-step inequality.} Substitute $s=u-\sigma/L_z$. Then $\int_0^{L_z}F=L_zC_0\int_{u-1}^u\rho(s)e^{a_V(s)}ds$,
where $a_V(s)$ is evaluated at $\ell'=sL_z$, and $\ell G_V(\ell)=uL_zC_0\rho(u)e^{a_V(u)}$. By the first identity of \ref{R1},
$u\rho(u)=\int_{u-1}^u\rho(s)\,ds$. Let $\alpha:=4K\big(\log(u+1)+(u-1)/(u+1)\big)/L_z$; then $a_V(u)-a_V(s)\ge\alpha(u-s)$ on
$[u-1,u]$, and $e^{-\alpha t}\le1-(1-e^{-\alpha})t$ for $t\in[0,1]$. Let $E_*$ be a lower bound for the mean of $u-s$ under the
density proportional to $\rho(s)$ on $[u-1,u]$. With $\rho(u-1)=u\,r(u)\,\rho(u)$, this gives
\[
  \ell G_V(\ell)-\int_0^{L_z}F-2K_MC_0\rho(u-1)e^{a_V(u)}\ \ge\ \mathcal N(u)\cdot
  \Big[L_z(1-e^{-\alpha})E_*-2K_M\,r(u)\Big],\qquad\mathcal N(u):=C_0e^{a_V(u)}\rho(u)\,u .
\]
Both sides are in the absolute units of the inequality $\ell v(\ell)\le\ell G_V(\ell)$. The bracket is of order $1$, not of
order $L_z$: $L_z(1-e^{-\alpha})\approx4K\log(u+1)$. Two choices of $E_*$ are available:
\begin{itemize}
\item $E_*=\frac12$ always, since that density is decreasing;
\item $E_*=1-1/r(u-1)$ for $u\ge2$: by \ref{R2} the density is dominated in likelihood ratio by the truncated exponential
  with rate $r(u-1)$, whose mean is at least $1-1/r(u-1)$.
\end{itemize}
Dividing the excesses by the normaliser $\mathcal N(u)$, the induction step closes if
\begin{equation}\label{eq:lemV-onestep}
  L_z(1-e^{-\alpha})\max\big(\tfrac12,\,1-1/r_{\mathrm{low}}(u-1)\big)-2K_M\,r_{\mathrm{up}}(u)-22e^{-0.4L_z}-X(u)>0,
\end{equation}
where $X(u):=0.7627/(C_0(1-\log u)u)$ for $u<1+\log72/L_z$ (using $\mathcal N(u)\ge C_0\rho(u)u$ and $\rho(u)=1-\log u$ on
$[1,2]$), and $X(u):=0$ otherwise. The left side of \eqref{eq:lemV-onestep} is increasing in $L_z$, and the range of $X$
shrinks as $L_z$ grows, so $L_z=\max(100,u)$ is the worst case.
\begin{itemize}
\item \emph{Computer-assisted check} (see Appendix~\ref{app:cert}). With $r$ enclosed in $[r_{\mathrm{low}},r_{\mathrm{up}}]$ (and
  exact on $[1,2]$), the left side of \eqref{eq:lemV-onestep} without $X$ (and with $15e^{-0.4L_z}$) is positive on
  $u\in[1,10^6]$: $13{,}823$ cells in ball arithmetic, no failures, minimum certified value $1.724$ (the difference
  $7e^{-0.4L_z}\le10^{-16}$ between the two prime-power allowances is immaterial). On $u\in[1,2]$, the left
  side \emph{with} $X$ and $22e^{-0.4L_z}$, and even with an additional allowance $2K_Mr_{\mathrm{up}}(u)(e^{0.141}-1)$, is
  positive: $1{,}387$ cells, no failures, minimum certified value $0.829$.
% Logged correction 10: the small-height excess X(u) and the [1,2] check (aux_cert_arb.py, third line) are new.
\item For $u>10^6$, use $r_{\mathrm{low}}\ge16.5$ and $r_{\mathrm{up}}(u)\le\log u+\log\log u+2$. The left side is then at least
  $2K_M\big[3\log(u+1)\cdot0.94-\log u-\log\log u-2\big]>0$, since $K\ge3K_M$.
\end{itemize}
This completes the induction. The range $\ell'\le L_z^2$ covers $u\le L_z$.
\end{proof}

% 04b-identity.tex — v3 §3.2.
\subsection{The margin identity and the push-through form}\label{sec:noise-identity}
% v3: §3.2

\paragraph{Bookkeeping} (Lemmas~\ref{lem:count}, \ref{lem:excess}, Corollary~\ref{cor:lag}). Rigid integers and
elements of $\Abs$ sit on their own unit jumps of $T$ (jump slots). Every level crossing of $T$ is either a jump slot or a
continuous slot, so on every window $W$
\[
  \#\{\text{continuous slots in }W\}=\int_WT'_{\mathrm{cont}}\pm2 .
\]
% Logged correction 95 (CORRECTIONS-vs-v3.md): v3 states the bookkeeping on slot windows with a +-2 for "the two slot
% boundaries", but the margin identity applies it to the physical window W_y, whose end cells are cut, and its nu has no edge
% term (Lemma noise adds one). On slot windows the identity is exact; for any interval the residual E_edge of the two
% partial end cells is defined and bounded below, and it now enters nu.
By Lemma~\ref{lem:count}, on a slot window $W=(k,j]$ (the slots $k+1,\dots,j$, with physical interval $[s_{k+1},s_{j+1})$) the
number of primes decided equals the number of continuous slots, minus the non-absorbed jittered composites landing in $W$,
plus the change in the queue:
\begin{equation}\label{eq:bookkeeping}
  \sigma(W)=\#\{\text{continuous slots in }W\}-C_{\mathrm{jit}}(W)+\Delta Q(W),\qquad\Delta Q(W):=Q_j-Q_k .
\end{equation}
Indeed $\sum_{i=k+1}^j(P_i+A_i-1)=Q_j-Q_k$ by the greedy rule, and the composites with value in $[s_{k+1},s_{j+1})$, counted by
$\sum_{i=k+1}^jA_i$, are the elements of the jump slots of $W$, one each, and the $C_{\mathrm{jit}}(W)$ others. Here $C_{\mathrm{jit}}$ counts the composites $n\notin\Abs$ whose
latest prime $q$ (in decision order) is jittered; the cofactor may be any element of $\NN$. Write
$C_{\mathrm{jit}}=C^{(1)}+C^{(\ge2)}$, according to whether $q\,\|\,n$ or $q^2\mid n$.

\emph{Physical windows.} For an arbitrary interval $W$, let $k$ be the last slot that starts before $W$ ($k:=j_0-1$ if there
is none), and let $k+1,\dots,j$ ($j\ge k$) be the slots that start in $W$. Then $\sigma(W)$ and the continuous slots in $W$
refer to these slots (\S\ref{sec:greedy}), while $C_{\mathrm{jit}}(W)$ counts by value. Put $\Delta Q(W):=Q_j-Q_k$, and let
$W^\circ:=[s_{k+1},s_{j+1})$ be the physical interval of the slot window $(k,j]$. The \emph{edge residual} of $W$ is
\begin{equation}\label{eq:edge}
\begin{aligned}
  E_{\mathrm{edge}}(W)&:=\sigma(W)-\#\{\text{continuous slots in }W\}+C_{\mathrm{jit}}(W)-\Delta Q(W)\\
  &=C_{\mathrm{jit}}(W\setminus W^\circ)-C_{\mathrm{jit}}(W^\circ\setminus W),
\end{aligned}
\end{equation}
where the second form is \eqref{eq:bookkeeping} for $(k,j]$. Here $W\setminus W^\circ$ lies in the cell $J_k$ (for $k=j_0-1$,
below $s_{j_0}$, where there are no jittered composites, Lemma~\ref{lem:count}), and $W^\circ\setminus W$ (empty if $j=k$) lies
in the cell $J_j$, which ends after $W$. These are the two partial end cells. Since $\sigma(W)$, the continuous-slot count and
$\Delta Q(W)$ refer to the same slots, the slot starts cancel, and only the composites in the cut parts of the two end cells
remain. As $A_i\le Q_i+1$ (by $Q_i=Q_{i-1}+A_i+P_i-1$ with $Q_{i-1},P_i\ge0$),
\[
  -A_j\le E_{\mathrm{edge}}(W)\le A_k,\qquad|E_{\mathrm{edge}}(W)|\le1+\max(Q_k,Q_j),
\]
with $A_{j_0-1}:=0$. If $W$ is the physical interval of a slot window, then $W^\circ=W$ and $E_{\mathrm{edge}}(W)=0$. (Its $k,j$ may differ from those of the slot window by an empty slot at an end; $W^\circ$ is the same.)
\begin{itemize}
\item $\Comp$ is the predictable compensator of $C^{(1)}$. Given $\mathcal F_{q-}$, the element $q\cdot w$ ($q\nmid w$)
  lands in $W$ with conditional probability $|J_q\cap W/w|/H(s_q)$, and $q\cdot w$ is formed only if $q$ is decided
  after every prime of $w$, i.e.\ $s_q>\tau_w$. Hence
  \[
    \Comp(W)=\sum_{w\ge2}\int_{W/w}J(\Pi_{\tau_w}\sigma)\cdot\mathbf 1[qw\notin\Abs],
  \]
  where $\Pi_\tau\sigma$ keeps the slots $s_q>\tau$.
% Logged correction 83 (CORRECTIONS-vs-v3.md): the sum defining Comp^(2) is over the jittered q (as is C^(>=2)), the m decided
% before q, and the elements outside Abs; v3 leaves the range of q, the latest-prime filter and the absorption indicator implicit.
\item $\Comp^{(2)}$ is the predictable compensator of $C^{(\ge2)}$:
  $\Comp^{(2)}(W):=\sum_q\sum_{a\ge2}\sum_m\Prob(q^am\in W,\ q^am\notin\Abs\mid\mathcal F_{q-})$, over the jittered primes $q$
  and the $m\in\NN$ with $q\nmid m$ and $\tau_m<s_q$ (including $m=1$).
\item $\Mart':=(C^{(1)}-\Comp)+(C^{(\ge2)}-\Comp^{(2)})=:\Mart+\Mart^{(2)}$. Both parts are martingales in decision
  order. They are controlled by separate Freedman bounds (\S\ref{sec:noise-perscale}), because the increment of $q$ enters
  $\Mart$ through $W/q$ and $\Mart^{(2)}$ through $W/q^a$.
\end{itemize}

Since $\int_{W_x/w}g=\frac1w\int_{W_x}g(y/w)\,dy$,
\[
  \Comp(W_x)/H(x)=(B\Kpt\rhot)(x)-(BK_{\Abs}\rhot)(x)+\NT(x).
\]
\begin{itemize}
\item $K_{\Abs}$ is the part of $\Kpt$ whose landing is absorbed ($y/w<y(y)$, $\log y\ge L_*$).
\item The cutoff $\mathbf 1[y/w\ge\Xm]$ in $\Kpt$ costs nothing, since $\rhot$ vanishes below $\Xm$.
\item $\NT$ replaces the slot filter $s_q>\tau_w$ by the landing filter $y/w>\tau_w$ on the non-absorbed part. It is
  supported on \emph{$\tau$-near-threshold} cofactors, i.e.\ those with $|\tau_w-y/w|\le2H(y/w)$, and
  \begin{equation}\label{eq:NT}
    |\NT(y)|\le2\eta(y)\cdot\sup\rhot,\qquad\eta(y):=\sum_{w\ \tau\text{-near at }y}\frac1w .
  \end{equation}
\item $\NT_{\mathrm{pt}}$ denotes $\NT$ before the $B$-average, and $m:=T'-\Kpt\rhot-\NT_{\mathrm{pt}}$ is the
  \emph{pointwise margin density}.
\end{itemize}

\begin{proposition}[the margin identity]\label{prop:margin-identity}
% v3: Identity (3.1)
% Logged correction 95: nu contains the edge residual of W_y (eq:edge) and the box term E_box, so the identity is exact.
For every $y\ge X_0$,
\begin{equation}\label{eq:margin-identity}
\begin{gathered}
  \rhot+B\Kpt\rhot=F,\qquad F:=\tilde T'-\nu-\NT,\\
  \nu(y):=\frac{\Mart'(W_y)-\Delta Q(W_y)-E_{\mathrm{edge}}(W_y)+\Comp^{(2)}(W_y)+E_{\mathrm{box}}(y)}{H(y)},
\end{gathered}
\end{equation}
where $\Delta Q(W_y)$ and $E_{\mathrm{edge}}(W_y)$ are those of the physical window $W_y$ \eqref{eq:edge}, and
$E_{\mathrm{box}}(y):=\sigma(W_y)-H(y)\rhot(y)$. Moreover $|E_{\mathrm{box}}|\le2.01$.
\end{proposition}

\begin{proof}
% Logged correction 95: v3 applies the slot-window bookkeeping to W_y and puts the box term into theta, although tilde T' is
% defined without it; both residuals are now explicit terms of nu.
By \eqref{eq:edge} for the physical window $W_y$, exactly,
\[
  \sigma(W_y)=\#\{\text{continuous slots in }W_y\}-C_{\mathrm{jit}}(W_y)+\Delta Q(W_y)+E_{\mathrm{edge}}(W_y).
\]
Divide by $H(y)$. On the left,
$\sigma(W_y)/H(y)=\rhot(y)+E_{\mathrm{box}}(y)/H(y)$. On the right, write $C_{\mathrm{jit}}=\Comp+\Comp^{(2)}+\Mart'$, use the display
for $\Comp(W_y)/H(y)$ above, and move $B\Kpt\rhot$ to the left. The terms $BK_{\Abs}\rhot$ and
$\#\{\text{continuous slots in }W_y\}/H$ make up $\tilde T'$, the term $-\NT$ remains, and the rest is $-\nu$; there is no error term.

The box term: $H(y)\rhot(y)=\sum_qH(y)H(s_q)^{-1}\mathbf 1[y-H(s_q)/2<s_q\le y+H(s_q)/2]$, while $\sigma(W_y)$ counts the $q$ with
$s_q\in W_y$. Every box is centred, with width $H(s_q)=H(y)(1+O(\eps H(y)/y))$ for these $q$ by \eqref{eq:RH}, and distinct
decided primes have slot starts at least $1-10^{-8}$ apart (they are continuous-slot starts, \S\ref{sec:greedy}). So at most one prime near each end of
$W_y$ is counted differently, and the weights $H(y)/H(s_q)$ change the rest by $O(\eps H(y)^2/y)\le0.01$. Hence
$|E_{\mathrm{box}}|\le2.01$.
\end{proof}

The compensator margin density at $y$ is
\begin{equation}\label{eq:M}
  M(y):=\tilde T'(y)-(B\Kpt\rhot)(y)=\rhot(y)+\nu(y)+\NT(y).
\end{equation}
The noise $\nu$ enters only at the current scale $y$: it depends on the data on $W_y$ and on the two slot cells at its ends (Lemma~\ref{lem:noise}). Every other
quantity in \eqref{eq:margin-identity} at $y$ involves $\sigma$ only at scales $\le y/2+H$.

\paragraph{The push-through form.} $\Kpt$ maps functions supported on scales $\ge v$ to functions supported on scales
$\ge2v$. So $1+B\Kpt$ and $1+\Kpt B$ are invertible, and their Neumann series terminate on bounded ranges. For any such
operators $B,K$,
\[
  (1+BK)^{-1}=(1-B)+B(1+KB)^{-1}\big(1-K(1-B)\big),
\]
since $(1+BK)(1-B)+B(1-K(1-B))=1-B+BK-BKB+B-BK+BKB=1$. With $K=\Kpt$ and $\Rres:=(1+\Kpt B)^{-1}$, \eqref{eq:margin-identity}
gives
\begin{equation}\label{eq:pushthrough}
  \rhot=(1-B)F+B\Rres F-B\Rres\Kpt(1-B)F .
\end{equation}
\begin{itemize}
\item $\Rres$ is the \emph{long} operator: at each step the sharp landing filter, followed by an $H(u)$-average at the
  landing scale $u$. Its kernels $C_n$ are studied in \S\ref{sec:kernels}.
\item $(1-B)$ acts only on the source. The response $\rhot$ is never differentiated, so no regularity lemma for it is needed.
\item Applying the pointwise inverse $(1+\Kpt)^{-1}$ to \eqref{eq:margin-identity} instead would leave the response term
  $(1-B)\Kpt\rhot$, which is of size $\asymp1/\log y$ in sup norm. This is why the push-through form is used.
\end{itemize}
Note also the push-through identity $(1+B\Kpt)^{-1}B=B\Rres$ and its consequence $\Kpt B\Rres=1-\Rres$.

\paragraph{The jump split.} Split $T'=\Tc+J$ with $J=\jz\mathbf 1_{[X_0,\infty)}$ (\S\ref{sec:noise-notation}). The
$J$-part of $F$ is $BJ$. By the push-through identity its contribution to $\rhot$, and hence to the margin
$M=\rhot+\nu+\NT$, is exactly $B\Rres J$. For the rest put
\begin{equation}\label{eq:decomp}
  F_c:=F-BJ=\Tc+F_r,\qquad F_r:=(B-1)\Tc+\vartheta/H-\nu-\NT,\qquad G:=\Kpt(1-B)F_c .
\end{equation}
By \eqref{eq:pushthrough} applied to $F_c$,
$M=B\Rres J+(1-B)F_c+B\Rres F_c-B\Rres G+\nu+\NT$.

\emph{Placement convention.} The absorbed source $\sigma_{\Abs}$ is part of $\Tc$, never of $F_r$. Only its
$H$-oscillation enters, through $(B-1)\Tc$ and $o_u$ (see \eqref{eq:osc}). The jump enters only through $B\Rres J$, that is,
through Lemma~\ref{lem:xb}.

% 04c-saias.tex — v3 §3.3: Lemma M, Theorem 3.4 (Beurling–Saias) with full proof, Lemma 3.4A, Corollary 3.5.
% Author: gap-localpnt.
\subsection{Sharp coefficients: Lemma M and the Beurling--Saias bound}\label{sec:noise-saias}
% v3: §3.3

For a start $y$, the sharp resolvent is $((1+\Kpt)^{-1}g)(y)=\sum_{n\in\NN}c_n(y)\,g(y/n)/n$, where $c_1=1$ and
\[
  c_n=-\sum_{d\mid n,\ d<n}\mathbf 1[\tau_{n/d}<y/n]\cdot\mathbf 1[y/n\ge\Xm]\cdot c_d .
\]
Expanding the recursion, $c_n(y)$ is the signed count of ordered factorisations $n=w_1\cdots w_k$ into
$w_i\in\NN\setminus\{1\}$ whose cumulative filters $\tau_{w_i}\cdot w_1\cdots w_i<y$ and cutoffs
$y/(w_1\cdots w_i)\ge\Xm$ all pass, with sign $(-1)^k$. Since the landing points $y/(w_1\cdots w_i)$ decrease along a
chain, the cutoffs reduce to the single condition $y/n\ge\Xm$.

\begin{lemma}[Lemma M]\label{lem:lemmaM}
% v3: Lemma 3.3
Let $t:\PP\to(0,\infty)$ be any weight, with ties broken by a fixed order, and $\tau_w:=\max_{p\mid w}t(p)$. Then
$c_n(y)=\mu(n)\,\chi^t_y(n)\,\mathbf 1[y/n\ge\Xm]$, where $\chi^t_y(n)=1$ if and only if $n$ is squarefree and
\[
  t(p)\cdot\prod_{q\mid n,\ t(q)\ge t(p)}q<y\qquad\text{for every prime }p\mid n .
\]
In words: strip the primes of $n$ in decreasing $t$-order; after removing $p$, the current point must still exceed $t(p)$.
\end{lemma}

\begin{proof}
We may assume $y/n\ge\Xm$. Write $q\succ q'$ for the $t$-order.

\emph{$n$ squarefree.} Chains are ordered set partitions $(B_1,\dots,B_k)$ of the set of primes of $n$, with
$w_i=\prod B_i$. Let $\Pi_i:=w_1\cdots w_i$ and $F(B):=\max_{p\in B}t(p)$. A chain is valid if and only if
$F(B_i)\,\Pi_i<y$ for all $i$.
\begin{itemize}
\item Let $q_1\succ\cdots\succ q_r$ be the primes of $n$. The canonical chain $(\{q_1\},\dots,\{q_r\})$ is valid if and only
  if $\chi^t_y(n)=1$, since the primes $q$ with $t(q)\ge t(q_i)$ are exactly $q_1,\dots,q_i$. Its sign is
  $(-1)^r=\mu(n)$.
\item For a non-canonical valid chain, let $i$ be the first index with $B_i\ne\{q_i\}$, put $q:=q_i$, and let $j\ge i$
  be the index with $q\in B_j$. Every block from position $i$ on consists of primes $\preceq q$, so $F(B_j)=t(q)$.
  \begin{itemize}
  \item[(S)] If $|B_j|\ge2$, split $B_j$ into the two consecutive blocks $B_j\setminus\{q\},\{q\}$. The first new
    condition is implied by the old condition of $B_j$, since both its $F$ and its partial product are smaller. The second
    new condition is $t(q)\,\Pi_j<y$, which is the old condition of $B_j$.
  \item[(M)] If $B_j=\{q\}$ (so $j>i$), merge $B_{j-1}$ and $\{q\}$ into one block. Its $F$ is $t(q)$ and its partial
    product is $\Pi_j$, so its condition is the old condition of $\{q\}$; the other conditions are unchanged.
  \end{itemize}
  The maps (S) and (M) are mutually inverse, preserve validity and the first deviation $(i,q)$, and change $k$ by $1$.
  So the non-canonical valid chains cancel in pairs.
\end{itemize}

\emph{$n$ not squarefree} (induction on $\Omega(n)$). By induction, $c_d=\mu(d)\chi^t_y(d)$ for $d<n$, so
$c_n=-\sum_{d\in\mathcal F_n,\ d\text{ squarefree}}\mu(d)\chi^t_y(d)$ with
$\mathcal F_n:=\{d\mid n,\ d<n:\ \tau_{n/d}\cdot n<y\}$.
Let $s$ be the $t$-largest prime with $s^2\mid n$.
\begin{itemize}
\item $\mathcal F_n$ is stable under $d\mapsto d\triangle\{s\}$ (toggling $s$). Indeed $d\cdot s\mid n$ whenever $d$ is
  squarefree, $d\mid n$ and $s\nmid d$; and $n/d$, $n/(ds)$ have the same prime support, because both still contain $s$.
  So $\tau_{n/d}=\tau_{n/(ds)}$. Moreover, if $\mathcal F_n\ne\emptyset$ then $t(s)\,n\le\tau_{n/d}\,n<y$.
\item $\chi^t_y(d)=\chi^t_y(d\triangle\{s\})$. The conditions at primes $p\succ s$ do not involve $s$. The conditions at
  the other primes $p$ hold for both, since $t(p)\prod_{t(q)\ge t(p)}q\le t(s)\,\mathrm{rad}(n)\le t(s)\,n<y$.
\item $\mu$ changes sign under $d\mapsto d\triangle\{s\}$, so the terms cancel in pairs and $c_n=0$.\qedhere
\end{itemize}
\end{proof}

\paragraph{Consequences.}
\begin{itemize}
\item $|c_n|\le1$, and $\sum_n|c_n|/n\le\sum_{n\in\NN,\,n\le y/\Xm}1/n\le\log y+B_1$.
\item \emph{Density.} With the decision order $t$ ($|t(p)-p|\le H(p)$, and $t(p)=p$ for rational primes),
  $\chi^t_y(n)=1$ and $y/n\in[v,2v)$ force $n$ to be squarefree and \emph{$4v$-dense}: $p_{j+1}\le4v\,p_1\cdots p_j$ for the
  primes $p_1<\dots<p_r$ of $n$, sorted by value. To see this, fix $j$ and let $p^*$ be the prime of smallest $t$ among
  $p_{j+1},\dots,p_r$. All of $p_{j+1},\dots,p_r$ have $t\ge t(p^*)$, so the condition at $p^*$ gives
  $t(p^*)\,p_{j+1}\cdots p_r<y<2v\,n$, that is, $t(p^*)<2v\,p_1\cdots p_j$. Hence
  $p_{j+1}\le p^*\le t(p^*)+H(p^*)\le 4v\,p_1\cdots p_j$, since $H(p^*)\le p^*/10^{10}$.
\item Lemma~\ref{lem:lemmaM} concerns \emph{sharp scalar} filters. For the true operator $\Rres$, whose filter is
  averaged across the threshold, the pointwise analogue is false: for $n=4$ with void filter and constant $H$,
  $\TV(C_4)=\frac14$ while $\mu(4)=0$. The transfer to $\Rres$ is Theorem~\ref{thm:kr} and its random-width form Theorem~\ref{thm:kr-random}, which use only the masses $c_n$.
\end{itemize}

\begin{theorem}[Beurling--Saias, upper bound]\label{thm:bs}
% v3: Theorem 3.4
Assume \textup{(H$_{\mathrm{BS}}$)} and $\log X_0\ge10^8$. For $t\ge X_0$ put $\tau:=\log t$. Call a squarefree $n\in\NN$
\emph{$t$-dense} if $p_{j+1}\le t\,p_1\cdots p_j$ for its primes $p_1<\dots<p_r$. Then for every $Y\ge t$,
\[
  W_t(Y):=\sum_{n\in\NN\cap(Y,2Y],\ t\text{-dense}}\frac1n\le C_{\mathrm{BS}}\,\frac{\log t}{\log Y},\qquad
  C_{\mathrm{BS}}:=C_0\max\big(e^{\kappa_\Phi/\tau},\,30e^{\kappa_\Phi/(30\tau)}\big),\quad\kappa_\Phi:=4\tau+1500 .
\]
In particular $C_{\mathrm{BS}}\le0.70\,e^{4+1.5\cdot10^{-5}}\le38.3$. The same bound holds with $t$ replaced by $4t$, at
the cost of a factor $1+\log4/\log X_0$.
\end{theorem}

For $\NN=\N$ the order of magnitude is due to Saias \cite{Saias1997}. The point here is a proof that uses only
Lemmas~\ref{lem:cm} and~\ref{lem:smooth-v}, so that it applies to the realised system $\NN$. The exponent of $1/\log Y$ is
exactly $1$; this is what makes the noise sums of \S\ref{sec:noise-perscale} independent of $x$.

\begin{proof}[Proof of Theorem~\ref{thm:bs}]
\emph{Step 1: recursion.} Put $W_t(Y;z):=\sum_{n\in(Y,2Y],\ t\text{-dense},\ \Pplus(n)<z}1/n$. If $n$ is $t$-dense and
$p=\Pplus(n)$, then $n/p$ is $t$-dense with $\Pplus(n/p)<p$, and $p\le t\cdot n/p$, so $p^2\le tn\le2tY$. Summing over $p$
first (for each $p$ the cofactor range $(Y/p,2Y/p]$ is a single dyadic block, so the sum is exact),
\begin{equation}\label{eq:bs-rec}
  W_t(Y;z)\le\sum_{p<z,\ p\le\sqrt{2tY}}\frac1p\,W_t(Y/p;p).
\end{equation}

\emph{Step 2: logarithmic variables.} Put $\ell:=\log Y$, $u:=\ell/\log z$, and
$\Psi(\ell,u):=W_t(e^\ell;e^{\ell/u})\cdot\max(1,\ell/\tau)$. For a prime $p$ write
\[
  \sigma:=\log p,\qquad m:=\ell-\sigma,\qquad w:=m/\sigma,\qquad \nu:=\ell/m=(1+w)/w,\qquad
  w_1:=\frac{\ell-\tau-\log2}{\ell+\tau+\log2}.
\]
Then $p<z$ if and only if $w>u-1$, and $p\le\sqrt{2tY}$ if and only if $w\ge w_1$.
% Logged correction 6 (CORRECTIONS-vs-v3.md): w_1 includes log 2 (v3: (l-tau)/(l+tau)); Z1 bound 0.265 tau - 2714.
For $\ell\ge3\tau+\log2$ every term has $m\ge(\ell-\tau-\log2)/2\ge\tau$, so
$\max(1,\ell/\tau)/\max(1,m/\tau)=\nu$, and \eqref{eq:bs-rec} becomes
\begin{equation}\label{eq:bs-rec-log}
  \Psi(\ell,u)\le\sum_{p:\ w\ge w_{\mathrm{lo}}}\frac1p\,\nu\,\Psi(m,w),\qquad w_{\mathrm{lo}}:=\max(u-1,w_1).
\end{equation}
In the continuum the kernel is $dw/w$. Its exact fixed point is $g(u)=\rho(u-1)$: by \ref{R1},
$\int_{u-1}^\infty g(w)\,dw/w=g(u)$ for $u\ge2$, while for $u<2$
$\int_{w_1}^\infty g(w)\,dw/w=1+\log(1/w_1)$. So the kernel is critical, and every error must be summable along chains.
This is why the exponent of $1/\log Y$ is exactly~$1$.

\emph{Step 3: the profile.} Fix the constants
\[
  B:=5K'=11,\qquad\kappa_\Phi:=4\tau+1500,\qquad\ell_1:=30\tau,
\]
and the functions
\[
  u_0(\ell):=\ell^{1/4},\qquad U(\ell):=u_0+\ell^{1/4}\sqrt{2\log(\ell/\tau)+4},\qquad
  Q(\ell,u):=A+a,\quad A:=\frac{B(1+u)^3\log(2+u)}{\ell},\quad a:=\frac{(u-u_0)_+^2}{2\sqrt\ell},
\]
\[
  \Phi(\ell):=\Phi_\infty e^{-\kappa_\Phi/\ell},\qquad\Phi_\infty:=C_0\max\big(e^{\kappa_\Phi/\tau},\,30e^{\kappa_\Phi/\ell_1}\big).
\]
We prove, by induction on $\ell$, the claim
\begin{equation}\label{eq:bs-claim}
  \Psi(\ell,u)\le G(\ell,u):=\Phi(\ell)\,\rho(u-1)\,e^{Q(\ell,u)}\qquad(\ell\ge\tau,\ 0\le u\le U(\ell)).
\end{equation}
At $u=0$ (that is, $z=\infty$) this gives the theorem, since
$W_t(Y)=\Psi(\ell,0)\tau/\ell\le\Phi_\infty e^{(B\log2-\kappa_\Phi)/\ell}\,\tau/\ell\le\Phi_\infty\tau/\ell$ and
$\Phi_\infty=C_{\mathrm{BS}}$. The three factors of $G$ play the following roles:
\begin{itemize}
\item $A$ absorbs the errors of the prime sums;
\item $a$ inflates the profile gradually, by $\log(\ell/\tau)+2$ at $u=U$; this pays for the density factor $\ell/\tau$ in
  the tail (Step~6);
\item $\Phi$ absorbs the boundary term $\log(1/w_1)\approx2\tau/\ell$ of the critical zone $u<2$. The ratio is exact:
  $\Phi(m)/\Phi(\ell)=e^{-\kappa_\Phi/(w\ell)}$, since $1/m-1/\ell=1/(w\ell)$.
\end{itemize}

\emph{Step 4: base region $\tau\le\ell\le\ell_1$.} Drop the density condition, which costs at most a factor
$\max(1,\ell/\tau)$, and apply Lemma~\ref{lem:smooth-v}. Its hypotheses hold: $u\le U(\ell)\le4.3\ell^{1/4}$, so
$L_z=\ell/u\ge\ell^{3/4}/4.3\ge100$ and $u\le L_z$. This gives $\Psi\le(\ell/\tau)C_0\rho(u)e^{a_V}\le G(\ell,u)$, because
$\rho(u)\le\rho(u-1)$, $a_V\le A$ (as $11(1+u)^3\ge12.8u^2$), and $(\ell/\tau)C_0\le\Phi_\infty e^{-\kappa_\Phi/\ell}$ on
$[\tau,\ell_1]$. The last inequality holds by the choice of $\Phi_\infty$: the function $(\ell/\tau)e^{\kappa_\Phi/\ell}$ is
convex in $\ell$, so its maximum on $[\tau,\ell_1]$ is at an endpoint.

\emph{Step 5: the induction step, discretised.} Let $\ell>\ell_1$ and $0\le u\le U(\ell)$, and assume
\eqref{eq:bs-claim} for all smaller $\ell$. Split the primes in \eqref{eq:bs-rec-log} into the range $w\le U(m)$, where
the induction hypothesis applies, and the tail $w>U(m)$ (Step~6). The range $w_{\mathrm{lo}}\le w\le U(m)$ is an interval
$\sigma\in[\sigma_J,\sigma_a]$, with $\sigma_a:=\ell/(1+w_{\mathrm{lo}})$. On it put
\[
  f(\sigma):=\nu\,G(m,w)=f_1f_2,\qquad f_1:=\nu\,\rho(w-1),\qquad f_2:=\Phi(m)e^{Q(m,w)} .
\]
Here $f_1$ is increasing in $\sigma$ (both $\nu$ and $\rho(w-1)$ increase as $w$ decreases), and
\[
  |\partial_\sigma\log f_2|\le\eta_f(m,w):=\frac{\kappa_\Phi}{m^2}+|\partial_mQ|+\frac{(1+w)^2}{\ell}|\partial_wQ|,
\]
since $dm/d\sigma=-1$ and $dw/d\sigma=-(1+w)^2/\ell$. Write $\sum_{p\le e^\sigma}1/p=\log\sigma+M_z+E(\sigma)$ with
$|E(\sigma)|\le K'/\sigma$ (Remark~\ref{rem:cm-chain}). Partial summation gives
$\sum_pf(\sigma_p)/p\le\int f\,d\sigma/\sigma+K'[f(\sigma_a)/\sigma_a+f(\sigma_J)/\sigma_J+\int|f'|\,d\sigma/\sigma]$.
Using $|f'|\le f'+2f\,|\partial_\sigma\log f_2|$ and integrating $\int f'\,d\sigma/\sigma$ by parts, the terms in
$f(\sigma_J)$ cancel:
\begin{equation}\label{eq:bs-D1}
  \sum_p\frac{f(\sigma_p)}p\le\int f\frac{d\sigma}\sigma+\frac{2K'f(\sigma_a)}{\sigma_a}+K'\int f\frac{d\sigma}{\sigma^2}
  +2K'\int f\,|\partial_\sigma\log f_2|\frac{d\sigma}\sigma .
\end{equation}
In $w$-variables $d\sigma/\sigma=dw/(1+w)$, so $\int f\,d\sigma/\sigma=\int G(m,w)\,dw/w$ and
$\int f\,d\sigma/\sigma^2=\ell^{-1}\int G(m,w)(1+w)\,dw/w$.

Now divide by $G(\ell,u)$. Put
\[
  D(w):=Q(m,w)-Q(\ell,u)-\frac{\kappa_\Phi}{w\ell},\qquad
  \pi(dw):=\frac{\rho(w-1)\,dw}{w\,\rho(u-1)}\ \text{on }[w_{\mathrm{lo}},\infty),\qquad
  g(w):=\frac{\rho(w-1)}{\rho(u-1)}e^{D(w)}\mathbf 1[w\le U(m)] .
\]
The main term is $\int\pi\,e^D\mathbf 1[w\le U(m)]$. Since $e^D>0$, it is at most $\int_{w_{\mathrm{lo}}}^\infty\pi\,e^D=1+\mu_0+J$,
where
\begin{itemize}
\item $\mu_0:=\log(1/w_{\mathrm{lo}})$ for $u<2$, and $\mu_0:=0$ for $u\ge2$, is the exact excess mass of $\pi$, by \ref{R1};
\item $J:=\int_{w_{\mathrm{lo}}}^\infty\pi\,(e^D-1)$.
\end{itemize}
In the same way, any positive majorant of $e^D\mathbf 1[w\le U(m)]$ may be integrated over the full range $[w_{\mathrm{lo}},\infty)$.
The three error terms are
\[
  E_1:=\frac{2K'(1+w_{\mathrm{lo}})^2g(w_{\mathrm{lo}})}{w_{\mathrm{lo}}\,\ell},\qquad
  E_2:=\frac{K'}\ell\int g\,\frac{(1+w)\,dw}w,\qquad E_3:=2K'\int g\,\eta_f\,\frac{dw}w ,
\]
and $T$ denotes the tail of Step~6, divided by $G(\ell,u)$. The induction closes if
\begin{equation}\label{eq:bs-Sigma}
  \Sigma(\ell,u):=-\big[\mu_0+J+E_1+E_2+E_3+T\big]\ge0\qquad\text{for all }\ell>\ell_1,\ 0\le u\le U(\ell).
\end{equation}

\emph{Step 6: the tail.} The primes with $w>U(m)$ correspond exactly to the $n$ that are $e^{\ell/u^*}$-smooth, with
$u^*:=U(\ell)+\frac12$. This uses $U(m)\ge U(\ell)-\frac12$ for $m\ge\ell(1-1/U)$. Dropping the density condition and
applying Lemma~\ref{lem:smooth-v} gives a tail contribution to $\Psi$ of at most $(\ell/\tau)C_0\rho(u^*)e^{a_V(u^*)}$.
Divide by $G(\ell,u)$, and use $\Phi_\infty\ge C_0$, $Q\ge a$ and
$\rho(u^*)/\rho(u-1)\le e^{-r(u-1)(u^*-u+1)}$ \ref{R3}. With $y:=\log(\ell/\tau)\ge\log30$,
\[
  T\le\exp\big(\kappa_\Phi/\ell+a_V(u^*)+y-a(\ell,u)-r(u-1)(U(\ell)-u+\tfrac32)\big).
\]
Maximising over the admissible $y$ gives $T\le\exp(\kappa_{\max}/\ell+a_V(u^*)+E_T)$, where
$\kappa_{\max}:=4\ell/30+1500\ge\kappa_\Phi$ and, with $U_{30}(\ell):=u_0+\ell^{1/4}\sqrt{2\log30+4}$,
\[
  E_T:=\begin{cases}\log30-a(\ell,u)-r(u-1)\big(U_{30}-u+\tfrac32\big),&u\le U_{30},\\
  -2-1.5\,r(u-1),&u>U_{30}.\end{cases}
\]
In the second case this uses $a(\ell,U(\ell))=y+2$.

\emph{Step 7: verification of \eqref{eq:bs-Sigma}.} There are three regions. Throughout, $\tau\ge10^8$ and $\ell>\ell_1=30\tau$.

\smallskip\noindent\emph{(i) Region Z1: $u<2$, every $\ell\ge30\tau$.} Here $\rho(u-1)=1$ and $\pi(dw)=\rho(w-1)\,dw/w$. Since
$u_0=\ell^{1/4}\ge234$, $D_a:=a(m,w)-a(\ell,u)=0$ for $w\le40$, so $D=D_A+D_\Phi$ with $D_A:=A(m,w)-A(\ell,u)$ and
$D_\Phi:=-\kappa_\Phi/(w\ell)\le0$; and $e^{D}-1\le(e^{D_A}-1)_++(e^{D_\Phi}-1)$.
\begin{itemize}
\item \emph{The $A$-part.} With $h(w):=(1+w)^4\log(2+w)/w$ one has $A(m,w)=(B/\ell)h(w)$ exactly, so
  $D_A\le(B/\ell)h(w)\le10^{-3}$ for $w\le40$ and $e^{D_A}-1\le1.001D_A$ there. On $[w_{\mathrm{lo}},2]$, $\rho(w-1)=1$ and
  $h\le h(2)=56.145$. On $[2,\infty)$, $\rho(w-1)\le e^{-(w-2)}$ by \ref{R3} with $r(1)=1$, and $h(w)/w\le(3+s)^4(\log4+s/4)/4$
  with $s=w-2$. The part $w>40$ contributes at most $10^{-40}$: there $\rho(w-1)\le\rho(39)e^{-r(39)(w-40)}$ with
  $r(39)\ge4.4$, while the slope of $D$ is at most $0.25$. So
  $\ell\,E_\pi[(e^{D_A}-1)_+]\le1.001\,B\,(56.145\log(2/w_{\mathrm{lo}})+J_2)$, with
  $J_2:=\int_0^\infty e^{-s}(3+s)^4(\log4+s/4)\,ds/4\le200.6$ (by the Gamma moments).
\item \emph{The $\Phi$-benefit.} Put $\Phi_{\mathrm{ben}}:=E_\pi[1-e^{-\kappa_\Phi/(w\ell)}]$. Using
  $1-e^{-y}\ge y(1-y/2)$, $\kappa_\Phi/\ell\le0.1334$ and $w\ge w_1\ge29/31-10^{-8}$,
  $\Phi_{\mathrm{ben}}\ge(\kappa_\Phi/\ell)(1-\kappa_\Phi/(2w_1\ell))E_\pi[1/w]\ge0.9287\,(\kappa_\Phi/\ell)\,c_2^-$, where
  \[
    E_\pi[1/w]\ge c_2^-:=\int_1^2\frac{dw}{w^2}+\int_2^3\frac{(1-\log(w-1))\,dw}{w^2}=\frac23+\frac{\log2}3-\log\frac43=0.6100 .
  \]
\item \emph{The excess mass.} With $x':=(\tau+\log2)/\ell\le1/30+10^{-17}$, $\mu_0=\log(1/w_{\mathrm{lo}})\le\log(1/w_1)\le2x'+x'^3$,
  so $\ell\mu_0\le2.0012\,\tau+1.4$.
\item The remaining terms satisfy $\ell(E_1+E_2+E_3+T)\le33$.
\end{itemize}
Altogether, since $\kappa_\Phi\ge4\tau$ and $\log(2/w_{\mathrm{lo}})\le0.76$,
\[
  \ell\Sigma\ge\tau\,(4\cdot0.9287\,c_2^--2.0012)-1.001\,B\,(56.145\cdot0.76+J_2)-33-1.4\ge0.265\,\tau-2714>0 .
\]

\smallskip\noindent\emph{(ii) Large $\ell$: $\ell\ge10^{18}$, every $u\ge2$.} This region is analytic, apart from a
one-variable inequality checked in ball arithmetic. Put
\begin{gather*}
  s:=w-u,\qquad r_0:=r(u-1),\qquad\Lambda(u):=(1+u)^2\log(2+u),\qquad h_A(u,w):=h(w)-(1+u)^3\log(2+u),\\
  D_A=(B/\ell)\,h_A,\qquad D_a:=a(m,w)-a(\ell,u),\qquad D_\Phi:=-\kappa_\Phi/(w\ell),\qquad D':=D_A+D_a .
\end{gather*}
Write $E_\pi$ for integration against $\pi$, whose total mass is $1+\mu_0\le1.7$. For $s\ge0$ the density of $\pi$ is at most
$e^{-r_0s}/u$, by \ref{R3}. The \emph{first-order profile} is
\[
  c(u):=\ell\big(\Phi_{\mathrm{ben}}-\mu_0\big)-B\,j_A(u)-2K'e_1(u)-K'e_2(u),
\]
where
\begin{itemize}
\item $\Phi_{\mathrm{ben}}:=E_\pi[1-e^{-\kappa_\Phi/(w\ell)}]$ (exact, not linearised);
\item $j_A:=E_\pi h_A$;
\item $e_1:=(1+w_{\mathrm{lo}})^2\rho(w_{\mathrm{lo}}-1)/(w_{\mathrm{lo}}\rho(u-1))$;
\item $e_2:=\int(1+w)\rho(w-1)\,dw/(w\rho(u-1))$.
\end{itemize}
For $u\ge2$, $c$ depends on $\ell$ only through $\Phi_{\mathrm{ben}}$.

\begin{lemma}[remainder, in ratio form]\label{lem:bs-remainder}
% v3: Theorem 3.4, proof (g), case l >= 10^18
For $\tau\ge10^8$, $\ell\ge\max(10^{18},30\tau)$ and $0\le u\le U(\ell)$,
\[
  \ell\Sigma\ge\big(1-\eps(\ell)\big)\big(c(u)+\ell s_a\big)-60\cdot\mathbf 1[u\ge u_0-2]-9\cdot10^{17}/\ell,
\]
where $s_a:=-E_\pi D_a\ge0$ is the slack of $a$, $a'(\ell,u):=\partial_ua=(u-u_0)_+/\sqrt\ell$, and
\[
  \eps(\ell):=2\cdot10^4\,U^2\log(2+U)/\ell+30\,a'(\ell,U)+\ell^{-1/4}\le0.03 ,
\]
with $U=U(\ell)$; $\eps$ is decreasing in $\ell$.
\end{lemma}

\begin{proof}
Write $e^D-1=(e^{D_\Phi}-1)+e^{D_\Phi}(e^{D'}-1)$, and use $e^{D'}-1\le D'+\frac12D'^2e^{D'_+}$. This gives
\[
  \ell\Sigma\ge c(u)+\ell s_a-\ell\big(\rho_1+\rho_2+\rho_3+\rho_4+E_3+T\big),
\]
where $\rho_1:=E_\pi[(1-e^{D_\Phi})|D'|]$, $\rho_2:=\frac12E_\pi[D'^2e^{D'_+}]$, $\rho_3:=E_1-E_1^{\mathrm{lin}}$ and
$\rho_4:=E_2-E_2^{\mathrm{lin}}$. Here $E_i^{\mathrm{lin}}$ denotes $E_i$ with $e^{D}$ replaced by $1$ in $g$, so that
$\ell E_1^{\mathrm{lin}}=2K'e_1$ and $\ell E_2^{\mathrm{lin}}=K'e_2$; and $\rho_3=E_1^{\mathrm{lin}}(e^{D'(w_{\mathrm{lo}})}-1)$. We use
the moments $M_k:=E_\pi|s|^k\le1.7\cdot2^k+e\,k!/2$, valid for $r_0\ge1$. The bounds below are for $u\ge2$; the case
$u<2$ is identical with $\Lambda(2)$ in place of $\Lambda(u)$.
\begin{enumerate}[label=(\alph*)]
\item $|h_A|\le\Lambda(u)\big(1.5+8|s|(1+|s|)^3\big)$. This uses
  $h(u)-(1+u)^3\log(2+u)=(1+u)^3\log(2+u)/u$ and $h'(w)\le8(1+w)^2\log(2+w)$ for $w\ge1$. For $|s|\le1$ and $u\ge27$ the
  bound is at most $10.3\Lambda(u)$. Hence, for $u\ge27$ (where $r_0\ge4.3$), $E_\pi|h_A|\le18\Lambda$ and
  $E_\pi h_A^2\le182\Lambda^2$. These bounds are in ratio form: the weight $e^{-r_0s}$ enters, not $\sup_w$.
\item $\rho_1\le6\,\Phi_{\mathrm{ben}}\,E_\pi|D'|$, because $(1-e^{-y})/y$ is decreasing, $(u+1)/(u-1)\le3$, and
  $\pi(w\le u+1)\ge\frac12$.
\item The $a$-part. For $u\ge u_0-2$: $r_0\ge15$, and $s_a\ge0.32\,a'(\ell,u)-1.01/\sqrt\ell$. The proof combines the
  chord bound on $[u-1,u]$, the identity $E_\pi(u-w)_+=1-1/r(u)$ from \ref{R1}, and the shift bound
  $(w-m^{1/4})_+\le(w-u_0)_++u_0/(3(1+w))$. Also
  $\frac12E_\pi D_a^2\le3a'(\ell,u)^2M_2+0.75M_4/\ell+0.01/\ell\le10\,a'(\ell,U)\,s_a+60/\ell$. For $u<u_0-2$:
  $\ell|E_\pi D_a|\le1.02\sqrt\ell\,e^{-r_0(u_0-u-2/3)}/(u\,r_0^3)$.
\item $\rho_3\le\ell E_1^{\mathrm{lin}}\cdot2\big(128B\Lambda/\ell+3a'(\ell,U)\big)$ and
  $\rho_4\le\ell E_2^{\mathrm{lin}}\cdot2\big(8B\Lambda M_4/\ell+3a'(\ell,U)\big)$, with $\ell E_1^{\mathrm{lin}}=2K'u^2r_0$ for
  $u\ge2$. Also
  $\ell E_3\le2K'\big[14\kappa_\Phi/\ell+12B(1+u)^4\log(2+u)M_4/\ell+4(1+u)^2M_2\,a'(\ell,U)\big]$; its
  $\kappa_\Phi$-piece is at most $10^{-3}\,\ell\Phi_{\mathrm{ben}}$.
\end{enumerate}
For $u\ge27$, every remainder is a small multiple of $c(u)\ge K'u^2\log(1+u)$ (proved below) or of $\ell\Phi_{\mathrm{ben}}$. For
$u<27$ the remainders are absolute, at most $9\cdot10^{17}/\ell$. On the range $u\ge u_0-2$,
$c(u)\ge K'(u_0-2)^2\log u_0\ge2\cdot10^9$, so the term $-60$ is negligible. Finally $\eps(10^{18})\le0.03$,
$\eps(10^{22})\le0.003$, and $\eps$ is decreasing. The constants are crude, and nothing is optimised.
\end{proof}

\emph{Positivity of $c(u)$, for $\tau\ge10^8$.}
\begin{itemize}
\item \emph{$u\ge27$: $c(u)\ge K'u^2\log(1+u)$.} Here $\mu_0=0$ and $\Phi_{\mathrm{ben}}\ge0$.
  \begin{itemize}
  \item On $[u-1,u]$, convexity of $h$ ($h''>0$ on $[1,\infty)$) gives the chord bound $h(w)\le h(u)-h'(u-1)(u-w)$, with
    $h'(u-1)\ge3u^2\log(1+u)$, and \ref{R1} gives $E_\pi(u-w)_+=1-1/r(u)$.
  \item On $w\ge u$, $h(w)-h(u)\le h'(w)s$ and \ref{R3} bound the remaining part by $144\Lambda(u)/(ur_0^2)$.
  \item Hence $E_\pi h_A\le(1+u)^3\log(2+u)/u-3u^2\log(1+u)(1-1/r(u))+144\Lambda(u)/(ur_0^2)$. Also $e_1=u^2r(u-1)$ and
    $e_2\le u+1$.
  \end{itemize}
  With $B=5K'$ this gives
  \begin{multline*}
    \frac{c(u)}{K'u^2}\ge5\log(1+u)\Big[3\big(1-1/r(u)\big)-\frac{(1+u)^3\log(2+u)}{u^3\log(1+u)}
    -\frac{144(1+u)^2\log(2+u)}{u^3r_0^2\log(1+u)}\Big]\\-2r(u-1)-\frac{u+1}{u^2}.
  \end{multline*}
  With $r$ enclosed in $[r_{\mathrm{low}},r_{\mathrm{up}}]$, the right side is $\ge\log(1+u)$ on $[27,4\cdot10^5]$
  (computer-assisted; see Appendix~\ref{app:cert}: $9{,}609$ cells in ball arithmetic, no failures, minimum margin
  $0.189$). Beyond $4\cdot10^5$ it is analytic: $r_{\mathrm{low}}\ge15.3$ gives
  $5\cdot1.79\log u-2(\log u+\log\log u+2)>\log(1+u)$.
\item \emph{$2\le u<27$: $\Phi$-dominance.} Here
  $\ell\Phi_{\mathrm{ben}}\ge\kappa_\Phi\cdot0.933\cdot0.69/u\ge4\cdot10^8\cdot0.933\cdot0.69/u\ge9.5\cdot10^6$, while
  $B\,j_A\le B(1.5+393.5/u)\Lambda(u)\le5\cdot10^5$ and $2K'u^2r_{\mathrm{up}}\le2\cdot10^4$.
\item \emph{$u<2$:} region Z1, case (i).
\end{itemize}
So, by Lemma~\ref{lem:bs-remainder}, $\Sigma\ge0$ for every $\ell\ge10^{18}$ and every $u\le U(\ell)$.

\smallskip\noindent\emph{(iii) Region Z2: $30\tau\le\ell<10^{18}$ and $u\ge2$ (computer-assisted; see Appendix~\ref{app:cert}).}
Here we use a closed-form lower bound $S(\ell,u)\le\Sigma(\ell,u)$, built only from \ref{R1}--\ref{R4}, the constant $K'$ of
Lemma~\ref{lem:cm}, and convexity. Since $u\ge2$, $w_{\mathrm{lo}}=u-1$ and $\mu_0=0$. Put $v:=u-u_0$,
$c_u:=u_0/(3u)$, $k_u:=1+1/(2(u-1))$ and $\varphi(w):=(w-u_0+c_u)_+^2/(2\sqrt\ell)$.
\begin{itemize}
\item \emph{Two elementary majorants.} First, $1/(2\sqrt m)=\sqrt{(1+w)/w}/(2\sqrt\ell)\le k_u/(2\sqrt\ell)$ for $w\ge u-1$.
  Second, $(w-m^{1/4})_+\le(w-u_0+u_0/(3(1+w)))_+\le(w-u_0+c_u)_+$. Together, $a(m,w)\le k_u\varphi(w)$, and $\varphi$ is
  convex. Also $A(m,w)=(B/\ell)h(w)$ with $h$ convex.
\item \emph{Part I, $w\in[u-1,u]$.} By the chord bounds for $h$ and $\varphi$, together with $h(u)-h(u-1)\ge h'(u-1)$,
  $\varphi(u)-\varphi(u-1)\ge\varphi'(u-1)$, and $D_\Phi\le-\kappa_{\min}/(u\ell)$, we get $D\le\beta_I-\alpha(u-w)$, where:
  \begin{itemize}
  \item $\beta_I:=A(\ell,u)/u+\beta_a-\kappa_{\min}/(u\ell)$;
  \item $\alpha:=(B/\ell)u^3(3u-4)\log(1+u)/(u-1)^2+k_u(v-1+c_u)_+/\sqrt\ell$;
  \item $\beta_a:=(k_u-1)(v+c_u)_+^2/(2\sqrt\ell)+c_u(2v_++c_u)/(2\sqrt\ell)$.
  \end{itemize}
  Then use $e^{-\alpha t}\le1-(1-e^{-\alpha})t$ and $E_\pi(u-w)_+=1-1/r(u)$ \ref{R1}.
\item \emph{Part II, $w\in[u,\bar U]$, $s=w-u$.} Here $D\le\beta_{II}+\gamma(s)s$, with $\beta_{II}:=A(\ell,u)/u+\beta_a$ and a
  two-piece slope: $\gamma(s)=\gamma_1$ for $s\le60$ and $\gamma(s)=\gamma_{\max}$ beyond. These are the slopes of
  $(B/\ell)h+k_u\varphi$ bounded via $h'(w)\le8(1+w)^2\log(2+w)$ at $w=u+60$ and at $w=\bar U$, respectively. This is
  combined with $\rho(w-1)/\rho(u-1)\le e^{-r(u-1)s}$ \ref{R3}. The integrals converge because $r(u-1)>\gamma_1$ and
  $r(u-1)>\gamma_{\max}$, which is checked in every cell.
\item \emph{What the induction needs.} It needs an upper bound for $\int\pi\,e^D\mathbf 1[w\le U(m)]-1$. The Part I/II
  functions are positive majorants of $e^D\mathbf 1[w\le U(m)]$ on $[u-1,\bar U]$, and $\mathbf 1[w\le U(m)]=0$ beyond
  $\bar U$. So their integrals bound it. The same majorants (with a two-piece bound for $\eta_f$) bound $E_1$, $E_2$ and
  $E_3$; $T$ is bounded as in Step~6. The difference $M-1$ is written as a sum of small terms, so there is no cancellation
  against~$1$.
\item \emph{Uniformity in $\tau\ge10^8$.} The bound uses $\kappa_{\min}:=4\cdot10^8$ where $\kappa_\Phi$ helps,
  $\kappa_{\max}=4\ell/30+1500$ where it hurts, and the slopes at $\bar U(\ell):=U(\ell;\tau=10^8)$, which is the largest
  junction.
\item \emph{The cover.} The range $\ell\in[3\cdot10^9,10^{18}]$ is split into $403$ bands of ratio $1.05$. In each band,
  $u$-cells cover $[2,\max(40,u_0/2)]$, and cells in $v=u-u_0(\ell)$ cover the rest up to $\bar U(\ell)$. The endpoints of
  the $v$-cells are rounded outward, so the seams between the two kinds of cells, and the top edge, are covered exactly
  for every $\ell$ in the band. On each cell, $S$ is evaluated in ball arithmetic, with $r_{\mathrm{low}}$ and $r_{\mathrm{up}}$
  computed in the same arithmetic. The enclosure bounds $S$ on the whole cell, so no Lipschitz or interpolation argument is
  needed.
\item \emph{Result.} $302{,}013$ cells (including adaptive bisections), no failures, and minimum certified $\ell S=2.491\cdot10^6$.
\end{itemize}

\emph{Step 8: conclusion.} By Steps 4--7, \eqref{eq:bs-claim} holds for all $\ell\ge\tau$ and $0\le u\le U(\ell)$, and
Step~3 gives the theorem. The threshold $\ell_1=30\tau$ is internal to the induction: the base region $[\tau,30\tau]$ is
covered by Lemma~\ref{lem:smooth-v}. So the conclusion holds for every $Y\ge t$, with the same $C_{\mathrm{BS}}$. For the
$4t$-dense version, replace $\tau$ by $\tau+\log4$ throughout; nothing else changes.
\end{proof}

The induction claim \eqref{eq:bs-claim} is itself used later, in the proof of the short-interval Saias bound
(Lemma~\ref{lem:si}). We record it.

\begin{lemma}[BS-profile]\label{lem:bs-profile}
% v3: Lemma 3.4A
Let $\tau=\log t\ge10^8$, $\ell\ge\tau$ and $0\le u\le U(\ell)$. Then
\[
  W_t(e^\ell;e^{\ell/u}):=\sum_{n\in\NN\cap(e^\ell,2e^\ell],\ t\text{-dense},\ \Pplus(n)<e^{\ell/u}}\frac1n
  \le\frac\tau\ell\,\Phi_\infty\,\rho(u-1)\,e^{Q(\ell,u)},
\]
with $U$, $Q$ and $\Phi_\infty\le38.3$ as in Step~3 of the proof of Theorem~\ref{thm:bs}. Moreover, $\Phi_\infty$ is
decreasing in $\tau$.
\end{lemma}

\begin{proof}
This is \eqref{eq:bs-claim}, divided by $\max(1,\ell/\tau)=\ell/\tau$. Steps 4--7 prove it on the whole range
$0\le u\le U(\ell)$, not only at $u=0$.

\emph{Uniformity in $\tau\ge10^8$.} The computer-assisted region is uniform in $\tau$ by construction ($\kappa_{\min}$,
$\kappa_{\max}$, $\bar U$). The analytic regions are monotone in $\tau$ in the safe direction: $\ell\Sigma\ge0.265\tau-2714$ for
$u<2$, and for $\ell\ge10^{18}$ the term $\ell\Phi_{\mathrm{ben}}$ increases with $\kappa_\Phi$. The base region uses
Lemma~\ref{lem:smooth-v} with $\Phi_\infty$ chosen for the given $\tau$.
\end{proof}

\begin{corollary}[Saias weights at every start]\label{cor:saias}
% v3: Corollary 3.5
For every start $y\ge X_0$ and every $\Xm/2\le v\le y$,
\[
  W^{\mathrm{sh}}_y(v):=\sum_{e\in\NN:\ y/e\in[v,2v)}\frac{|c_e(y)|}e\le C_S\,\frac{\log(2v)}{\log y},\qquad
  C_S:=\max(2.01\,C_{\mathrm{BS}},\,1.4)\le77 .
\]
\end{corollary}

\begin{proof}
If $c_e(y)\ne0$, then $y/e\ge\Xm$ (the cutoff in $\Kpt$), and $e$ is squarefree and $4v$-dense (Lemma~\ref{lem:lemmaM}).
\begin{itemize}
\item For $v\le\sqrt y/4$, apply Theorem~\ref{thm:bs} with $t:=4v\ge2\Xm$ and $Y:=y/(2v)$. Then $Y\ge t$ and
  $\log Y\ge\frac12\log y$, and
  \[
    W^{\mathrm{sh}}_y(v)\le C_{\mathrm{BS}}\frac{\log(4v)}{\log(y/2v)}\le2C_{\mathrm{BS}}\Big(1+\frac{\log2}{\log X_0}\Big)
    \frac{\log(2v)}{\log y}\le2.01\,C_{\mathrm{BS}}\frac{\log(2v)}{\log y}.
  \]
\item For $v>\sqrt y/4$, put $A:=y/(2v)\in[\frac12,2\sqrt y)$. Then $W^{\mathrm{sh}}_y(v)\le\sum_{e\in(A,2A]}1/e$.
  This sum is at most $1$ always (on $\N$ below $\Xm$), and at most $\log2+2\eta_1+1/A\le0.6952$ if $A\ge y^{1/4}$. In the
  first case $\log(2v)/\log y\ge\frac34$, and in the second $\log(2v)/\log y\ge\frac12-\log2/\log y$. Either way
  $W^{\mathrm{sh}}_y(v)\le1.4\log(2v)/\log y$.\qedhere
\end{itemize}
\end{proof}

% 04k-kernels.tex -- v3 §3.4: the operator kernels (Kantorovich--Rubinstein transfer of Lemma M).
% Author: fix-bprime. Input from 04-noise.tex. Contents: kernels and steps, KR facts, Theorem 3.6, Corollary 3.7,
% Proposition 3.8, Lemma 3.9, Lemma 3.10.
\subsection{The operator kernels}\label{sec:kernels}
% v3: §3.4

The long operator $\Rres=(1+\Kpt B)^{-1}$ of \eqref{eq:pushthrough} differs from the sharp resolvent $(1+\Kpt)^{-1}$ in two
ways. After every latest-prime step it averages over a window of length $H$ at the landing scale, and it evaluates the
filters at the actual (averaged) centre, not at the nominal landing point. This subsection shows that the first
difference costs only a constant factor in the Saias weights of Corollary~\ref{cor:saias}, and that the second is
supported on near-threshold steps, which are rare. The main tool is the Kantorovich--Rubinstein distance: it contracts
under dilation by the factor $d/n$ at every step, so that, unlike total variation, it can be iterated along chains of
arbitrary length.

\subsubsection{Kernels and steps}\label{sec:kernels-steps}

Fix a start $x$. For a cofactor $w\in\NN$, $w\ge2$, and a position $v$, define the \emph{step}
\begin{equation}\label{eq:step}
  Q_w(v,\cdot):=\mathbf 1[\tau_w<v/w]\,\mathbf 1[v/w\ge\Xm]\;\bx_{H(v/w)}(\cdot-v/w),
\end{equation}
where $\bx_h$ is the uniform probability measure on $[-h/2,h/2]$. Thus the filter is evaluated at the actual centre
$v/w$, and the window average comes after the filter. By the definitions of $\Kpt$ and $B$ in
\S\ref{sec:noise-notation}, for every function $g$,
\[
  (\Kpt Bg)(v)=\sum_{w\ge2}\frac1w\,\mathbf 1[\tau_w<v/w]\,\mathbf 1[v/w\ge\Xm]\,(Bg)(v/w)=\sum_{w\ge2}\frac1w\int g\,dQ_w(v,\cdot).
\]

\begin{definition}[Kernels]\label{def:kernels}
% v3: §3.4, Kernels
The kernels of $\Rres$ from the start $x$ are the signed measures $C_n=C_n^x$ ($n\in\NN$) given by the last-step recursion
\begin{equation}\label{eq:kernel-rec}
  C_1^x=\delta_x,\qquad C_n^x=-\sum_{d\mid n,\;d<n}\int Q_{n/d}(v,\cdot)\,dC_d^x(v)\quad(n\ge2).
\end{equation}
\end{definition}

Expanding $\Rres=\sum_{k\ge0}(-\Kpt B)^k$ (the series terminates at every point, by the remark after
\eqref{eq:pushthrough}) and grouping the chains $w_1,\dots,w_k$ according to their product $n=w_1\cdots w_k$ gives
\begin{equation}\label{eq:R-kernels}
  (\Rres g)(x)=\sum_{n\in\NN}\frac1n\int g\,dC_n^x .
\end{equation}
Each $C_n^x$ is supported in an interval of radius $R_n\le1.02\,H(x/n)$ around $x/n$. Indeed, a step $d\to n$ maps
the support of $C_d$, of radius $R_d$ around $x/d$, to radius $R_d/w$ around $x/n$ and then adds half a window, so by
induction and \eqref{eq:RH}, $R_n\le H(x/n)/2+\max_d 1.02\,H(x/d)/w\le H(x/n)\,(\tfrac12+1.02\cdot2^{\eps-1})\le1.02\,H(x/n)$.

\paragraph{Nominal steps.} Write $H_n:=H(x/n)$. For a step $d\to n$ with cofactor $w=n/d$, the \emph{nominal} step
\[
  Q^{\nom}_{d\to n}(v,\cdot):=\kappa_{d\to n}\,\bx_{H(v/w)}(\cdot-v/w),\qquad
  \kappa_{d\to n}:=\mathbf 1[\tau_w<x/n]\,\mathbf 1[x/n\ge\Xm],
\]
freezes the scalar filter at the nominal landing point $x/n$; the window width stays $H(v/w)$. The nominal kernels
$C_n^{x,\nom}$ are defined by \eqref{eq:kernel-rec} with $Q^{\nom}$ in place of $Q$. Their masses
$c_n:=C_n^{x,\nom}(\R)$ satisfy $c_1=1$ and $c_n=-\sum_{d\mid n,d<n}\kappa_{d\to n}c_d$, which is the recursion of the
sharp resolvent. Hence, by Lemma~\ref{lem:lemmaM},
\begin{equation}\label{eq:nominal-mass}
  c_n=\mu(n)\chi^t_x(n).
\end{equation}

\begin{definition}[Near-threshold steps and terminal levels]\label{def:near}
% v3: §3.4, τ-near / X₀-near / terminal
A step $d\to n$ ($w=n/d$) is \emph{$\tau$-near} if $|\tau_w-x/n|\le2H_n$, and \emph{$X_0$-near} if
$|x/n-\Xm|\le R_d/w$. A level $n$ is \emph{terminal} if $x/n<\Xm+0.6H(X_0)$.
\end{definition}

\begin{lemma}[Structure of the near-threshold steps]\label{lem:near-structure}
% v3: §3.4, bullets on near-threshold steps and terminal levels
\begin{enumerate}[label=(\roman*)]
\item On a step that is neither $\tau$-near nor $X_0$-near, $Q_w(v,\cdot)=Q^{\nom}_{d\to n}(v,\cdot)$ for every $v$ in the
  support of $C_d$.
\item $R_d/w\le0.52\,H(X_0)$ for every $X_0$-near step, and every $X_0$-near step lands at a terminal level.
\item A terminal level has no children in either family: $Q_{w'}(v,\cdot)=Q^{\nom}_{n\to nw'}(v,\cdot)=0$ for every $w'\ge2$
  and every $v$ in the support of $C_n$.
\end{enumerate}
\end{lemma}

\begin{proof}
(i) On the support of $C_d$, the centre is $v/w\in[x/n-R_d/w,\,x/n+R_d/w]$, and $R_d/w\le1.02H(x/d)/w\le2H_n$ by
\eqref{eq:RH}. Off $\tau$-near steps the filter $\mathbf 1[\tau_w<v/w]$ is therefore constant there and equal to
$\mathbf 1[\tau_w<x/n]$; off $X_0$-near steps the same holds for $\mathbf 1[v/w\ge\Xm]$.

(ii) By \eqref{eq:RH}, $R_d/w\le1.02\,H(w\cdot x/n)/w\le1.02\cdot2^{\eps-1}H(x/n)\le0.52\,H(x/n)$. If the step is
$X_0$-near, then $x/n\le\Xm+0.52H(x/n)$. Since $H(v)/v$ is decreasing and $H(X_0)/X_0<10^{-10}$, this forces $x/n<X_0$,
so $H(x/n)\le H(X_0)$, $R_d/w\le0.52H(X_0)$ and $x/n\le\Xm+0.52H(X_0)$: the level $n$ is terminal.

(iii) The support of $C_n$ lies in $[x/n-R_n,x/n+R_n]$ with $R_n\le1.02H_n$, so every child centre is at most
$(\Xm+0.6H(X_0)+1.03H(X_0))/2<\Xm$; both filters $\mathbf 1[v/w'\ge\Xm]$ and $\mathbf 1[x/(nw')\ge\Xm]$ vanish.
\end{proof}

\paragraph{Width variation.} Put $\vartheta_u:=\sup_{v\ge u/2}H'(v)$. By \eqref{eq:RH} and the bound
$d\log H/d\log v\le\max(16\log\log v/\log v,\,4\cdot10^{-6})\le4\cdot10^{-6}$ for $v\ge X_0/2$,
$\vartheta_u\le10^{-5}H(u)/u$. Along any chain the nominal landing scales $u_i$ decrease at least geometrically (by a
factor $\ge2$) and stay $\ge\Xm$, and $H(u)/u$ is decreasing, so
\begin{equation}\label{eq:width-product}
  \prod_i(1+2\vartheta_{u_i})\le\exp\Bigl(2\sum_{j\ge0}\vartheta_{2^j\Xm}\Bigr)\le\exp\bigl(10^{-4}H(X_0)/X_0\bigr)
  \le\lambda_0:=\exp\bigl(2.1H(X_0)/X_0\bigr).
\end{equation}
Every position $v$ in the support of $C_d$ satisfies $v/w\ge(x/n)/2$, so $|H(v/w)-H_n|\le\vartheta_{x/n}|v/w-x/n|$.

\subsubsection{The Kantorovich--Rubinstein norm}\label{sec:kr}

For a finite signed measure $\nu$ on $\R$ with $\nu(\R)=0$ and compact support, put $F_\nu(t):=\nu((-\infty,t])$ and
\[
  \KRn{\nu}:=\int_\R|F_\nu(t)|\,dt=\sup\Bigl\{\int f\,d\nu:\ \operatorname{Lip}f\le1\Bigr\}.
\]
The two expressions agree: for Lipschitz $f$, $\int f\,d\nu=-\int f'(t)F_\nu(t)\,dt$ (integration by parts, $\nu(\R)=0$),
which is at most $\int|F_\nu|$, with equality for $f'=-\operatorname{sgn}F_\nu$. Let $D_w$ denote the push-forward under
$t\mapsto t/w$, and $\bx_h$ the uniform probability on $[-h/2,h/2]$.

\begin{lemma}[Kantorovich--Rubinstein facts]\label{lem:kr-facts}
% v3: §3.4, (K1)-(K5)
Let $\nu$ have mass $0$ and let $\rho$ be a finite signed measure. Then:
\begin{enumerate}[label=(K\arabic*)]
\item $\KRn{\rho*\nu}\le\TV(\rho)\,\KRn{\nu}$;
\item $\KRn{D_w\nu}=\KRn{\nu}/w$;
\item $\TV(\bx_h*\nu)\le(2/h)\,\KRn{\nu}$;
\item $\KRn{\rho-\rho(\R)\delta_a}\le\int|t-a|\,d|\rho|(t)$; in particular $\KRn{\bx_h-\delta_0}=h/4$;
\item multiplying a measure by a $[0,1]$-valued (or $[-1,1]$-valued) function does not increase its total variation.
\end{enumerate}
\end{lemma}

\begin{proof}
(K1) $F_{\rho*\nu}(t)=\int F_\nu(t-s)\,d\rho(s)$, so $\int|F_{\rho*\nu}|\le\iint|F_\nu(t-s)|\,dt\,d|\rho|(s)=\TV(\rho)\KRn{\nu}$.
(K2) $F_{D_w\nu}(t)=F_\nu(wt)$, and $\int|F_\nu(wt)|\,dt=\KRn{\nu}/w$.
(K3) $\bx_h*\nu$ has density $h^{-1}\bigl(F_\nu(t+h/2)-F_\nu(t-h/2)\bigr)$, whose $L^1$ norm is at most $(2/h)\int|F_\nu|$.
(K4) For $\operatorname{Lip}f\le1$, $\int f\,d(\rho-\rho(\R)\delta_a)=\int(f(t)-f(a))\,d\rho(t)\le\int|t-a|\,d|\rho|$. For
$\rho=\bx_h$ and $a=0$, $F(t)=t/h+\tfrac12-\mathbf 1[t\ge0]$ on $[-h/2,h/2]$, and $\int|F|=2\int_0^{h/2}(\tfrac12-t/h)\,dt=h/4$.
(K5) is immediate.
\end{proof}

\paragraph{Ordered factorisations.} Write $\Hord(m)$ for the number of ordered factorisations of $m\in\NN$ into factors
$\ge2$ ($\Hord(1)=1$), so that $\Hord(m)=\sum_{g\mid m,\,g<m}\Hord(g)$ for $m\ge2$, and put
\[
  A_0(\eps):=\sum_{m\ge2}\Hord(m)m^{\eps-2},\qquad A_1(\eps):=\sum_{m\ge2}\Hord(m)m^{\eps-2}\log m .
\]
For $\NN=\N$, $\sum_{m\ge1}\Hord(m)m^{-s}=1/(2-\zeta(s))$, whose abscissa of convergence is $\rho_0=1.7286\ldots<1.97$, the
root of $\zeta(\rho_0)=2$; numerically $A_0(0.03)=2.0673$ and $A_1(0.03)=9.409$. For $\NN$ the same values hold up to a
factor $1+O(\eta_1)$, by \eqref{eq:HBS}.

\subsubsection{The KR bound for nominal kernels}\label{sec:kr-bound}

We first prove the bound for nominal kernels whose window widths are also frozen, and then transfer it to the nominal
kernels of \S\ref{sec:kernels-steps}, whose widths $H(v/w)$ vary with the position $v$. The \emph{frozen-width} nominal
kernels $C_n^{x,\mathrm{fr}}$ are defined by \eqref{eq:kernel-rec} with the steps
$Q^{\mathrm{fr}}_{d\to n}(v,\cdot):=\kappa_{d\to n}\,\bx_{H_n}(\cdot-v/w)$; their masses are again the $c_n$ of
\eqref{eq:nominal-mass}.

\begin{theorem}[KR bound, nominal filters]\label{thm:kr}
% v3: Theorem 3.6 (stated, as in v3's proof, for widths frozen at H_n; random widths: Theorem 3.6′ below)
For every start $x$ and every $n\ge2$,
\begin{align*}
  \TV(C_n^{x,\mathrm{fr}})&\le|c_n|+\tfrac12\sum_{e\mid n,\,1<e<n}|c_e|\,\frac{H_e}{H_n}\,\Hord(n/e)\,\frac en,\\
  \KRn{C_n^{x,\mathrm{fr}}-c_n\delta_{x/n}}&\le\tfrac14\sum_{e\mid n,\,e>1}|c_e|\,H_e\,\Hord(n/e)\,\frac en .
\end{align*}
\end{theorem}

\begin{proof}
Write $C_n=C_n^{x,\mathrm{fr}}$ and $D_n:=C_n-c_n\delta_{x/n}$, a measure of mass $0$, with $D_1=0$. Since the dilation
$D_{n/d}$ maps $\delta_{x/d}$ to $\delta_{x/n}$, \eqref{eq:kernel-rec} gives
\begin{equation}\label{eq:Dn-rec}
  D_n=c_n\bigl(\bx_{H_n}-\delta_0\bigr)(\cdot-x/n)-\sum_{d\mid n,\,d<n}\kappa_{d\to n}\,\bx_{H_n}*D_{n/d}D_d ,
\end{equation}
where $c_n=-\sum_{d<n}\kappa_{d\to n}c_d$ was used.

\emph{The KR bound.} By (K4), (K1), (K2) and $|\kappa|\le1$,
\[
  \KRn{D_n}\le\tfrac14|c_n|H_n+\sum_{d\mid n,\,d<n}\frac dn\KRn{D_d}.
\]
We show by induction on $n$ that $\KRn{D_n}\le\frac14\sum_{e\mid n,e>1}|c_e|H_e\Hord(n/e)\,e/n$. Inserting this for $d<n$
(only $d>1$ contribute, since $D_1=0$) and writing $d=eg$,
\[
  \sum_{d\mid n,\,d<n}\frac dn\sum_{e\mid d,\,e>1}|c_e|H_e\Hord(d/e)\frac ed=\sum_{e\mid n,\,1<e<n}|c_e|H_e\frac en\sum_{g\mid n/e,\,g<n/e}\Hord(g)
  =\sum_{e\mid n,\,1<e<n}|c_e|H_e\Hord(n/e)\frac en,
\]
by the recursion $\Hord(m)=\sum_{g\mid m,\,g<m}\Hord(g)$. Adding the term $\frac14|c_n|H_n$, which is the term $e=n$
($\Hord(1)=1$), proves the claim.

\emph{The TV bound.} From \eqref{eq:Dn-rec} and $C_n=c_n\bx_{H_n}(\cdot-x/n)+(D_n-c_n(\bx_{H_n}-\delta_0)(\cdot-x/n))$, by (K3)
and (K2),
\[
  \TV(C_n)\le|c_n|+\frac2{H_n}\sum_{d\mid n,\,d<n}\frac dn\KRn{D_d}\le|c_n|+\frac12\sum_{e\mid n,\,1<e<n}|c_e|\frac{H_e}{H_n}\Hord(n/e)\frac en,
\]
by the KR bound for $d<n$ and the same divisor-sum identity.
\end{proof}

The bound is attained at $n=4$ (void filter, constant $H$), where $\TV(C_4)=\frac14$ while $\mu(4)=0$. The point of the proof is
that it never iterates the triangle inequality on total variation. It iterates on the KR distance, which dilation contracts
by $d/n$ at every step, and the filters enter only through the exact masses $c_e$.

\begin{theorem}[KR bound, nominal filters, varying widths]\label{thm:kr-random}
% v3: Theorem 3.6 for v3's nominal step (random widths H(v/w)); factor λ₀ per CORRECTIONS-vs-v3 entry 3
Let $\lambda_0=\exp(2.1H(X_0)/X_0)$ be as in \eqref{eq:width-product}. For every start $x$ and every $n\ge2$, the nominal kernels
$C_n^{x,\nom}$ of \S\ref{sec:kernels-steps} satisfy
\begin{align*}
  \TV(C_n^{x,\nom})&\le|c_n|+\tfrac12\lambda_0\sum_{e\mid n,\,1<e<n}|c_e|\,\frac{H_e}{H_n}\,\Hord(n/e)\,\frac en,\\
  \KRn{C_n^{x,\nom}-c_n\delta_{x/n}}&\le\tfrac14\lambda_0\sum_{e\mid n,\,e>1}|c_e|\,H_e\,\Hord(n/e)\,\frac en .
\end{align*}
\end{theorem}

\begin{proof}
For a cofactor $w$ let $\mathcal Q_w\nu:=\int\bx_{H(v/w)}(\cdot-v/w)\,d\nu(v)$ be the nominal step without its scalar filter.
We first prove two properties of $\mathcal Q_w$, for a mass-$0$ measure $\nu$ supported in an interval $I$ with
$\inf I/w\ge u/2$ and $H_*:=\inf_{v\in I}H(v/w)$:
\begin{enumerate}[label=(Q\arabic*)]
\item $\KRn{\mathcal Q_w\nu}\le(1+2\vartheta_u)\KRn{\nu}/w$;
\item $\TV(\mathcal Q_w\nu)\le(2/H_*)(1+2\vartheta_u)\KRn{\nu}/w$.
\end{enumerate}
\emph{Proof of (Q1).} For $\operatorname{Lip}f\le1$ put $g(v):=\int f\,d\bx_{H(v/w)}(\cdot-v/w)=\int_{-1/2}^{1/2}f(v/w+sH(v/w))\,ds$.
Then $|g'(v)|\le\int_{-1/2}^{1/2}(1/w)(1+|s|H'(v/w))\,ds\le(1+\vartheta_u)/w$, so
$\int f\,d\mathcal Q_w\nu=\int g\,d\nu\le(1+\vartheta_u)\KRn{\nu}/w$.

\emph{Proof of (Q2).} Write $k_v(t):=H(v/w)^{-1}\mathbf 1[|t-v/w|\le H(v/w)/2]$ for the density of the window. Integrating
by parts in $v$ ($\nu$ has mass $0$ and $F_\nu$ is its distribution function),
$\mathcal Q_w\nu=\int k_v\,d\nu(v)=-\int F_\nu(v)\,\partial_vk_v\,dv$. For each $v$, $\partial_vk_v$ consists of two point
masses at the ends of the window, of sizes $(1\pm H'(v/w)/2)/(wH(v/w))$, and the density $-H'(v/w)/(wH(v/w)^2)$ on the window;
its total variation is $\le(2+2H'(v/w))/(wH(v/w))\le(2/H_*)(1+\vartheta_u)/w$. Hence
$\TV(\mathcal Q_w\nu)\le(2/H_*)(1+\vartheta_u)\int|F_\nu|/w$, which gives (Q2).

\emph{The recursion.} With $D_n:=C_n^{x,\nom}-c_n\delta_{x/n}$, the step from the point $x/d$ lands at $x/n$ with width
exactly $H_n$, so
\[
  D_n=c_n\bigl(\bx_{H_n}-\delta_0\bigr)(\cdot-x/n)-\sum_{d\mid n,\,d<n}\kappa_{d\to n}\,\mathcal Q_{n/d}D_d .
\]
The support of $D_d$ lies within $R_d$ of $x/d$, so its image centres satisfy $v/w\ge(x/n)/2$, and on it
$H(v/w)\ge H_n(1-\vartheta_{x/n})$ by the width-variation estimate of \S\ref{sec:kernels-steps}. By (K4) and (Q1),
\[
  \KRn{D_n}\le\tfrac14|c_n|H_n+(1+2\vartheta_{x/n})\sum_{d\mid n,\,d<n}\frac dn\KRn{D_d},
\]
and by (Q2), $\TV(C_n^{x,\nom})\le|c_n|+(2/H_n)(1+3\vartheta_{x/n})\sum_{d<n}(d/n)\KRn{D_d}$. Let $\Theta_n$ be the largest
product $\prod_i(1+3\vartheta_{u_i})$ over the nominal landing scales $u_i$ of chains ending at $n$. The induction of
Theorem~\ref{thm:kr}, with every term multiplied by these factors, gives
$\KRn{D_n}\le\frac14\Theta_n\sum_{e\mid n,e>1}|c_e|H_e\Hord(n/e)\,e/n$ and the TV bound with $\frac12\Theta_n$. Finally
$\Theta_n\le\exp(3\sum_j\vartheta_{2^j\Xm})\le\exp(10^{-4}H(X_0)/X_0)\le\lambda_0$, as in \eqref{eq:width-product}.
\end{proof}

\begin{remark}[Other step kernels]\label{rem:kr-variants}
% v3: §3.4, "Variants" bullet
The proof of Theorem~\ref{thm:kr-random} uses only (Q1), the bound $\KRn{Q(x/d,\cdot)-\delta_{x/n}}\le\sigma H_n$ (for windows
$\sigma=\frac14$), and (Q2) with a constant $\tau_0'/H_*$ in place of $2/H_*$ (for windows $\tau_0'=2$); the constant in front
of the sum is then $\sigma\tau_0'\lambda_0$. The step order ``box, then filter'' and the exact compensator kernel
$\bx(H(x)/m)*\mathrm{jitter}(H(x/m))$ (for which $\sigma\le\frac12$, the sum of the two widths over $4$) satisfy these with
$\sigma\tau_0'\le1$, so they obey Theorem~\ref{thm:kr-random} with $\frac12\lambda_0$ replaced by $\lambda_0$.
\end{remark}

\subsubsection{Dyadic weights of the kernels}\label{sec:kr-weights}

\begin{corollary}\label{cor:kr-weights}
% v3: Corollary 3.7
For a start $x$ put $W^{\nom}_x(u):=\sum_{x/n\in[u,2u)}\TV(C_n^{x,\nom})/n$ and $W^{\mathrm{sh}}_x(v):=\sum_{x/e\in[v,2v)}|c_e|/e$.
\begin{enumerate}[label=(\alph*)]
\item $W^{\nom}_x(u)\le W^{\mathrm{sh}}_x(u)+\frac12\lambda_0\sum_{m\ge2}\Hord(m)m^{\eps-2}W^{\mathrm{sh}}_x(um)$.
\item Hence, by Corollary~\ref{cor:saias}, for every start $y\in[X_0,x]$ and every $u\ge X_0/2$,
  \[
    W^{\nom}_y(u)\le C_S\bigl[(1+\tfrac12\lambda_0A_0)\log(2u)+\tfrac12\lambda_0A_1\bigr]/\log y\le C_N\log(2u)/\log y,\qquad
    C_N:=C_S\bigl(2.04+4.71/\log X_0\bigr)\le158 .
  \]
\item $\sum_n\TV(C_n^{x,\nom})/n\le2.04(\log x+B_1)$; this needs no Saias bound.
\end{enumerate}
\end{corollary}

\begin{proof}
Multiply Theorem~\ref{thm:kr-random} by $\mathbf 1[x/n\in[u,2u)]/n$, sum over $n$, and substitute $n=em$ ($m\ge2$). By
\eqref{eq:RH}, $H_e/H_{em}\le m^\eps$, and $x/n\in[u,2u)$ if and only if $x/e\in[um,2um)$. This gives (a). For (b), insert
$W^{\mathrm{sh}}_y(v)\le C_S\log(2v)/\log y$ and use $\log(2um)=\log(2u)+\log m$; then
$1+\frac12\lambda_0A_0(0.03)\le1+1.03365(1+3\delta)\le2.034\le2.04$ and $\frac12\lambda_0A_1(0.03)\le4.7045(1+3\delta)\le4.71$,
and $\log(2u)\ge\log X_0$. So $C_N$ is the same constant as with frozen widths. For (c), sum (a) over all dyadic $u$ and use
$\sum_e|c_e|/e\le\log x+B_1$ (Lemma~\ref{lem:lemmaM}). Theorem~\ref{thm:kr-random} and the chain-sum identity hold verbatim
from every start $y\le x$, with $c_m(y)=\mu(m)\chi^t_y(m)$.
\end{proof}

\subsubsection{The true kernels}\label{sec:true-kernels}

The true kernels $C_n=C_n^x$ use the filter at the actual centre. They differ from the nominal ones only through the
near-threshold steps of Definition~\ref{def:near}. Write $P_{d\to n}$ and $P^{\nom}_{d\to n}$ for the step operators
$\nu\mapsto-\int Q_{n/d}(v,\cdot)\,d\nu(v)$ and its nominal version, so that $C_n=\sum_{d<n}P_{d\to n}C_d$.

\begin{proposition}[True kernels]\label{prop:true-kernels}
% v3: Proposition 3.8
Put $W^{\mathrm{op}}_x(u):=\sum_{x/n\in[u,2u)}\TV(C_n^x)/n$. If $\lambda_0C_NE_\eta\le\frac12$ (Lemma~\ref{lem:tau-near}), then on
$\Good(<x)$
\begin{align*}
  W^{\mathrm{op}}_x(u)&\le2C_N\log(2u)/\log x\quad(u\ge2X_0),\qquad
  W^{\mathrm{op}}_x(u)\le5.1\,C_N\log(2u)/\log x\quad(u\in\{X_0/2,X_0\}),\\
  \sum_n\TV(C_n^x)/n&\le\frac{2.04(\log x+B_1)}{1-2.04\lambda_0\theta_*}+3.2\le2.1(\log x+B_1)+3.2 .
\end{align*}
\end{proposition}

\begin{proof}
\emph{Duhamel form.} Put $\tilde E_n:=C_n-C_n^{\nom}$. Subtracting the two recursions,
\[
  \tilde E_n=S_n+\sum_{d<n}P^{\nom}_{d\to n}\tilde E_d,\qquad S_n:=\sum_{d<n}(P_{d\to n}-P^{\nom}_{d\to n})C_d .
\]
By Lemma~\ref{lem:near-structure}(i), only $\tau$-near and $X_0$-near steps contribute to $S_n$. On a step $d\to n$ the true
and nominal filters at a centre $c$ are $\varphi_\tau(c)\varphi_{X_0}(c)$ and $\kappa_\tau\kappa_{X_0}$; write their difference as
$(\varphi_\tau-\kappa_\tau)\varphi_{X_0}+\kappa_\tau(\varphi_{X_0}-\kappa_{X_0})$, and let $S^\tau_n$ and $S^{X_0}_n$ collect the
first and second parts. The first vanishes off $\tau$-near steps and the second off $X_0$-near steps. Each part multiplies the
measure $D_wC_d$ of centres by a function with values in $[-1,1]$ and then places a window (a probability), so by (K5)
\[
  \TV(S_n^\tau)\le\sum_{d\to n\ \tau\text{-near}}\TV(C_d),\qquad \TV(S_n^{X_0})\le\sum_{d\to n\ X_0\text{-near}}\TV(C_d).
\]
By induction, $\tilde E_n=\sum_{n'\mid n}\mathcal T^{\nom}_{n'\to n}S_{n'}$, where $\mathcal T^{\nom}_{n'\to n}$ is the signed sum
over nominal chains from level $n'$ to level $n$ ($\mathcal T_{n\to n}=\mathrm{id}$).

\emph{(i) $X_0$-near sources are terminal.} By Lemma~\ref{lem:near-structure}(ii)--(iii), $S^{X_0}_{n'}$ sits at a terminal level,
which has no children; so $\mathcal T^{\nom}_{n'\to n}S^{X_0}_{n'}=0$ for $n\ne n'$, and the parent $d$ of an $X_0$-near step is
non-terminal.

\emph{(ii) $\tau$-sources: the fixed point.} Fix $n'$ and a point $v$ in the support of $S^\tau_{n'}$, so
$|v-x/n'|\le R_{n'}$. The measure $\mathcal T^{\nom}_{n'\to n}\delta_v$ is built by nominal steps from the start $v$ with the
scalar filters $\kappa$ of the original chains. These are frozen filters, so the proof of Theorem~\ref{thm:kr-random}
applies verbatim to it, with masses given by the sharp recursion of those filters, i.e. by $c_{n/n'}(x/n')$, and with widths
$H(v/g)=H((x/n')/g)(1+O(\vartheta))$ for the cofactors $g$ of the chain. Hence $\TV(\mathcal T^{\nom}_{n'\to n}\delta_v)$ is
at most $\lambda_0$ times the right side of Theorem~\ref{thm:kr-random} for the start $x/n'$ and level $n/n'$, and
Corollary~\ref{cor:kr-weights}(b) at the start $x/n'$ gives
\[
  \sum_{x/n\in[u,2u)}\frac{\TV(\mathcal T^{\nom}_{n'\to n}S^\tau_{n'})}{n}\le\lambda_0\frac{\TV(S^\tau_{n'})}{n'}\,
  \frac{C_N\log(2u)}{\log(x/n')}.
\]
Summing over $n'$ and grouping by the parent $d$ of each $\tau$-near step $d\to n'$ ($n'=dw$),
\[
  \sum_{x/n\in[u,2u)}\frac{\TV(\tilde E_n^\tau)}n\le\lambda_0C_N\log(2u)\sum_d\frac{\TV(C_d)}d\,\tilde\eta(x/d),\qquad
  \tilde\eta(v)=\sum_{w\ \tau\text{-near at }v}\frac1{w\log(v/w)},
\]
with $d$ non-terminal. Let $M$ be the least constant with $W^{\mathrm{op}}_x(v)\le M\log(2v)/\log x$ for all $v$, counting
non-terminal levels and the $\tau$-parts at terminal levels; $M<\infty$ since $n\le x/\Xm$. Then
$\sum_d\TV(C_d)/d\cdot\tilde\eta(x/d)\le\sum_v W^{\mathrm{op}}_x(v)\sup_{[v,2v)}\tilde\eta\le ME_\eta/\log x$, and
$W^{\mathrm{op}}_x(u)\le W^{\nom}_x(u)+\lambda_0C_NME_\eta\log(2u)/\log x$. Hence $M\le C_N(1+\lambda_0ME_\eta)$, and
$M\le2C_N$ since $\lambda_0C_NE_\eta\le\frac12$.

\emph{(iii) Terminal levels.} They gain
\[
  \sum_{n\ \text{terminal}}\frac{\TV(S^{X_0}_n)}n\le\sum_d\frac{\TV(C_d)}d\,\eta_{X_0}(x/d),\qquad
  \eta_{X_0}(v):=\sum_{w\in\NN:\ |v-w\Xm|\le1.02H(v)}\frac1w .
\]
\begin{itemize}
\item For $v\le X_0^2/2$ the $w$-window has length $2.04H(v)/\Xm<1$ and lies below $\Xm$, where $\NN=\N$; so it holds at most
  one $w$, and $\eta_{X_0}(v)\le1.01X_0/v$.
\item For $v>X_0^2/2$, the window $I_v$ has length $\le2.1H(v)/X_0$ at height $v/X_0$. By the lag bound
  Corollary~\ref{cor:lag}, $|N-T|\le H$ at that height, and under the target \eqref{eq:target}, $T(I_v)\le2|I_v|+2+\#\Abs(I_v)$;
  by Lemma~\ref{lem:lemmaP}, $\#\Abs(I_v)\le2(v/X_0)^{\theta/2}$. So $I_v$ holds at most
  $2.1H(v)/X_0+2H(v/X_0)+4+2(v/X_0)^{\theta/2}$ elements, and
  $\eta_{X_0}(v)\le1.01\bigl(2.1H(v)+2X_0H(v/X_0)+4X_0+2X_0(v/X_0)^{\theta/2}\bigr)/v$.
\item Since $\eta_{X_0}(v)=0$ for $v<2\Xm-1.1H(X_0)$, these give
  $\sum_{v\ \text{dyadic}}\log(2v)\sup_{[v,2v)}\eta_{X_0}\le1.52\log(2X_0)$ and $\sum_v\sup_{[v,2v)}\eta_{X_0}\le1.52$: the first
  bullet contributes $0.506\log(2X_0)$ from the block $[X_0,2X_0)$ and $1.01\sum_{j\ge1}2^{-j}(\log(2X_0)+j\log2)$ from the
  others, and the second bullet less than $e^{-10^7}$.
\end{itemize}
Terminal levels lie in the blocks $u\in\{X_0/2,X_0\}$. With (ii) they add at most $2C_N\cdot1.52\log(2X_0)/\log x$ there.
Since $1.52\log(2X_0)\le1.53\log(2u)$ for these $u$, the total there is $\le5.1\,C_N\log(2u)/\log x$.

\emph{The $\infty$-norm.} Run (ii) with Corollary~\ref{cor:kr-weights}(c) in place of (b): the transfer from $n'$ adds
$\le2.04\lambda_0(\log(x/n')+B_1)\TV(S^\tau_{n'})/n'$, and grouping by parents gives the weight
$\theta(v)=\sum_{w\ \tau\text{-near at }v}(\log(v/w)+B_1)/w$; this yields the first term. For (iii) use the block bound
$\sup_vW^{\mathrm{op}}_x(v)\le2.1$, which needs no Saias bound: $W^{\mathrm{sh}}_x(v)\le\sum_{e\in(x/2v,x/v]}1/e\le1$, so
$W^{\nom}_x(v)\le1+\frac12\lambda_0A_0\le2.04$ by Corollary~\ref{cor:kr-weights}(a), and the $\tau$-part adds
$\le2.04\sum_v\sup\eta\le10^{-6}$. That gives $\le2.1\cdot1.52\le3.2$.

\emph{No circularity.} The inputs (Lemma~\ref{lem:lemmaM}, Corollary~\ref{cor:saias}, Lemma~\ref{lem:tau-near},
Corollary~\ref{cor:lag}) concern $\NN$ on $[1,x/\Xm]$ only.
\end{proof}

\begin{remark}
% v3: remark after Proposition 3.8
Treating the $X_0$-near steps like the $\tau$-near ones in the fixed point would put
$\sum_v\log(2v)\sup\eta_{X_0}\approx1.5\log X_0$ into $E_\eta$, and $C_NE_\eta\le\frac12$ would fail. They are harmless only
because they are terminal.
\end{remark}

\begin{lemma}[$\tau$-near energy]\label{lem:tau-near}
% v3: Lemma 3.9
For $v\ge\Xm$ put $\eta(v):=\sum_{w\ \tau\text{-near at }v}1/w$, $\tilde\eta(v):=\sum_{w\ \tau\text{-near at }v}1/(w\log(v/w))$
and $\theta(v):=\sum_{w\ \tau\text{-near at }v}(\log(v/w)+B_1)/w$, where $w$ is $\tau$-near at $v$ if
$|\tau_w-v/w|\le2H(v/w)$ and $v/w\ge\Xm$. Put
\[
  E_\eta:=\sum_{v\ \text{dyadic}}\log(2v)\sup_{[v,2v)}\tilde\eta,\quad
  E_\eta^0:=\sum_{v\ \text{dyadic}}\log(2v)\sup_{[v,2v)}\eta,\quad
  E_\theta:=\sum_{v\ \text{dyadic}}\log(2v)\sup_{[v,2v)}\theta,\quad \theta_*:=\sup_v\theta(v).
\]
Then, on $\Good(<y)$ at all heights used,
\[
  E_\eta,\ E_\eta^0,\ E_\theta,\ \theta_*\le10^4H(X_0)\frac{\log^2X_0}{X_0}=10^4\Lambda^2e^{8(\log\Lambda)^2-\Lambda}\le e^{-0.99\Lambda},\qquad
  E_\eta^{(2)}:=\sum_{v\ \text{dyadic}}(\log3v)^2\sup_{[v,2v)}\eta\le10^4H(X_0)\frac{\log^3X_0}{X_0},
\]
and $\sup_y\log y\cdot\eta(y)\le e^{-0.99\Lambda}$. (At $\Lambda=10^8$ these are $e^{-(10^8-2762)}$ and $e^{-(10^8-2770)}$.)
\end{lemma}

\begin{proof}
\emph{Localisation.} Let $w$ be $\tau$-near at $v'$, and write $w=pm$ with $p$ the prime of $w$ that is largest in decision
order, so that $\tau_w=t(p)$ and $\Pplus(m)\le p+H(p)$. Since $|t(p)-p|\le H(p)$, the condition $|t(p)-v'/(pm)|\le2H(v'/w)$
gives $|p-v'/(pm)|\le3H(p)(1+o(1))$, i.e. $p^2\in v'/m\pm3.1\,pH(p)$. So $p$ lies in a window of length $3.1H(p_0)$
around $p_0:=\sqrt{v'/m}\ge\Xm$.

\emph{Counting.} Let $[P,2P)$ be the dyadic block of $p_0$. The window lies in $[P,4P]$ and has length
$\le3.2H(P)\le6H(P)$, so by Corollary~\ref{cor:localcount} it holds at most $13.4\cdot3.2H(P)+K_{\mathrm{loc}}(P)+2\le
43H(P)+K_{\mathrm{loc}}(P)+2$ primes, each with $1/p\le1.01/P$. For $v'\in[v,2v)$ this forces $m\in(v/4P^2,2v/P^2]$, and a
pair $(P,m)$ occurs for at most three dyadic $v$. The $m$ are $4P$-smooth, so by Lemma~\ref{lem:cm},
$\sum_m\log(8mP^2)/m\le5.4\log^2P(1+5/\log P)$ and $\sum_m1/m\le1.9\log P$.

\emph{$E_\eta$.} Since $\log(2v)/\log(v/w)\le1.01\log(8mP^2)/\log P$, each $P$-block contributes at most
$16.5(43H(P)+K_{\mathrm{loc}}(P)+2)\log P/P$.

\emph{$E_\eta^0$, $E_\eta^{(2)}$, $E_\theta$, $\theta_*$.} $E_\eta^0$ has no division by $\log(v/w)$; $E_\eta^{(2)}$ has one
more factor $\le1.01\log(8mP^2)$; $E_\theta$ has the extra factor $\log(v/w)+B_1\le1.01\log4P$. So each $P$-block contributes
$\le16.5(43H(P)+K_{\mathrm{loc}}(P)+2)\cdot1.02\log^2P/P$ to $E_\eta^0$ and $E_\theta$ (one more $\log P$ to $E_\eta^{(2)}$),
and $\theta_*$ is bounded by the same expression without the weight $\log(2v)$.

\emph{Summation.} The blocks decay geometrically (ratio $\le2^{\eps-1}\le0.52$, by \eqref{eq:RH}), and
$K_{\mathrm{loc}}(P)\le H(P)$. Summing from $P=X_0/2$ gives $\le2.1\cdot16.5\cdot44\cdot1.02\cdot H(X_0)\cdot4\log^2X_0/X_0
\le10^4H(X_0)\log^2X_0/X_0$.

\emph{$\log y\cdot\eta(y)$.} The $m$ with $p_0\in[P,2P)$ form one block $(y/4P^2,y/P^2]$ of $4P$-smooth elements, whose
$1/m$-weight is $\le2C_0\rho(\log(y/P^2)/\log4P)e^{o(1)}$ by Lemma~\ref{lem:smooth-v}. This absorbs the factor $\log y$
when $\log y\ge3\log P$, and $\log y\le3\log P$ otherwise.

\emph{Numerics.} $H(X_0)=e^{8(\log\Lambda)^2}$, which is $\le e^{2715}$ at $\Lambda=10^8$, and $K_{\mathrm{loc}}(X_0)\le H(X_0)$. For
every $\Lambda\ge10^8$, $9.3+2\log\Lambda+8(\log\Lambda)^2\le0.01\Lambda$, which gives the bounds $e^{-0.99\Lambda}$.
\end{proof}

Consequently $C_NE_\eta\le\frac12$ and $2.04\theta_*\le0.02$, with room $e^{-10^8}$. The per-scale constant $K_{\mathrm{loc}}(P)$
is essential: a constant growing with the current scale would make $E_\eta$ grow with $x$. The quantity $E_\eta^0$ is what
makes the $\tau$-near term $\NT$ summable with the weights $\log(2u)$ of Lemma~\ref{lem:predictable}; a pointwise bound on
$\log y\cdot\eta(y)$ alone would give a sum growing like $\log\log x$.

\begin{lemma}[Filtered cofactor weight]\label{lem:S}
% v3: Lemma 3.10 (Lemma S)
For a start $y$ and a scale $v\ge X_0/2$ put $\omega_y(v):=\sum_{w\in\NN:\ y/w\in[v,2v),\ \tau_w<y/w}1/w$. Then:
\begin{enumerate}[label=(\roman*)]
\item $\omega_y(v)\le1.2$;
\item for $y\ge4v^2$ the filter forces $\Pplus(w)\le3v$. Put $u_y:=\log(y/2v)/\log3v$. If $u_y\le\log3v$, then
  $\omega_y(v)\le2C_0\rho(u_y)e^{a_V}$ with $a_V\le13u_y\log(u_y+2)/\log3v$; so $\omega_y(v)\le2C_0\rho(u_y)e^{0.02}$ for
  $u_y\le30$ and $\omega_y(v)\le2C_0e^{-0.9u_y\log u_y}$ for $30<u_y\le\log3v$. If $u_y>\log3v$, then
  $\omega_y(v)\le5(\log3v)^{1.65}e^{-u_y/2}$;
\item $\sum_{Y=2^k\ge v}\log(2Y)\sup_{y\in[Y,2Y)}\omega_y(v)\le C_S'(\log3v)^2$ with $C_S':=10$.
\end{enumerate}
\end{lemma}

\begin{proof}
(i) $\omega_y(v)\le\sum_{w\in(y/2v,y/v]}1/w\le\log2+2\eta_1+2v/y\le1.2$ by \eqref{eq:HBS} (and $\le1$ below $\Xm$).

(ii) If $w$ passes the filter then $\tau_w<y/w\le2v$, so every prime of $w$ is decided before slot $2v$ and is $\le2v+H(2v)\le3v$.
Since $w\in(y/2v,y/v]$, Lemma~\ref{lem:smooth-v} with $z=3v$ and $\ell=\log(y/2v)$ gives the first bound (its hypothesis
$u\le L_z$ is $u_y\le\log3v$); the numerical forms follow from $a_V\le13u\log(u+2)/\log3v$ and $\log3v\ge10^8$, and from
$\rho(u)\le e^{-u\log u}$. For $u_y>\log3v$, Rankin's bound with exponent $1-1/(2\log3v)$ and Lemma~\ref{lem:cm} gives
$\omega_y(v)\le5(\log3v)^{1.65}e^{-u_y/2}$.

(iii) The range $Y<4v^2$ contributes at most $1.2\cdot\frac3{2\log2}(1+o(1))(\log3v)^2\le2.7(\log3v)^2$. The range
$Y\ge4v^2$ contributes at most $\frac{2\cdot0.70\cdot e^{0.02}}{\log2}(\log3v)^2\int_1^\infty\frac{1+u}{\Gamma(u+1)}\,du\le
7.11(\log3v)^2$, using $\rho(u)\le1/\Gamma(u+1)$ (from $\rho(u)\le\rho(u-1)/u$) and
$\int_1^\infty(1+u)/\Gamma(u+1)\,du=3.4479$. The parts with $u_y>30$ add less than $10^{-20}(\log3v)^2$. The total is
$\le9.82(\log3v)^2$.
\end{proof}

The filter forces a dense cofactor ($v<\sqrt y$) to be $3v$-smooth, hence Dickman-small. No Riemann-sum or
short-interval smooth-number input is needed, and Lemma~\ref{lem:smooth-v} avoids the Rankin loss $(\log3v)^{2+e}$.

% 04d-noise.tex — v3 §3.5: the noise, per scale.
\subsection{The noise, per scale}\label{sec:noise-perscale}
% v3: §3.5

\paragraph{Inputs from \S\ref{sec:queue}.} Each is used frozen at the scale where it is read.
\begin{itemize}
\item The queue bound $Q\le B(u)$ (Theorem~\ref{thm:queue}, \eqref{eq:queueB}); so $|\Delta Q(W)|\le B(u)$ on windows at
  scale $u$.
\item The per-prime increment bound (Proposition~\ref{prop:ecl}, Corollary~\ref{cor:increment}). On $E_{\mathrm{cl}}^{(u)}$,
  for every slot window $I\subset[u,2u]$ with $|I|\le H(u)$, every non-absorbed latest prime $q\ge y(u)$, every $a\ge1$ and
  every jitter offset, $X_q^{(a)}(I)\le2M_1(u)$. The event $E_{\mathrm{cl}}^{(u)}$ fails on $\Good$ with probability
  $\le L_u^{-2}$.
\item The local count (Corollary~\ref{cor:localcount}). For dyadic $P$ and windows $J\subset[P,4P]$ of length $\le6H(P)$
  (part (i)), the non-absorbed count, and so the prime count, is $\le13.4|J|+K_{\mathrm{loc}}(P)+2$, with
  $K_{\mathrm{loc}}(P)=2B(4P)+7\lambda_P/3+23$ and $\lambda_P=4\log P$. For real $P\ge X_0$ (part (ii)) the same holds with
  $K_{\mathrm{loc}}(P)+K_{\mathrm{loc}}(2P)+4$ in place of $K_{\mathrm{loc}}(P)+2$.
\item Absorption (Lemmas~\ref{lem:lemmaP}, \ref{lem:rhoA}, Proposition~\ref{prop:S1S2}). Non-absorbed latest primes satisfy
  $q\ge y(u)=H(u)L_0L_u^6$, so $|I/q|\le L_u^{-6}$ for $|I|\le H(u)$. The small-$q$ case S1 does not arise in the present
  regime.
\end{itemize}

\paragraph{The composites with $q^2\mid n$: the source $s_2$.} Recall $C^{(\ge2)}=\Comp^{(2)}+\Mart^{(2)}$
(\S\ref{sec:noise-identity}). Fix $u$ and a window $I\subset[u,2u]$, and put $\sigma:=\frac12\sqrt{L_u\log L_u}$.
% Logged correction 52 (CORRECTIONS-vs-v3.md): v3 bounds the Dickman factor by 1/Gamma(max(1, L_u/s - 2)), which exceeds 1 on
% (1,2) (up to 1.129), and starts its sup at log X_0, although latest primes can lie in [X_0^-, X_0); it also applies the
% smooth-element weight with a smaller Dickman variable and without its factor e^{a_V}. The bound is re-derived below, with
% the endpoint X_0^-, and s_2 is defined by its closed form.
\begin{itemize}
\item The latest-prime filter forces $\tau_m<t(q)$ for $q^am$, so $\Pplus(m)\le2q$.
\item \emph{Interchange the sums.} Fix $a\ge2$ and $m$. Then $q^am\in I$ means $q\in I_{m,a}:=(I/m)^{1/a}$, and the $q$-sum of
  the landing probabilities is $\int_{I_{m,a}}\rhot\le1.01\,|I_{m,a}\cap[\Xm,\infty)|$, since $\rhot\le1.01$ surely (see the
  variance bound below) and $\rhot=0$ below $\Xm$. With $y=q^am\in I$ we have $dq=dy/(amq^{a-1})$, and on
  $I_{m,a}\cap[\Xm,\infty)$, $q\ge Q_{m,a}:=\max(\Xm,(u/m)^{1/a})$. Hence
  \[
    \frac{\Comp^{(2)}(I)}{|I|}\le1.01\sum_{a\ge2}\frac1a\sum_m\frac1{m\,Q_{m,a}^{a-1}},
  \]
  where $m$ runs over $\NN$ with $m\le2u/(\Xm)^a$ and $\Pplus(m)\le2(2u/m)^{1/a}$.
\item \emph{Blocks.} Split the $m$-sum into blocks $m\in(Y,2Y]$, $Y=2^{k-1}$ ($k\ge0$); for each $a$, at most $1.443L_u$ of
  them are nonempty. Define $s$ by $e^{as}=u/(2Y)$ and put $s_Y:=\max(\log\Xm,s)$. In the block, $Q_{m,a}\ge e^{s_Y}$,
  $\log Y=L_u-as-\log2$, and the $m$ are $z$-smooth with $\log z:=s+\log2+(2\log2)/a\le s+1.39$. A nonempty block has
  $Y<2u/(\Xm)^a$, so $s>\log\Xm-0.7$ and $\log z\ge\log\Xm-0.01$. The block contributes at most
  $\frac1ae^{-(a-1)s_Y}v(\log Y,z)$, where $v(\ell,z)$ is the $1/n$-weight of the $z$-smooth $n\in\NN\cap(e^\ell,2e^\ell]$, as in
  Lemma~\ref{lem:smooth-v}.
\item \emph{First bound.} Always $v\le1$ (the small heights in the proof of Lemma~\ref{lem:smooth-v}; the block $Y=\frac12$
  holds only $m=1$), and $s_Y\ge\log\Xm$. So
  $\Comp^{(2)}(I)/|I|\le1.01\cdot1.443L_u(\Xm)^{-1}\bigl(\frac12+\sum_{a\ge3}(\Xm)^{2-a}/a\bigr)\le0.73L_u/\Xm$.
\item \emph{Second bound, for $\sigma>\log\Xm$.} Then $L_u\ge10^{15}$, since $\log\Xm\ge10^8-1$. We show that every block term
  is at most $\frac1ae^{-\sigma}(\Xm)^{-(a-2)/2}$, that is, $(a-1)s_Y-\log v\ge\sigma+\frac12(a-2)\log\Xm$.
  \begin{itemize}
  \item If $s_Y\ge\sigma$, or if $(a-2)s_Y\ge2\sigma$, this holds with $v\le1$.
  \item Otherwise $s\le s_Y<\sigma$ and $as<4\sigma$, so $\ell:=\log Y\ge L_u-4\sigma-0.7\ge(1-4\cdot10^{-7})L_u$ and
    $\log z<\sigma+1.39$. Since $(a-1)s_Y\ge\frac12(a-2)\log\Xm$, it suffices that $-\log v\ge\sigma$. Put $u':=\ell/\log z$
    and $\omega:=2\sqrt{L_u/\log L_u}$, so that $\sigma=\frac14\omega\log L_u$ and
    $u'\ge(L_u-4\sigma-0.7)/(\sigma+1.39)\ge(1-4\cdot10^{-7})\,\omega$. For every $w\ge1$, iterating $\rho(w)\le\rho(w-1)/w$
    (from \ref{R1}) down to $w-\lfloor w\rfloor\in[0,1)$, where $\rho=1$, gives
    $\rho(w)\le\Gamma(w-\lfloor w\rfloor+1)/\Gamma(w+1)\le1/\Gamma(w+1)\le e^{-w\log w+w}$ (Stirling).
  \item If $u'\le\log z$: Lemma~\ref{lem:smooth-v} gives $v\le C_0\rho(u')e^{a_V}$ with
    $a_V=12.8\,u'\log(u'+2)/\log z\le1.3\cdot10^{-7}u'\log(u'+2)$. So
    \[
      -\log v\ge u'(\log u'-1)-1.3\cdot10^{-7}u'\log(u'+2)\ge(1-10^{-6})\,\omega\bigl(\tfrac12\log L_u-\tfrac12\log\log L_u-0.31\bigr),
    \]
    which is $\ge1.75\sigma$ for $\log L_u\ge34.5$: with $x:=\log L_u$ the ratio to $\sigma$ is $(1-10^{-6})(2-(2\log x+1.24)/x)$,
    which increases in $x$, and $\sigma\ge\log X_0^-$ forces $L_u\ge10^{15}$.
  \item If $u'>\log z$: then $\sqrt\ell>\log z$, and $v(\ell,z)\le v(\ell,z')$ with $\log z':=\sqrt\ell$. Lemma~\ref{lem:smooth-v} at
    the Dickman variable $\sqrt\ell$, where $a_V=12.8\log(\sqrt\ell+2)$, gives
    $-\log v\ge\sqrt\ell(\log\sqrt\ell-1)-12.8\log(\sqrt\ell+2)\ge5.5\sigma$; its ratio to $\sigma$ also increases in $\log L_u$, and
    is $\ge5.53$ at $L_u=10^{15}$.
  \end{itemize}
  Summing, $\Comp^{(2)}(I)/|I|\le1.01\cdot1.443L_u\,e^{-\sigma}\bigl(\frac12+\sum_{a\ge3}(\Xm)^{-(a-2)/2}/a\bigr)\le0.73L_u\,e^{-\sigma}$.
\item Together, for every window $I\subset[u,2u]$,
  \begin{equation}\label{eq:s2}
    \frac{\Comp^{(2)}(I)}{|I|}\le s_2(u):=9L_u\exp\Big(-\max\big(\log\Xm,\tfrac12\sqrt{L_u\log L_u}\big)\Big).
  \end{equation}
  In particular $s_2(u)\le9L_u/\Xm\le9(1+10^{-9})L_u/X_0$ always, since $X_0/\Xm\le1+H(X_0)/X_0$.
\item $\sum_{u\text{ dyadic}}\log(2u)\,s_2(u)\le e^{-9\cdot10^7}$, and the same holds with the weight $(\log3u)^2$. For
  $L_u\le\Lambda^2$, use $s_2(u)\le9L_u/\Xm$: these at most $1.443\Lambda^2$ scales contribute $\le14\Lambda^8/\Xm$. For
  $L_u>\Lambda^2$, $\frac12\sqrt{L_u\log L_u}\ge3.03\sqrt{L_u}$, and $10L^3e^{-3.03\sqrt L}$ is decreasing, so these scales
  contribute $\le10\Lambda^6e^{-3.03\Lambda}+\frac{10}{\log2}\int_{\Lambda^2}^\infty L^3e^{-3.03\sqrt L}\,dL\le10\Lambda^7e^{-3.03\Lambda}$.
  The total is $\le e^{-0.99\Lambda}$.
\item Put $A_u:=\#\{a\ge2:(\Xm)^a\le2u\}\le1.01L_u/\Lambda$. By Corollary~\ref{cor:increment} (which covers every
  $a\ge1$), the increment of a prime $q$ to $\Mart^{(2)}(I)$ is at most $b_2:=2A_uM_1(u)$.
\end{itemize}

\paragraph{The per-scale Freedman events.} Put $\lambda'_u:=3\log u+2\log\log u+2+\log H(u)$. For dyadic $u\ge X_0$, let
$\mathcal W'_u$ be the set of windows $[a,b]\subset[u,2u]$ with integer endpoints and $b-a\le1.03H(u)+1$; there are at most
$(u+1)(1.03H(u)+2)$ of them. Let $\mathcal W_u\subset\mathcal W'_u$ consist of those with $b-a\ge H(u)-1$.
\begin{itemize}
\item $G_u^{\mathrm{NV}}$ (the event E3 of Definition~\ref{def:good}): every $W\in\mathcal W_u$ satisfies \eqref{eq:F1}
  and \eqref{eq:F2}.
\item $G_u^{\mathrm{sw}}$ (the event E3$'$): every $W\in\mathcal W'_u$ satisfies \eqref{eq:F1} and \eqref{eq:F2}.
\end{itemize}
\begin{align}
  |\Mart(W)|&\le\sqrt{2\bar V_W\lambda'_u}+\tfrac43M_1(u)\lambda'_u,&\bar V_W&:=2M_1(u)\cdot50\log(2u)\cdot|W|;\label{eq:F1}\\
  |\Mart^{(2)}(W)|&\le\sqrt{2\bar V^2_W\lambda'_u}+\tfrac23b_2\lambda'_u,&\bar V^2_W&:=b_2\cdot s_2(u)\cdot|W|.\label{eq:F2}
\end{align}

\paragraph{Truncation.} As in Theorem~\ref{thm:queue}(i), Freedman's inequality is applied to the truncated increments
$\tilde X_q(W):=\min(X_q(W),2M_1(u))$ and $\tilde X^{(2)}_q(W):=\min(X^{(2)}_q(W),b_2)$, each compensated by its own
conditional mean. These are martingales in decision order, with increments $\le2M_1$ and $\le b_2$ surely. Freedman's
inequality is used in the form
\[
  \Prob(\exists n:\ S_n\ge a,\ V_n\le v)\le\exp\Big(-\frac{a^2}{2(v+ba/3)}\Big),
\]
for increments bounded by $b$ \cite[(4.1)]{Freedman1975} (with Bennett's bound for the exponent), with the deterministic $a$
and $v=\bar V_W,\bar V^2_W$ of \eqref{eq:F1}, \eqref{eq:F2}.
\begin{itemize}
\item \emph{Truncation is inactive on $\Good(<t(W))$.} For every latest prime $q$ contributing to $W$, the set of cofactors
  (elements whose primes are all decided before $q$) is $\mathcal F_{q-}$-measurable. So ``$X_q(W)\le2M_1$ for every jitter
  offset of $q$'' is an $\mathcal F_{q-}$-measurable event. It may not yet be certified at time $s_q$ (for $q<\sqrt{2u}$).
  But on $\Good(<t(W))$ the event $E_{\mathrm{cl}}^{(u)}$ holds (its time $8u/L^{11}+H$ is $<t(W)$), and
  Corollary~\ref{cor:increment} then gives it for every $q$ and every offset. So $\tilde X_q=X_q$ and
  $\E[\tilde X_q\mid\mathcal F_{q-}]=\E[X_q\mid\mathcal F_{q-}]$. The truncated and the true martingales therefore agree on
  $\Good(<t(W))$, and likewise for $\Mart^{(2)}$.
\item \emph{The variance bound holds on $\Good(<t(W))$.} We have
  $\mathrm{Var}(\tilde X_q\mid\mathcal F_{q-})\le2M_1\E[X_q\mid\mathcal F_{q-}]$, so $V_W\le2M_1\Comp(W)$; likewise
  $V^2_W\le b_2\Comp^{(2)}(W)\le b_2s_2|W|$. Next,
  $\Comp(W)=\int_W(\Kpt\rhot-K_{\Abs}\rhot+\NT_{\mathrm{pt}})\le1.01\sum_v\omega_y(v)\cdot|W|+\int_W|\NT_{\mathrm{pt}}|\le50\log(2u)\,|W|$.
% Logged correction 73: v3 says (H_BS) "at heights <= u"; a short window W may close at t(W) < u + 12 log u + 3, where that is
% not certified. Lemma S(i) at y in W reads only cofactors w <= y/v <= 4u/X_0, far below the horizon of Good(<t(W)).
  This uses $\sum_v\omega_y(v)\le1.2\log_2(4u/X_0)$ (Lemma~\ref{lem:S}(i), that is, \textup{(H$_{\mathrm{BS}}$)} at heights
  $\le4u/X_0$, below the certified horizon of $\Good(<t(W))$, Remark~\ref{rem:horizon}), and the sure bound $\rhot\le1.01$. The latter holds because each continuous
  slot decides at most one prime, and the starts of consecutive continuous slots are at least $1-10^{-8}$ apart (between them
  $\int T'_{\mathrm{cont}}=1$, \S\ref{sec:greedy}, with $T'_{\mathrm{cont}}\le1+10^{-8}$). No upper bound on the length of a slot is used
  in this section.
% Logged correction 77: one unit of continuous mass lies between consecutive continuous-slot starts, not in each continuous slot
% (a continuous slot followed by a jump slot, or an empty one in the tie case, carries less).
\item So on $\Good(<t(W))$ the event $\{|\Mart(W)|>$ the bound of \eqref{eq:F1}$\}$ lies in Freedman's event for the
  truncated martingale; likewise for \eqref{eq:F2}. No positivity of the margin is used.
\end{itemize}

\paragraph{Probability.} With $a=\sqrt{2v\lambda'_u}+\frac23b\lambda'_u$, Freedman's inequality gives, per window and per
martingale, a two-sided failure probability $\le2e^{-\lambda'_u}=2e^{-2}u^{-3}L_u^{-2}H(u)^{-1}$. With four bounds per window,
\[
  \Prob\big((G_u^{\mathrm{NV}}\cap G_u^{\mathrm{sw}})^c\cap\Good(<t)\big)\le8e^{-2}(u+1)(1.03H(u)+2)u^{-3}L_u^{-2}H(u)^{-1}
  \le1.2\,u^{-2}L_u^{-2}.
\]
Each window's bound is an event whose time is the moment the window closes ($\le2u$). These are the events E3 and E3$'$ of
Definition~\ref{def:good}. Together with $E_{\mathrm{cl}}$ they enter $\Good(<t)$ and the budget $\Prob(\tau<\infty)\le1.5/\Lambda$
of Lemma~\ref{lem:budget}. The budget is dominated by $E_{\mathrm{cl}}$: $\sum_uL_u^{-2}\le1.45/\Lambda$, while the
Freedman part is $\le2X_0^{-2}$.

\begin{lemma}[per-scale noise]\label{lem:noise}
% v3: Lemma 3.11
For $y\in[u,2u]$, work on $\Good(<y+H(y))$. This event contains the E3/E3$'$ windows of the blocks
$u'\in\{u/2,u,2u\}$ that close by $y+H(y)/2$, and $E_{\mathrm{cl}}^{(u')}$. Put
\begin{multline*}
  \nu_0(u'):=\sqrt{\frac{206M_1(u')\lambda'_{u'}\log(2u')}{H(u')}}\Big(1+\frac1{H(u')}\Big)
  +\sqrt{\frac{4.2A_{u'}M_1(u')s_2(u')\lambda'_{u'}}{H(u')}}\\
  +\frac{\frac43(1+A_{u'})M_1(u')\lambda'_{u'}+3B(u')+100\log(2u')+6}{H(u')}+1.1s_2(u').
\end{multline*}
Then
\[
  |\nu(y)|\le\bar\nu(u):=2.05\max_{u'\in\{u/2,u,2u\}}\nu_0(u').
\]
% Logged correction 73 (R17, C73-6): the good event of the bound on E(I) is stated.
For every slot window $I\subset[u,2u]$ with $|I|\le H(u)$, put $E(I):=\sigma(I)-\int_Im$, where
$m=T'-\Kpt\rhot-\NT_{\mathrm{pt}}$ is the pointwise margin density (\S\ref{sec:noise-identity}). Then, on $\Good(<\max I+H(u))$,
\begin{multline*}
  |E(I)|\le\sqrt{200M_1(u)\lambda'_u\log(2u)(|I|+2)}+\sqrt{4A_uM_1(u)s_2(u)\lambda'_u(|I|+2)}\\
  +\tfrac43(1+A_u)M_1(u)\lambda'_u+B(u)+100\log(2u)+s_2(u)(|I|+2)+2 .
\end{multline*}
Write $\hat\nu(u):=\sup_{|I|\le H(u)}|E(I)|/H(u)$. Both bounds depend on $u$ only: nothing at the current scale $x$
enters.
\end{lemma}

\begin{proof}
\emph{The bound on $\nu$.} By \eqref{eq:margin-identity},
$H(y)\nu(y)=\Mart(W_y)+\Mart^{(2)}(W_y)-\Delta Q(W_y)+\Comp^{(2)}(W_y)+E_{\mathrm{edge}}$.
\begin{itemize}
\item $E_{\mathrm{edge}}$ counts the two partial slot cells at the ends of $W_y$. Each holds at most one slot and at most
  $A_j\le B(u)+1$ composites, so $|E_{\mathrm{edge}}|\le2B(u)+4$.
\item $|\Delta Q|\le B(u)$ by Theorem~\ref{thm:queue}, and $\Comp^{(2)}(W_y)\le s_2|W_y|$ by \eqref{eq:s2}.
\item Choose $W^-\subset W_y\subset W^+$ in $\mathcal W_u$ with $|W^+\setminus W^-|\le2$. Counts and compensators are
  monotone in the window, so $\Mart(W^-)-\Comp(W^+\setminus W^-)\le\Mart(W_y)\le\Mart(W^+)+\Comp(W^+\setminus W^-)$, with
  $\Comp(W^+\setminus W^-)\le100\log(2u)$. The same holds for $\Mart^{(2)}$, with $2s_2$ in place of $100\log(2u)$.
\item Apply \eqref{eq:F1} and \eqref{eq:F2} with $|W^\pm|\le1.03H(u)+1$, and divide by $H(y)\ge H(u)$.
\item If $W_y$ crosses $u$ or $2u$ (that is, $y$ is within $H(y)/2$ of it), split $W_y$ there. Each piece lies in one
  dyadic block $u'$ and has length $\le1.03H(u')$, so the above applies to it in block $u'$. The two bounds add, and
  $H(y)/H(u')\le1.02$ by \eqref{eq:RH}; the factor $2.05$ covers this.
\end{itemize}

\emph{The bound on $E(I)$.} The bookkeeping \eqref{eq:bookkeeping} on the slot window $I$ has no partial cells. By the
computation of \eqref{eq:margin-identity} without the $B$-average,
$\sigma(I)=\int_Im-\Comp^{(2)}(I)-\Mart(I)-\Mart^{(2)}(I)+\Delta Q(I)\pm2$. Now sandwich $I$ between windows of
$\mathcal W'_u$ as above, and apply \eqref{eq:F1}, \eqref{eq:F2} and \eqref{eq:s2}.
\end{proof}

\paragraph{The \textup{(SW)} sums.} Under the tolerance $M_1\le H/L^{14}$ of \eqref{eq:params}, with $B\le28M_1L^2$,
$A_u\le1.01L_u/\Lambda$ and $\lambda'_u\le18.2L_u$ (since $\log H\le15.1\Lambda$ on the raised range),
\[
  \bar\nu(u)\le128L_u^{-6}+2.3\max_{u'}s_2(u'),\qquad\hat\nu(u)\le62L_u^{-6}+1.1s_2(u),
\]
and
\begin{equation}\label{eq:sw-sums}
  \sum_{u\text{ dyadic}\ge X_0}\log(2u)\,\bar\nu(u/2)\le10^{-30},\qquad
  \sum_v(\log3v)^2\,\bar\nu(v/2)\le10^{-22},\qquad
  \sum_v\log(3v)\,\hat\nu(v)\le10^{-30}.
\end{equation}
\begin{itemize}
\item The leading term is $\sqrt{206\cdot18.2}\,L\,(M_1/H)^{1/2}\le61.3L^{-6}$. The linear terms are $\le140L^{-12}$:
  $\frac43(1+A_u)M_1\lambda'\le50L^2M_1$ and $3B\le84L^2M_1$. Then $\sum_j128(\Lambda+j\log2)^{-5}\le46.2/\Lambda^4$ and
  $\sum_j128(\Lambda+j\log2)^{-4}\le61.6/\Lambda^3$.
\item On the standard range, where $H=H_{\mathrm{std}}$ and $M_1\le\frac12H^{1/2}$, the same sums are $\le e^{-600}$; the first
  term is then $\approx18L\,e^{-2(\log L)^2}$.
\item The sums converge whenever $H(v)/M_1(v)\ge(\log v)^{8+\delta}$. This is what fixes the tolerance $M_1\le H/L^{14}$.
\end{itemize}

% 04e-lp.tex — v3 §3.6 (first part): Lemma 3.12 and Theorem 3.13 (local PNT) with full proof.
\subsection{A local prime number theorem}\label{sec:noise-lp}
% v3: §3.6 (Lemma 3.12, Theorem 3.13)

\begin{lemma}[threshold mass]\label{lem:threshold-mass}
% v3: Lemma 3.12
For squarefree $n$ put $\tau(n):=\max_{p\mid n}t(p)\prod_{q\mid n,\,t(q)\ge t(p)}q$, so that $\chi^t_y(n)=1$ if and only if
$\tau(n)<y$. Assume $|N-T|\le D:=B+\frac12H(1+\kappa_+/\log x)+4$ at the relevant heights (Corollary~\ref{cor:lag}). For $x\ge X_0^3$ and
% Logged correction 38 (CORRECTIONS-vs-v3.md): x >= X0^3 (v3: x >= X0), so that every p in case (ii) exceeds X0.
% Logged correction 36 (CORRECTIONS-vs-v3.md): +4 (lower lag constant -4); D <= H still, as B << H.
$\delta\in(0,\frac14]$,
\[
  \sum_{n:\ \tau(n)\in[x(1-\delta),x)}\frac1n\le55\,\delta\log x\,\big(\log H(x)+\log(1/\delta)+2\big)
  +14\big(K_{\mathrm{loc}}(\sqrt x)+2\big)\log x\,\sqrt{\frac{H(x)}{\delta x}} .
\]
It is used only with $\delta\ge6H(x)x^{-1/3}$, and only in Proposition~\ref{prop:g1}(B), where $x\ge X_{\mathrm{sw}}$ with
$\log X_{\mathrm{sw}}=X_0^{1/100}>3\Lambda$.
\end{lemma}

\begin{proof}
Let $p$ be the prime attaining $\tau(n)=t(p)\,m$, with $m:=\prod_{t(q)\ge t(p)}q$, and write $n=k\cdot m$, where all
primes of $k$ are $t$-below $p$. The map $n\mapsto(k,m)$ is injective. A prime $q$ with $t(q)<t(p)$ has $q<P:=p+3H(p)$,
and $q<p$ if $p<\Xm$ (then $q$ and $p$ are rational primes, with $t(q)=q$). Hence
$\sum_k1/k\le\prod_{q<P}(1-1/q)^{-1}$, which we bound as follows.
\begin{itemize}
\item If $p<\Xm$: for $p<400$ a direct computation gives $\prod_{q<p}(1-1/q)^{-1}\le1.8\log p+0.25$ (Appendix~\ref{app:cert}, \texttt{finite-n}); for $p\ge400$,
  Rosser and Schoenfeld \cite[Thm.~8, (3.29)]{RosserSchoenfeld1962} give
  $\prod_{q\le p}(1-1/q)^{-1}\le e^\gamma\log p\,(1+1/(2\log^2p))\le1.7811\log p+0.16$.
\item If $p\ge\Xm$: split at $y_0:=\Xm-1$, below which all primes are rational. The rational part is at most
  $e^\gamma\log y_0\,(1+1/(2\log^2y_0))$ by the same bound. For the rest, $-\log(1-1/q)\le1/q+1/q^2$, and
  Lemma~\ref{lem:cm}(iv) with $y=y_0$, $x=P$ gives $\sum_{y_0<q<P}1/q\le\log(\log P/\log y_0)+2.2/\log y_0$, while
  $\sum_{q>y_0}1/q^2\le3/y_0$. So the product is at most $1.7811\log P\,(1+3\cdot10^{-8})\le1.782\log p$.
\end{itemize}
In all cases $\sum_k1/k\le1.8\log p+1\le2\log x$, as $p\le x$ and $\log x\ge10^8$. Also $1/m\le2p/x$.

Write $m=p\,m'$, so that $p^2m'\in I:=[x(1-\delta),x]\,(1+O(H(p)/p))$. In case (ii) below, $H(p)/p\le\delta/10$, since
$\delta\ge6H(x)x^{-1/3}$.
\begin{enumerate}[label=(\roman*)]
\item $p\le\sqrt{\delta x/(2H(x))}$. The interval $I/p^2$ has length $\ge2D$, so it holds at most $4\delta x/p^2$
  elements $m'$ (by $|N-T|\le D$). These pairs contribute at most $4\delta x\sum_{p\le\sqrt x}1/p$ to $\sum p$.
\item $p$ larger. Then $m'\le2H(x)/\delta$. For each $m'$, the $p$-interval $\sqrt{I/m'}$ has length
  $\ell_{m'}\le0.6\,\delta\sqrt{x/m'}$. Since $p^2m'=pm=\tau(n)\,p/t(p)$ and $|t(p)-p|\le H(p)\le\delta p/10$, the factor
  $1+O(H(p)/p)$ in $I$ lies within $1\pm0.11\delta$. So, with $m'\le2H(x)/\delta$, $\delta\le\frac14$, $\delta\ge6H(x)x^{-1/3}$ and
  $x\ge X_0^3$, the interval has minimum $\ge\sqrt{(1-\delta)(1-0.11\delta)}\,\sqrt{\delta x/(2H(x))}\ge0.85\sqrt3\,x^{1/3}\ge1.47X_0$
  and endpoint ratio $\le\big((1+0.11\delta)/((1-\delta)(1-0.11\delta))\big)^{1/2}\le1.19$. Let $P_1\le\sqrt x$ be the dyadic scale
  with $P_1\le\min\sqrt{I/m'}<2P_1$; then $P_1\ge X_0$ and the interval lies in $[P_1,2.38P_1)\subset[P_1,4P_1)$. Cut it into
  $\lceil\ell_{m'}/(6H(P_1))\rceil$ windows of length $\le6H(P_1)$ and apply the local count at the dyadic scale $P_1$
  (Corollary~\ref{cor:localcount}(i)) to each. This gives at most
  $13.4\ell_{m'}+\lceil\ell_{m'}/(6H(P_1))\rceil(K_{\mathrm{loc}}(P_1)+2)\le13.5\ell_{m'}+K_{\mathrm{loc}}(\sqrt x)+2$ primes, since
  $K_{\mathrm{loc}}\le H/100$ and $K_{\mathrm{loc}}$ is non-decreasing.
% Logged correction 37 (CORRECTIONS-vs-v3.md): v3 applies the local count at the non-dyadic height sqrt(x). Each p-interval
% lies in one block [P_1,4P_1) with P_1 dyadic, so part (i) applies there and the constant K_loc(sqrt x) is unchanged.
% (Part (ii), for real P, would give K_loc(P)+K_loc(2P) and turn the constant 14 of the statement into 23.)
\end{enumerate}
Each pair $(p,m')$ contributes at most $2\log x\cdot2p/x$. Collecting, with $p\le\sqrt{x/m'}$, the sum is at most
\[
  \frac{4\log x}x\Big[4\delta x(\log\log x+1)+8.1\,\delta x\big(\log(2H(x)/\delta)+1\big)+2(K_{\mathrm{loc}}+2)\sqrt x\,\sqrt{2H(x)/\delta}\Big].
\]
Since $\log H\ge8(\log\log x)^2\ge16\log\log x$, this is at most the stated bound.
\end{proof}

\paragraph{Notation for the local prime number theorem.}
\begin{itemize}
\item $\Pi$ is the realised prime measure (prime powers weighted $1/k$), so that $N=\exp^*(\Pi)$; and
  $\Pi_0:=\log^*(\delta_1+du)$, whose density is $(1-1/u)/\log u$.
\item $T_0:=T-(\#\Abs-\int\bar\rho_{\Abs})$ is the $\Abs$-free target. Its Mellin transform is
  $\hat T_0(s)=s/(s-1)+\hat R(s)+\hat E(s)$ (Lemma~\ref{lem:ZBlow}), where $\hat R=\sum_k\hat R_k$ is the transform of the
  designed oscillation \eqref{eq:Rhat}.
\item $\theta':=\max(\theta+0.01,\ 8(\log\log X_0)^2/\log X_0+0.005)\le0.51$ is a fixed exponent (only $\theta'<1$ is
  needed), and $\Lambda=\log X_0$.
\end{itemize}

\begin{theorem}[local prime number theorem]\label{thm:lp}
% v3: Theorem 3.13
Let $y\ge X_0$, and assume:
\begin{enumerate}[label=\textup{(LP-\alph*)}]
\item $N=T_0$ below $\Xm$, and $|N(u)-T_0(u)|\le D_0(u):=B(u)+3H(u)+4+3.1u^{\theta/2}(\log u)^{-4}\le u^{\theta'}$ for
  $\Xm\le u\le y$. (On $\Good(<y)$ this holds by Theorem~\ref{thm:queue}, Corollary~\ref{cor:lag} and \eqref{eq:T1}.)
% Logged correction 34 (CORRECTIONS-vs-v3.md): 3.1 (v3: 2) in D_0, for the accumulated mass deficit of the smoothed
% absorbed elements. D_0 <= u^{theta'} keeps huge room: theta' >= 0.26, and 3H(u) <= 3exp(8(loglog u)^2 + 1e-6 log u).
\item $|\hat E(s)|\le C_8(1+\kappa|s|)e^\gamma\Lambda$ for $\sigma\ge1-8/\Lambda$, with $\log C_8<440$
  (Lemma~\ref{lem:ZBlow} with $a=8$).
\item $|\hat R(s)|\le|s|$ for $\sigma\ge1-8/\Lambda$. (By \eqref{eq:Rhat}: there $\mathrm{Re}\,w=\sigma-\theta\ge0.4$, so
  $|\hat R_k(s)|\le C\alpha_k|s|x_k^{\theta-\sigma}\le2C|s|x_k^{-0.4}$, which is summable over $k$ since $x_1\ge\exp(X_0^2)$;
  see also \eqref{eq:Rhatbound}.)
\end{enumerate}
Then there are $c_{\mathrm{zf}}\ge0.088/\Lambda$, $c_{\mathrm{LP}}:=c_{\mathrm{zf}}/5\ge0.0176/\Lambda$, an effective $K$ with
$\log K=O(\Lambda)$, and a threshold $X_{\mathrm{LP}}$ with $\log X_{\mathrm{LP}}=O(\Lambda^2\log^2\Lambda)$, such that for
$X_{\mathrm{LP}}\le v\le y$:
\begin{enumerate}[label=\textup{(\roman*)}]
\item $\displaystyle\Big|\int_{[1,v]}\frac{d(\Pi-\Pi_0)(u)}u\Big|\le K\,e^{-c_{\mathrm{LP}}\sqrt{\log v}}$;
\item $|\Pi(v)-\Pi_0(v)|\le K\,v\,e^{-c_{\mathrm{LP}}\sqrt{\log v}}$;
\item \emph{(short windows)} let $\Lambda_p:=\log v\ge76\Lambda^3$, $\eta\in[e^{-2\Lambda},\frac12]$, and let $\varphi\ge\mathbf 1_{[v,v(1+\eta)]}$ be the
  smooth window of relative width $\eta(1+2\Lambda_p^{-2})$ with transition layers $\eta\Lambda_p^{-2}$. Put
  $L_1:=2\log\Lambda_p+\log(1/\eta)+5$. Then $\int\varphi\,d\Pi\le(1+\mathrm{Rel})\int\varphi\,d\Pi_0$, with
  \begin{equation}\label{eq:lpplus}
    \mathrm{Rel}(\Lambda_p,\eta)\le18\,G(L_1)\,\Lambda_p^3\,\eta^{-1}\exp\Big(-\frac{0.9\,c_{\mathrm{zf}}\Lambda_p}{L_1+442+2\log\Lambda+\log L_1}\Big),
  \end{equation}
  where $G(L):=\sup|\log(\zeta^\sharp(s)(s-1)/s)|$ on the shifted contour up to height $e^L$. Moreover
  $G(L)\le400(L+450+2\log\Lambda)$ for $5\le L\le\log T_y$, with $\log T_y:=\log y/(51\Lambda\log\log y)-500$; so
  $\log G(L)\le\Lambda$ for such $L$ with $L\le e^{\Lambda/2}$.
% Logged correction 24 (CORRECTIONS-vs-v3.md): (iii) restricted to Lambda_p >= 76 Lambda^3; G defined via the integrand
% log(zeta#(s)(s-1)/s) with the explicit bound above (v3: L_1 <= 3000 Lambda and log G <= Lambda unconditionally).
\end{enumerate}
\end{theorem}

\begin{proof}[Proof of Theorem~\ref{thm:lp}]
\begingroup\emergencystretch=3em
This is Landau's argument \cite[Lemmas~6.2--6.3 and the proof of Thm.~6.6]{MontgomeryVaughan2007}, with four changes. It is
run on an extension of $N$ beyond $y$, which is not a Beurling system. Positivity is therefore imported from a finite
Euler product. Landau's discs have radius $\asymp1/\Lambda$. And the region near the pole is handled with the fine form of
$\hat E$. Throughout, $s=\sigma+it$.\par\endgroup

\emph{Step 1: the extension.} Put $D:=N-T_0$, so $D=0$ below $\Xm$ and $|D(u)|\le u^{\theta'}$ on $[\Xm,y]$ by (LP-a). Define
$N^\sharp:=N$ on $[1,y]$ and $N^\sharp(u):=N(y)+T_0(u)-T_0(y)$ for $u>y$. Then $D^\sharp:=N^\sharp-T_0$ equals
$D(\min(u,y))$, and
\[
  \zeta^\sharp(s):=\int_{[1,\infty)}u^{-s}\,dN^\sharp=\hat T_0(s)+\hat D^\sharp(s),\qquad
  \hat D^\sharp(s)=s\int_{\Xm}^yD(u)u^{-s-1}du+D(y)y^{-s}.
\]
So $\hat D^\sharp$ is entire, with $|\hat D^\sharp(s)|\le|s|(\Xm)^{\theta'-\sigma}/(\sigma-\theta')+y^{\theta'-\sigma}$ for
$\sigma>\theta'$. Let $\Pi^\sharp:=\log^*(dN^\sharp)$, a signed measure (the series converges in the weighted norm
$\int u^{-\sigma}|\cdot|$ for large $\sigma$). Since $\log^*$ is causal, $\Pi^\sharp=\Pi$ on $[1,y]$.

\emph{Step 2: upper bounds on $[1,y]$.} Lemma~\ref{lem:cm} holds on $[1,y]$: (LP-a) and \eqref{eq:T1} give \textup{(H$_{\mathrm{BS}}$)} there. Write $d\psi=\log u\,d\Pi$ and
$S(z)=\int_{[1,z]}d\psi/u\le\log z+0.43$. By Lemma~\ref{lem:cm}(iv) and the values for $\N$,
$\int_{[1,y]}d\Pi/u\le\log\log y+4.5$. Let $N_y:=\exp^*(\Pi|_{[1,y]})$ count the $y$-smooth elements, so $N_y=N$ on
$[1,y]$, and put $P_y(s):=\int u^{-s}dN_y=\exp(\int_{[1,y]}u^{-s}d\Pi)$. Then $P_y(1)\le e^{4.5}\log y$.

\emph{Step 3: positivity transfer.} Put $\Delta:=\zeta^\sharp-P_y=\int_{(y,\infty)}u^{-s}(dT_0-dN_y)$. For $\sigma=1+\eta$
with $0<\eta\le1/8$, Rankin's trick ($u^{-\eta}\le y^{-\eta/2}u^{-\eta/2}$ for $u>y$) and (LP-a) give
\begin{gather*}
  P_y(\sigma)\le1/\eta+6+y^{-\eta/2}P_y(1),\qquad|\Delta(s)|\le2y^{-\eta}/\eta+y^{-\eta/2}P_y(1),\\
  |\Delta'(s)|\le2y^{-\eta}(\log y/\eta+1/\eta^2)+y^{-\eta/2}P_y(1)\,S(y),\qquad|P_y(s)|\ge1/P_y(\sigma),\qquad
  |P_y'/P_y(s)|\le S(y).
\end{gather*}
The last two hold because $\Pi\ge0$. Set $\eta_y:=17\log\log y/\log y$, so that $y^{-\eta/2}\le(\log y)^{-8.5}$ for
$\eta\ge\eta_y$. Then, for $1+\eta_y\le\sigma\le1+1/8$,
\begin{equation}\label{eq:lp-transfer}
  P_y(\sigma)\le2/\eta,\qquad|\Delta/P_y|\le10^{-2},\qquad\Big|\frac{d}{ds}\log(1+\Delta/P_y)\Big|\le(\eta\log y)^{-1}.
\end{equation}
Let $F:=(dT_0-dN_y)|_{(y,\infty)}*(dN_y)^{*-1}$. Since $|(dN_y)^{*-1}|\le dN_y$, the same bounds give
$\int u^{-\sigma}|dF|\le10^{-2}$. Hence $\Pi^\sharp=\Pi|_{[1,y]}+\log^*(\delta_1+F)$ satisfies $\int u^{-\sigma}|d\Pi^\sharp|<\infty$,
$\zeta^\sharp$ has no zeros in $\sigma\ge1+\eta_y$, and $\log\zeta^\sharp=\log P_y+\log(1+\Delta/P_y)$ is the Mellin transform of
$\Pi^\sharp$ there. The coefficients of $-P_y'/P_y$ form the positive measure $d\psi|_{[1,y]}$. So for the nonnegative
polynomial
\[
  (1+\cos\phi)^4=\tfrac18\big(35+56\cos\phi+28\cos2\phi+8\cos3\phi+\cos4\phi\big),
\]
with $(a_0,\dots,a_4):=(35,56,28,8,1)$, we get
\begin{equation}\label{eq:lp-341}
  \sum_{k=0}^4a_k\,\mathrm{Re}\Big(-\frac{\zeta^{\sharp\prime}}{\zeta^\sharp}\Big)(\sigma+ikt)\ge-\frac{128\,\eps_f}{\eta},
  \qquad\eps_f:=5\cdot10^{-4}\ \ (\ge1/\log y),\qquad\sigma=1+\eta\ge1+\eta_y .
\end{equation}

\emph{Step 4: the pole term.} By Step~3 and Lemma~\ref{lem:cm}(ii), with partial summation,
$-P_y'/P_y(1+\eta)=\int_{[1,y]}u^{-\eta}dS(u)\le\eta\int_1^\infty(\log u+0.43)u^{-1-\eta}du=1/\eta+0.43$. Hence, by
\eqref{eq:lp-transfer},
\begin{equation}\label{eq:lp-pole}
  -\frac{\zeta^{\sharp\prime}}{\zeta^\sharp}(1+\eta)\le\frac{1+10^{-8}}\eta+0.43\qquad(\eta\ge\eta_y).
\end{equation}
This replaces the usual use of the pole of $\zeta^\sharp$; the crude bound $|\hat E|\le C_8e^\gamma\Lambda$ is useless near
$s=1$.

\emph{Step 5: Landau's lemma.} Let $f$ be analytic on $|z-z_0|\le r$, with $|f|\le M|f(z_0)|$ there, and suppose every zero
of $f$ in the disc has real part $\le\mathrm{Re}\,z_0$. Then for any zero $\rho_1$ in the disc,
\[
  -\mathrm{Re}\frac{f'}f(z_0)\le\frac2r\log M-\mathrm{Re}(z_0-\rho_1)\Big(\frac1{|z_0-\rho_1|^2}-\frac1{r^2}\Big).
\]
\emph{Proof.} Divide $f$ by the Blaschke factors of its zeros, normalised on the circle of radius $r$; the quotient $g$ has
no zeros and $|g/g(z_0)|\le M$ on the circle. Borel--Carath\'eodory for $h=\log(g/g(z_0))$ on radii $R\uparrow r$ gives
$|h'(z_0)|\le2\log M/r$. Each zero contributes $\mathrm{Re}(z_0-\rho)(1/|z_0-\rho|^2-1/r^2)\ge0$ to $-\mathrm{Re}\,f'/f(z_0)$;
drop all but $\rho_1$.

\emph{Step 6: the region near the pole.} For $|t|\le X_0$ and $\sigma\ge1-8/\Lambda$, the proof of
Lemma~\ref{lem:ZBlow} gives the fine form
$\hat E(s)=\zeta(s)-1/(s-1)-1+O(X_0^{-\sigma})+O(\kappa(1+|s|)e^{\mathrm{Ein}(8)}e^\gamma\Lambda)$. Since
$s/(s-1)=1+1/(s-1)$, \eqref{eq:Rhatbound} (that is, $|\hat R(s)|\le|s|e^{-0.39X_0^2}$) and Step~1 give
\[
  \zeta^\sharp(s)=\zeta(s)+E(s),\qquad|E(s)|\le3X_0^{-0.3}(1+|t|)\qquad(|t|\le X_0,\ \sigma\ge1-8/\Lambda),
\]
using $\kappa\le X_0^{-1/3}$ and $e^{\mathrm{Ein}(8)}\le e^{438}$. On $|s-1|\le1/8$ the Laurent series gives
$|\zeta(s)-1/(s-1)|\le\gamma+0.02\le0.6$ (as $|\gamma_1|\le0.073$), so $|(s-1)\zeta^\sharp(s)-1|\le0.6/8+10^{-9}<1$, and
$\zeta^\sharp$ has no zero there. This covers every point with $|t|<1/16$ and $1-8/\Lambda\le\sigma\le1.1$; for
$\sigma>1.1$, $\zeta^\sharp\ne0$ by Step~3.

\emph{Step 7: the zero-free region for $|t|\ge1/16$.} Take Landau discs of radius $r:=a/\Lambda$ with $a:=6.8\le8$, so that
they stay in $\sigma\ge1-8/\Lambda$, where (LP-b) and (LP-c) apply. For $|t|\ge1/16$ they do not meet the pole. Put
\[
  L_t:=\log(|t|+3),\qquad X(t):=L_t+442+2\log\Lambda,\qquad Y(t):=X(t)+\log X(t)+0.7 .
\]
\begin{itemize}
\item \emph{Growth.} Let $|t|\ge1/16$, $\eta_y\le\eta\le10^{-3}/\Lambda$, $k\in\{1,2,3,4\}$ and $z_0=1+\eta+ikt$. Step~3 gives
  $|\zeta^\sharp(z_0)|\ge(1-10^{-2})/P_y(1+\eta)\ge0.495\,\eta$. On the disc $|z-z_0|\le r$:
  \begin{itemize}
  \item $|z/(z-1)|\le1+1/(k|t|-r)\le17.1$;
  \item $|\hat R(z)|\le|z|\le k|t|+1.01$, by (LP-c);
  \item $|\hat E(z)|\le e^\gamma C_8\Lambda+10^{-10^7}|z|$, by (LP-b), since $\kappa\le X_0^{-1/3}$ and $\log C_8<440$;
  \item $|\hat D^\sharp(z)|\le2.1|z|X_0^{-0.48}+1$, by Step~1, since $\sigma-\theta'\ge0.48$.
  \end{itemize}
  So $|\zeta^\sharp|\le2(1+3X_0^{-0.48})\max(k|t|+19.2,\ e^\gamma C_8\Lambda)$ there, for every $|t|\ge1/16$. Now
  $\log(4|t|+19.2)\le L_t+1.857$ and $\log(e^\gamma C_8\Lambda)<440.578+\log\Lambda$; since $\max(a,b)\le a+b-2.975$ when
  $a,b\ge2.975$, the ratio $M_k:=\max_{|z-z_0|\le r}|\zeta^\sharp(z)|/|\zeta^\sharp(z_0)|$ satisfies (the constant is
  $440.855+\log(1+3X_0^{-0.48})\le440.86$)
  \begin{equation}\label{eq:lp-growth}
    \log M_k\le L_t+440.86+\log\Lambda+\log(1/\eta).
  \end{equation}
\item \emph{The inequality.} Suppose $\zeta^\sharp(\beta+it)=0$ with $1-\beta\le\eta/\kappa_L$ for a parameter $\kappa_L>0$, and
  $\eta\ge\eta_y$. Apply Step~5 at $z_0=1+\eta+it$ with this zero, and at $z_0=1+\eta+ikt$ ($k=2,3,4$) with no zero term.
  This is legitimate because all zeros have real part $<1+\eta_y\le1+\eta$ (Step~3), so the dropped terms are $\ge0$. Insert
  these and \eqref{eq:lp-pole} into \eqref{eq:lp-341}. Using $1/(\eta+1-\beta)\ge\kappa_L/((\kappa_L+1)\eta)$, we get
  \[
    \frac{N_0(\kappa_L)}{\eta}\le15.05+\frac2r\sum_{k=1}^4a_k\log M_k+\frac{56\,\eta(1+1/\kappa_L)}{r^2},\qquad
    N_0(\kappa_L):=\frac{56\kappa_L}{\kappa_L+1}-35(1+10^{-8})-128\eps_f .
  \]
\item \emph{The constant.} Since $\sum_{k=1}^4a_k=93$, the inequality gives, as $|t|\to\infty$, the zero-free width
  $\mathrm{merit}\cdot r/(2\log|t|)$, with
  \[
    \mathrm{merit}:=\max_{\kappa_L}\frac{N_0(\kappa_L)}{93\,\kappa_L}=0.0262291\qquad(\kappa_L\approx3.79).
  \]
  So $\mathrm{merit}\cdot r/2=0.0891790/\Lambda$, and $\mathrm{merit}\cdot a/2\ge0.088$ needs $a\ge6.71$. We fix $\kappa_L:=3.79$,
  for which $N_0\ge9.2449$, and put
  \[
    c_{\mathrm{zf}}:=0.088/\Lambda .
  \]
  The slack between $0.08917$ and $0.088$ absorbs the lower-order terms of \eqref{eq:lp-growth} at finite height, as follows.
  The constant $\mathrm{merit}$ does not depend on $\Lambda$.
\item \emph{Explicit width.} Let $|t|\ge1/16$, and suppose that $\eta:=0.33769/(\Lambda Y(t))$ satisfies $\eta\ge\eta_y$. Then
  $\zeta^\sharp(\sigma+it)\ne0$ for
  \begin{equation}\label{eq:lp-width}
    \sigma\ge1-\delta_t,\qquad\delta_t:=\eta/\kappa_L=\frac{0.0891}{\Lambda\,Y(t)} .
  \end{equation}
  Indeed, otherwise the inequality holds with this $\eta$. Here $\log(1/\eta)=\log\Lambda+\log Y+1.0857$, and
  $\log Y-\log X\le(\log X+0.7)/X\le0.0144$ (as $X\ge479.96$), so the right side of \eqref{eq:lp-growth} is at most
  $X-0.054+\log Y\le X+\log X-0.039\le Y$.
  Multiply the inequality by $\eta$, and use $\eta\Lambda=0.33769/Y\le7.1\cdot10^{-4}$: the right side becomes at most
  $(2\cdot93/6.8)\cdot0.33769+10^{-6}\le9.2369<9.2449\le N_0$, a contradiction.
\item \emph{Two consequences.} Write $\delta_{\mathrm{zf}}(t):=\delta_t$.
  \begin{itemize}
  \item Since $0.088(X+\log X+0.7)\le0.0891(X+18)$ for $X\ge479$, we have $\delta_{\mathrm{zf}}(t)\ge c_{\mathrm{zf}}/(L_t+460+2\log\Lambda)$.
  \item \emph{Finite height.} Let $L\ge5$ and $|t|\le e^L$. Then $L_t\le L+3e^{-L}$, and
    $0.088(X+\log X+0.7)\le0.0891(L+442+2\log\Lambda+\log L)$: the difference of the two sides increases with $L$ and with
    $\Lambda$, and is positive at $L=5$, $\Lambda=10^8$. So at height $e^L$ the zero-free width is
    $\delta_{\mathrm{zf}}\ge c_{\mathrm{zf}}/(L+442+2\log\Lambda+\log L)$. Every use below has $L\ge40$.
  \end{itemize}
\item \emph{The height constraint.} Step~3 requires $\eta=\kappa_L\delta_{\mathrm{zf}}(t)\ge\eta_y=17\log\log y/\log y$, and under the
  present scaling this binds. Assume
  \begin{equation}\label{eq:lp-yrange}
    \log y\ge10^4\Lambda^2\log\log y,\qquad\text{so that}\quad Q_y:=\frac{\log y}{\Lambda\log\log y}\ge10^4\Lambda ,
  \end{equation}
  and put $\log T_y:=Q_y/51-500$. For $|t|\le T_y$ we have $X(t)\le Q_y$ and
  $Y(t)\le Q_y/51-57.3+2\log\Lambda+\log Q_y\le Q_y/50.35$, because $(1/50.35-1/51)Q_y\ge2.53\cdot10^{-4}Q_y\ge2\log\Lambda+\log Q_y$.
  Since $0.33769\cdot50.35\ge17.002$, this gives $\eta=0.33769/(\Lambda Y)\ge17/(\Lambda Q_y)=\eta_y$. So the zero-free region is
  \begin{equation}\label{eq:lp-zf}
    \sigma\ge1-\frac{c_{\mathrm{zf}}}{\log(|t|+3)+460+2\log\Lambda}\qquad\text{for }|t|\le T_y,\qquad
    \log T_y:=\frac{\log y}{51\Lambda\log\log y}-500 .
  \end{equation}
  Together with Step~6 this covers all $|t|\le T_y$. Every application below satisfies \eqref{eq:lp-yrange}.
\item By (LP-c), $\hat R$ enters only through the growth bound, like $\hat D^\sharp$. The region \eqref{eq:lp-zf} is
  consistent with the designed zeros, which have $1-\beta_k\asymp\log\log\gamma_k/\log\gamma_k$ (\S\ref{sec:zeros}).
\end{itemize}

\emph{Step 8: $\log\zeta^\sharp$ in the region.} Assume \eqref{eq:lp-yrange}, and let $\mathfrak G_1(s):=\log(\zeta^\sharp(s)(s-1)/s)$,
the branch that is real on $(1,\infty)$; it is analytic in the zero-free region, including near $s=1$. Fix $f:=0.9$, and
for $1/16\le|t|\le T_y$ let $s_1:=1+\eta_t+it$ with $\eta_t:=\kappa_L\delta_t$. Consider the discs centred at $s_1$ of
radii $R_1=\eta_t+\frac{1+f}2\delta_t$ and $r_1=\eta_t+f\delta_t$. They avoid the pole, and they lie in the zero-free
region, since $\delta_{\mathrm{zf}}$ changes by a relative $10^{-8}$ over them. The function
$h:=\log\zeta^\sharp-\log\zeta^\sharp(s_1)$ vanishes at $s_1$, and $\mathrm{Re}\,h\le\log M_1\le Y(t)$ on the larger disc, by
\eqref{eq:lp-growth} with $\eta=\eta_t$ (as in Step~7). Borel--Carath\'eodory gives
$|h|\le\frac{2r_1}{R_1-r_1}Y(t)=\frac{4(\kappa_L+f)}{1-f}Y(t)\le187.6\,Y(t)$ on the smaller disc. At the anchor, Step~3 gives
$|\log\zeta^\sharp(s_1)|\le\log P_y(1+\eta_t)+0.02\le\log(2/\eta_t)+0.02\le0.5\,Y(t)$. Finally
$|\log((s-1)/s)|\le\log17+\pi\le6$ for $|t|\ge1/16$ and $\sigma\le1.1$. Hence
\[
  |\mathfrak G_1(s)|\le190\,Y(t)\qquad\text{for }|s-s_1|\le r_1 .
\]
A point $s=\sigma+it$ of the contour $\sigma=1-f\delta_{\mathrm{zf}}(T_1)$, with $1/16\le|t|\le T_1\le T_y$, satisfies
$|s-s_1|=\eta_t+f\delta_{\mathrm{zf}}(T_1)\le r_1$, because $\delta_{\mathrm{zf}}$ decreases in $|t|$. For $|t|<1/16$ the contour
lies in $|s-1|\le1/16+10^{-8}$. There Step~6 gives $|\log((s-1)\zeta^\sharp(s))|\le-\log(1-0.075)\le0.078$, and
$|\log s|\le-\log(1-|s-1|)\le0.0646$, so $|\mathfrak G_1|\le0.1426$. Since
$Y\le2X$, it follows that for $5\le L\le\log T_y$ the quantity $G(L)$ of (iii) satisfies
\[
  G(L)\le190\,Y(e^L)\le400\,(L+450+2\log\Lambda);
\]
in particular $\log G(L)\le\Lambda$ for $L\le e^{\Lambda/2}$.

\emph{Step 9: Perron's formula; proof of (i) and (ii).} Let
$\mathfrak G(s):=\log\big(\zeta^\sharp(s+1)\,s/(s+1)\big)$. It is analytic in the shifted zero-free region, $\mathfrak G(0)=0$
(the residue of $\zeta^\sharp$ at $1$ is $1$, and $\zeta^\sharp>0$ on $(1,\infty)$), and $\mathfrak G$ is the Mellin transform of
$(\Pi^\sharp-\Pi_0)/u$ on $\mathrm{Re}\,s\ge\eta_y$. Let $\Phi^\sharp(w):=\int_{[1,w]}d(\Pi^\sharp-\Pi_0)/u$; it equals the left
side of (i) for $w\le y$.
\begin{itemize}
\item Take $h_0\in C_c^\infty(-2,-1)$ with $h_0\ge0$, $\int h_0=1$, and put $h:=Th_0(T\cdot)$. Smoothing to the left keeps
  everything inside $[1,y]$:
  \[
    \Phi^h(v):=\int\Phi^\sharp(ve^\tau)h(\tau)\,d\tau=\frac1{2\pi i}\int_{(\eta_y)}\mathfrak G(s)\,\hat h_0(s/T)\,v^s\,\frac{ds}s,
    \qquad|\hat h_0(w)|\le C_me^{2|\mathrm{Re}\,w|}(1+|w|)^{-m}.
  \]
\item Shift the contour to $\mathrm{Re}\,s=-f\delta_{\mathrm{zf}}(T_1)$ for $|t|\le T_1\le T_y$. There is no residue at $s=0$,
  because $\mathfrak G(0)=0$. Take $T:=e^{\sqrt{\log v}}$, $T_1:=e^{2\sqrt{\log v}}(\log y)^{2/m}$ and $m:=19$.
\item The contributions are as follows. The shifted line gives at most
  $K_8\,v^{-f\delta_{\mathrm{zf}}(T_1)}\log^2(T_1+3)$. The horizontal segments, and the line $\mathrm{Re}\,s=\eta_y$ beyond
  $T_1$, give at most
  \[
    K_9v^{\eta_y}(\log\log y)(T/T_1)^m\le e^{-\sqrt{\log v}},
  \]
  using $v^{\eta_y}\le(\log y)^{17}$. The
  smoothing error is
  \[
    |\Phi^\sharp(v)-\Phi^h(v)|\le\int_{(ve^{-2/T},v]}\frac{d\Pi+d\Pi_0}u\le\frac6T+3v^{\theta'-1},
  \]
  using $d\Pi\le dN$.
\item The heights used lie below $T_v$ once $\log v\ge1.1\cdot10^4\Lambda^2(\log\log v)^2$; here we apply the argument with
  $y':=v$ in place of $y$, which is allowed since (LP-a) holds on $[1,v]$; and $y'=v$ satisfies \eqref{eq:lp-yrange}.
  For $\log v\ge(1000+4\log\Lambda)^2$ the offset in
  \eqref{eq:lp-zf} at most doubles the denominator. Together this gives (i) with $c_{\mathrm{LP}}=c_{\mathrm{zf}}/5$. The
  threshold $X_{\mathrm{LP}}$ is chosen above both heights, so $\log X_{\mathrm{LP}}=O(\Lambda^2\log^2\Lambda)$, and
  $\log K=O(\Lambda)$.
\item (ii) is the same argument with $\mathfrak G_1(s):=\log(\zeta^\sharp(s)(s-1)/s)$, the kernel
  $v^s\hat h_0((s-1)/T)/s$, and the contour near $\sigma=1$.
\end{itemize}

\emph{Step 10: short windows; proof of (iii).} With the smooth window $\varphi$, write $\int\varphi\,d(\Pi^\sharp-\Pi_0)$ as the
contour integral of $\log(\zeta^\sharp(s)(s-1)/s)$ against the Mellin transform $\tilde\varphi$ of $\varphi$. Since $\varphi$ has
transition layers $\eta\Lambda_p^{-2}$, $\int|\tilde\varphi(\sigma+it)|\,dt\le4.04\,C_\varphi v^\sigma\Lambda_p^2$ with
$C_\varphi\le4$, and the integrand is negligible beyond height $T_1=e^5\Lambda_p^2/\eta$, i.e.\ $L_1=\log T_1$. No $1/T$
smoothing term arises, so $T_1$ is fixed by $\eta$ alone. Shifting to $\sigma=1-f\delta_{\mathrm{zf}}(T_1)$ and using Step~7 for
the width at height $e^{L_1}$ gives \eqref{eq:lpplus}.
\begin{itemize}
\item \emph{Range.} Step~7 gives the width at height $e^{L_1}$ if \eqref{eq:lp-yrange} holds and $L_1\le\log T_y$. Both hold for
  $\Lambda_p\ge76\Lambda^3$. First, $\log y\ge\Lambda_p\ge76\Lambda^3$, and $\log y/\log\log y$ is increasing, so
  \eqref{eq:lp-yrange} holds. Second, $L_1+500\le2\Lambda+2\log\Lambda_p+505$, while $\log T_y+500\ge\Lambda_p/(51\Lambda\log\Lambda_p)$;
  the latter exceeds the former for every $\Lambda_p\ge76\Lambda^3$, since both sides increase and at $\Lambda_p=76\Lambda^3$ the
  right side is $\ge\Lambda^2/(3\log\Lambda+5)\gg2\Lambda$. Also $L_1\ge2\log\Lambda_p\ge40$. So $5\le L_1\le\log T_y$, and the
  bound on $G(L_1)$ of Step~8 applies.
\item \emph{Size.} For $\eta\ge e^{-2\Lambda}$ and $L_1\le e^{\Lambda/2}$ (so that $\log G(L_1)\le\Lambda$ by Step~8),
  $\mathrm{Rel}\le0.1$ needs $0.9c_{\mathrm{zf}}\Lambda_p/(L_1+442+2\log\Lambda+\log L_1)\ge3\Lambda+3\log\Lambda_p+6$, that is,
  $\Lambda_p\ge76\Lambda^3$. If $L_1>e^{\Lambda/2}$, then $\log\Lambda_p\ge e^{\Lambda/2}/3$ and $G(L_1)\le\Lambda_p$ by Step~8, and the
  same comparison holds with room to spare. The threshold $76\Lambda^3$ is below $\Lambda_1=100\Lambda^3$ of
  Lemma~\ref{lem:si}.\qedhere
\end{itemize}
\end{proof}

\begin{remark}\label{rem:lp}
Positivity cannot be dropped: $N=\lfloor u\rfloor+\frac23\cos(3\log u+\pi/2)$ has $|N-u|\le2$, but its $\zeta^\sharp$ has a zero
near $1+3i$. The theorem uses only the realised $N$ on $[1,y]$, so it holds pathwise on $\Good(<y)$, and the adaptedness of $T$
plays no role. Parts (i) and (ii) are used only in Proposition~\ref{prop:g1}(B), whose range starts at
$X_{\mathrm{sw}}\gg X_{\mathrm{LP}}$; part (iii) is used in Lemma~\ref{lem:si}.
\end{remark}

% 04f-bridge.tex — v3 §3.6 (second part): Prop 3.14 (BR), the X0 band (input of 04x-x0band.tex), (Osc), Lemma 3.15,
% Prop 3.16 (G1').
\subsection{The bridge, the jump at \texorpdfstring{$X_0$}{X0}, and positivity of the sharp reference}\label{sec:noise-bridge}
% v3: §3.6 (Proposition 3.14 – Proposition 3.16)

\begin{proposition}[the bridge: operator versus sharp]\label{prop:br}
% v3: Proposition 3.14
Let $o_u$ be the supremum of $\osc(\Tc)$ over windows of length $2.04H(u)$ meeting $[u,2u)$: the $H$-window oscillation of
the jump-free source. Only the kernels $C_n$ with $n\ge2$ use it, so it is read at heights $\le x/2+1.1H(x)$. For $x\ge2\Xm$
(in particular at every point of $W_x$ for $x\ge2X_0$, since $x-H(x)/2\ge2X_0-H(2X_0)/2>2\Xm$; Theorem~\ref{thm:g2} and
Proposition~\ref{prop:pointwise}(b) use it there) split
% Logged correction 90 (CORRECTIONS-vs-v3.md): stated for x >= 2X_0 but applied at the points of W_x; the proof needs only
% x >= 2X_0^- (starts >= X_0 for Corollary kr-weights(b), and x/n >= X_0^- for every level used).
\[
  (\Rres T')(x)-\rhosh(x)=\mathrm{BR}_c(x)+J_{X_0}(x),\qquad
  J_{X_0}(x):=\jz\Big[(\Rres\mathbf 1_{[X_0,\infty)})(x)-\sum_nc_n(x)\mathbf 1[x/n\ge X_0]/n\Big],
\]
where $\mathrm{BR}_c:=\Rres\Tc-\sum_nc_n\Tc(\cdot/n)/n$. Then
\[
  \log x\cdot|\mathrm{BR}_c(x)|\le(5.1C_N+C_S)\sum_{u\text{ dyadic}}\log(2u)\,o_u+8.4\,C_NE_\theta=:E_{\mathrm{BR}} .
\]
\end{proposition}

\begin{proof}
Write $\mathrm{BR}_c(x)=\sum_n\frac1n\big[\int\Tc\,dC_n^{\mathrm{nom}}-c_n\Tc(x/n)\big]+\sum_n\frac1n\int\Tc\,d\tilde E_n$, in the
notation of \S\ref{sec:kernels}.
\begin{itemize}
\item \emph{The Kantorovich--Rubinstein part.} $\int\Tc\,dC_n^{\mathrm{nom}}-c_n\Tc(x/n)=\int(\Tc-\Tc(x/n))\,dC_n^{\mathrm{nom}}$, an
  integral over $[x/n\pm R_n]$ with $R_n\le1.02H(x/n)$. Let $u$ be the dyadic scale with $x/n\in[u,2u)$; then $u\ge X_0/2$,
  since $x/n\ge\Xm$. Each half $[x/n-R_n,x/n]$, $[x/n,x/n+R_n]$ has length $\le1.02\cdot2^{\eps}H(u)\le2.04H(u)$ and meets
  $[u,2u)$, so the integral is at most $2o_u\TV(C_n^{\mathrm{nom}})$. Summing over $n$ with Corollary~\ref{cor:kr-weights}(b)
  (start $x$), this part is at most $\sum_u2o_uW^{\mathrm{nom}}_x(u)\le2C_N\sum_u\log(2u)\,o_u/\log x$, which is inside the
  first term.
\item \emph{The near-threshold part.} For $\tau$-near steps, by the Duhamel form in Proposition~\ref{prop:true-kernels} and
  $|\Tc|\le2$, it is at most $4.2\sum_d(\TV(C_d)/d)\,\theta(x/d)\le8.4C_NE_\theta/\log x$ (Lemma~\ref{lem:tau-near}). For
  $X_0$-near (terminal) steps, $\Tc$ is within $o_{X_0}$ of $\Tc(X_0)=0$ on the support, so their part is at most
  $o_{X_0}\cdot2C_N\cdot1.52\log(2X_0)/\log x$, which is inside the first term.
\end{itemize}
Only the oscillation of the smooth source enters, never the $H$-scale roughness of the noise. The jump
$\jz\mathbf 1_{[X_0,\infty)}$ has no small oscillation, and its whole effect is $J_{X_0}$, which is treated next.
\end{proof}

% 04x-x0band.tex -- v3 §3.6: the jump of the target at X_0 (Lemma 3.14A = XB, its three ranges, Lemma 3.14B = SI(Λ)).
% Input from 04-noise (bridge subsection), after Proposition 3.14.
\subsubsection{The jump of the target at $X_0$}\label{sec:x0band}
% v3: §3.6, Lemmas 3.14A and 3.14B

The source $T'$ vanishes below $X_0$ and equals $\jz=1-\rhodet(X_0)$ just above it (\S\ref{sec:noise-notation}). A jump
cannot be handled by the oscillation bounds of Proposition~\ref{prop:br}: the kernels $C_n^x$ with $x/n$ within about
$H(X_0)$ of $X_0$ straddle it, and their weight in a dyadic block is of order $\log X_0/\log x$, which is not small
compared with the margin: $\rhosh(x)\log x\ge0.97-o(1)$ for $x\ge X_0^{1.01}$, but only $\ge0.82-o(1)$ on
$[X_0+H(X_0),X_0^{1.01}]$ (Proposition~\ref{prop:g1} and Lemma~\ref{lem:margin}(d)).
% Logged correction 47: piecewise margin
We therefore treat the jump separately and exactly. Recall
$J=\jz\mathbf 1_{[X_0,\infty)}$, $\Tc=T'-J$, and put
\[
  \rhosh_1:=(1+\Kpt)^{-1}\mathbf 1_{[X_0,\infty)},\qquad\text{so that}\qquad \rhosh=\rhosh_{\mathrm c}+\jz\rhosh_1,
  \qquad J_{X_0}=\jz\bigl(\Rres\mathbf 1_{[X_0,\infty)}-\rhosh_1\bigr),
\]
with $\rhosh_{\mathrm c}:=\sum_n c_n\Tc(\cdot/n)/n$ and $J_{X_0}$ as in Proposition~\ref{prop:br}. Throughout this
subsection
\[
  \Lambda:=\log X_0\ge10^8,\qquad \delta:=H(X_0)/X_0=e^{8(\log\Lambda)^2-\Lambda},\qquad \tau_0:=\log(4.01X_0),\qquad N:=y/X_0 .
\]
By \eqref{eq:RH}, $\eps\le0.03$. Near $X_0$ the jitter width is the standard one, and its local exponent
$\eps_*:=\sup_{X_0\le v\le X_0^2}d\log H/d\log v\le16\log\Lambda/\Lambda\le2.95\cdot10^{-6}$ ($16\log\log v/\log v$ is decreasing), so that
$H(vw)/H(v)\le w^{\eps_*}$ for $X_0\le v\le vw\le X_0^2$. By Huxley's theorem \eqref{eq:huxley}, $\jz\le1.0001/\Lambda$
(\S\ref{sec:noise-notation}).

\begin{lemma}[The $X_0$ jump in exact chain form]\label{lem:xb}
% v3: Lemma 3.14A (Lemma XB)
% Logged correction 90 (CORRECTIONS-vs-v3.md): stated for y >= 2X_0, but used at every y in W_x with x >= 2X_0 (eq:jX0), so down
% to 2X_0 - H/2. The proof uses no lower bound on y; below 2X_0^- the left side of eq:xb vanishes.
Let $y\ge2\Xm$. For a level $n$ write $H_n:=H(y/n)$. For a pair $(d,w)$ with $d,w\in\NN$, $d,w\ge2$, $n=dw$, let
$D_d:=C_d^{y,\nom}-c_d\,\delta_{y/d}$ (a measure of mass $0$; $C^{y,\nom}_d$ as in \S\ref{sec:kernels-steps}), and define
\begin{align*}
  \mathbb B_0&:=\{n:\ |y/n-X_0|\le H_n/2\},\\
  \mathbb B_1&:=\{(d,w):\ \text{the step $d\to n$ is $X_0$-near, i.e. } |y/(dw)-\Xm|\le R_d/w\},\\
  \mathbb B_2&:=\{(d,w)\notin\mathbb B_1:\ |y/n-X_0|\le(1+\vartheta_{X_0})H_n/2+R_d/w\}.
  % CORRECTIONS #29 (factor 1+ϑ_{X₀})
\end{align*}
Then, on $\Good(<y)$,
\begin{equation}\label{eq:xb}
  \bigl|\Rres\mathbf 1_{[X_0,\infty)}(y)-\rhosh_1(y)\bigr|\le\frakB(y)+\frac{4.3\,C_N E_\theta}{\log y},
\end{equation}
\begin{equation}\label{eq:frakB}
  \frakB(y):=\sum_{n\in\mathbb B_0}\frac{|c_n|}{2n}+\sum_{(d,w)\in\mathbb B_1}\frac{\TV(D_d)}{2dw}
  +\lambda_0\sum_{(d,w)\in\mathbb B_2}\frac{\KRn{D_d}}{dw\cdot wH_n}.
\end{equation}
By Lemma~\ref{lem:tau-near}, $E_\theta\le10^4\Lambda^2e^{8(\log\Lambda)^2-\Lambda}$, so with $C_N\le158$
\begin{equation}\label{eq:tau-term-bound}
  4.3\,C_NE_\theta\le4.3\cdot158\cdot10^4\,\Lambda^2e^{8(\log\Lambda)^2-\Lambda}\le e^{-0.99\Lambda}\qquad(\Lambda\ge10^8),
\end{equation}
since $15.8+2\log\Lambda+8(\log\Lambda)^2\le0.01\Lambda$. If no step from the start $y$ is
$\tau$-near, the last term of \eqref{eq:xb} may be omitted.
\end{lemma}

\begin{proof}
All kernels below start at $y$; write $C_n=C_n^y$ for the true and $C_n^{\nom}$ for the nominal kernels, so that
$\Rres\mathbf 1_{[X_0,\infty)}(y)=\sum_n n^{-1}C_n([X_0,\infty))$ by \eqref{eq:R-kernels}, and
$\rhosh_1(y)=\sum_n n^{-1}c_n\mathbf 1[y/n\ge X_0]$.

\emph{Step 1: the decomposition of Proposition~\ref{prop:true-kernels}.} Write $\tilde E_n:=C_n-C_n^{\nom}$. Then
$\tilde E_n=S_n+\sum_{d<n}P^{\nom}_{d\to n}\tilde E_d$ with sources $S_n=S_n^\tau+S_n^{X_0}$ on $\tau$-near and $X_0$-near
steps, where $P_{d\to n}$, $P^{\nom}_{d\to n}$ are the step operators of \S\ref{sec:true-kernels}. By Lemma~\ref{lem:near-structure}, the $X_0$-near sources sit at terminal levels and are never
transferred, so
\[
  C_n=C_n^{\nom}+S_n^{X_0}+\tilde E_n^\tau,\qquad \tilde E_n^\tau:=\sum_{n'\mid n}\mathcal T^{\nom}_{n'\to n}S^\tau_{n'},
\]
where $\mathcal T^{\nom}_{n'\to n}$ is the nominal transfer (the signed sum over nominal chains from level $n'$ to $n$).
On a step $d\to n$ the true and nominal filters at the centre $c$ are $\varphi_\tau(c)\varphi(c)$ and
$\kappa_\tau\kappa_{X_0}$, where $\varphi(c):=\mathbf 1[c\ge\Xm]$, $\kappa_{X_0}:=\mathbf 1[y/n\ge\Xm]$, and
$\varphi_\tau,\kappa_\tau$ are the $\tau$-filters. As in the proof of Proposition~\ref{prop:true-kernels}, on $X_0$-near steps the
whole difference goes to $S^{X_0}_n$, and on the other steps the difference $(\varphi_\tau-\kappa_\tau)\kappa_{X_0}$ goes to $S^\tau_n$;
so $S_n^{X_0}=-\sum_{d\to n\ X_0\text{-near}}\mathcal Q\bigl((\varphi_\tau\varphi-\kappa_\tau\kappa_{X_0})\,D_wC_d\bigr)$, where $\mathcal Q$
places the window $\bx_{H(c)}$ at each centre $c$. Split $C_d=C_d^{\nom}+(C_d-C_d^{\nom})$ and write
$S_n^{X_0}=S_n^{X_0,\nom}+S_n^{X_0,\tau}$ accordingly.

\emph{Step 2: the nominal part, level by level.} For a centre $c$ let $U_c$ be uniform on $[-H(c)/2,H(c)/2)$ and
$G(c):=\Prob(c+U_c\ge X_0)$, so that $\int\mathbf 1_{[X_0,\infty)}\,d(\mathcal Q\mu)=\int G\,d\mu$ for every measure $\mu$ of
centres. Near $X_0$, $G(c)=\min\bigl(1,\max\bigl(0,\frac12+(c-X_0)/H(c)\bigr)\bigr)$, which is Lipschitz with constant
$\le(1+2\vartheta_{X_0})/H_n\le\lambda_0/H_n$ on the support of $D_wD_d$ (width variation, \S\ref{sec:kernels-steps}), and
$G(y/n)=\Prob(y/n+U\ge X_0)$ with $U$ uniform on $[-H_n/2,H_n/2)$. By \eqref{eq:kernel-rec} for
the nominal kernels and $D_wC_d^{\nom}=c_d\delta_{y/n}+D_wD_d$,
\[
  C_n^{\nom}([X_0,\infty))=c_n\,G(y/n)-\sum_{d\mid n,\,d<n}\kappa_{d\to n}\int G\,d(D_wD_d).
\]
Adding
\[
  S_n^{X_0,\nom}([X_0,\infty))=-\sum_{X_0\text{-near}}\int G\,(\varphi_\tau\varphi-\kappa_\tau\kappa_{X_0})\,d(D_wC_d^{\nom}),
\]
and using $(\varphi_\tau\varphi-\kappa_\tau\kappa_{X_0})\,c_d\delta_{y/n}=0$ (the filters at the nominal point equal their nominal values),
\begin{equation}\label{eq:level-defect}
\begin{split}
  &C_n^{\nom}([X_0,\infty))+S_n^{X_0,\nom}([X_0,\infty))-c_n\mathbf 1[y/n\ge X_0]\\
  &\qquad=c_n\bigl(G(y/n)-\mathbf 1[y/n\ge X_0]\bigr)-\sum_{d\mid n,\,d<n}\int G\psi_{d\to n}\,d(D_wD_d),
\end{split}
\end{equation}
where $\psi_{d\to n}:=\varphi_\tau\varphi$ on $X_0$-near steps and $\psi_{d\to n}:=\kappa_\tau\kappa_{X_0}$ (a constant) on the
others. Since $D_1=0$, only $d\ge2$ occurs.
\begin{itemize}
\item $|G(y/n)-\mathbf 1[y/n\ge X_0]|\le\frac12$, and it vanishes unless $|y/n-X_0|\le H_n/2$, i.e. unless $n\in\mathbb B_0$.
\item $D_wD_d$ has mass $0$, total variation $\TV(D_d)$, and support of radius $R_d/w$ around $y/n$. The function
  $G\psi_{d\to n}$ takes values in $[0,1]$, so $|\int G\psi_{d\to n}\,d(D_wD_d)|\le\frac12\TV(D_d)$ (subtract the constant $\frac12$).
\item If the step is not $X_0$-near, then $G\psi_{d\to n}=\kappa_\tau\kappa_{X_0}G$, and $G$ is Lipschitz with constant
  $\le\lambda_0/H_n$; by Lemma~\ref{lem:kr-facts}(K2), $|\int G\,d(D_wD_d)|\le\lambda_0\KRn{D_wD_d}/H_n=\lambda_0\KRn{D_d}/(wH_n)$. If moreover
  $|y/n-X_0|>(1+\vartheta_{X_0})H_n/2+R_d/w$, the support lies where $G$ is constant (there $H(c)\le(1+\vartheta_{X_0})H_n$)
  and the term vanishes.
  % CORRECTIONS #29
\end{itemize}
Multiplying \eqref{eq:level-defect} by $1/n=1/(dw)$ and summing over $n$ gives at most $\frakB(y)$.

\emph{Step 3: the $\tau$-part.} What remains is
\[
  \sum_n n^{-1}\bigl(\tilde E_n^\tau+S_n^{X_0,\tau}\bigr)([X_0,\infty)),
\]
which is at most $\sum_n\TV(\tilde E_n^\tau)/n+\sum_n\TV(S_n^{X_0,\tau})/n$. As in the proof of Proposition~\ref{prop:true-kernels}
(the $\infty$-norm step), the first sum is $\le2.04\lambda_0\sum_d(\TV(C_d)/d)\,\theta(y/d)\le4.1\,C_NE_\theta/\log y$, by
Corollary~\ref{cor:kr-weights}(c) at every start, the bound on $W^{\mathrm{op,nt}}_y$ in Proposition~\ref{prop:true-kernels}
(the parents of $\tau$-near steps are non-terminal), and Lemma~\ref{lem:tau-near}. The
second is $\le\sum_d(\TV(\tilde E_d)/d)\,\eta_{X_0}(y/d)$ with $\eta_{X_0}$ as in the proof of
Proposition~\ref{prop:true-kernels}; with the block bound $\sum_{y/d\in[v,2v)}\TV(\tilde E_d)/d\le2\lambda_0C_N^2E_\eta\log(2v)/\log y$
from step (ii) of that proof and $\sum_v\log(2v)\sup\eta_{X_0}\le1.52\log(2X_0)$, it is
\[
  \le3.1\,C_N^2\tau_0E_\eta/\log y\le0.1\,C_NE_\theta/\log y,
\]
because $\tilde\eta\le\theta/\log^2\Xm$ pointwise. Together,
the $\tau$-part is at most $4.3\,C_NE_\theta/\log y$; if no step is $\tau$-near, then $S^\tau=0$, $\tilde E^\tau=0$ and
$C_d=C_d^{\nom}$ for every parent of an $X_0$-near step, so the $\tau$-part vanishes.
\end{proof}

\begin{remark}
% v3: remark after Lemma 3.14A
A cruder bound by the level sum $\sum_{n}(\TV(C_n)+|c_n|)/n$ over the band does not suffice: at $n_0=6$ it is about
$0.74$, while $\frakB(6X_0)\le5/12$ (Proposition~\ref{prop:xb-ranges}). The gain comes from three facts used above: the
direct chain ($d=1$) has no centre smear; the defect of a mass-$0$ measure against a $[0,1]$-valued function is at most
half its total variation; and the window average contributes only half a jump.
\end{remark}

\paragraph{Two counting tools.} The ranges below use two elementary bounds. The first is a local count on the realised
semigroup.

\begin{lemma}[Reciprocal sums over short intervals]\label{lem:short-recip}
% v3: Lemma 3.14A, middle range (Corollary 4.2, Lemma 4.18, Corollary 1.4); fix-Bprime §12.2 counting input
On $\Good(<y)$, for every $1\le M\le y/2$ and $0<a\le0.01$,
\[
  \sum_{e\in\NN\cap M[1-a,1+a]}\frac1e\le28\,a+7M^{-3/4}.
\]
\end{lemma}

\begin{proof}
Let $I:=M[1-a,1+a]\subset[M/2,2M]$; it suffices to show $\#(\NN\cap I)\le13.5|I|+K^*(M)$ with $K^*(M)/M\le6M^{-3/4}$,
since then $\sum_{e\in I}1/e\le(27aM+K^*(M))/(M(1-a))\le28a+7M^{-3/4}$.
\begin{itemize}
\item If $M<X_0/4$, then $I\subset[1,\Xm)$, where $\NN=\N$, and $\#(\NN\cap I)\le|I|+1$; take $K^*(M):=1$.
\item If $M\ge2X_0$, cut $I\subset[P,4P]$ ($P=M/2\ge X_0$) into $\lceil|I|/6H(P)\rceil$ pieces $J$ of length $\le6H(P)$.
  % Logged correction 37: Corollary localcount(ii) at the non-dyadic scale P = M/2; the range X_0 <= M < 2X_0 moves to the third case
  By Corollary~\ref{cor:localcount}(ii), $N(J)\le T(J)+11.4|J|+K_{\mathrm{loc}}(P)+K_{\mathrm{loc}}(2P)$, and under the target \eqref{eq:target},
  $T(J)\le2|J|+2+\#\Abs(J)$. Since $K_{\mathrm{loc}}\le H/100$, summing over the pieces gives
  $\#(\NN\cap I)\le13.42|I|+K_{\mathrm{loc}}(P)+K_{\mathrm{loc}}(2P)+2+\#\Abs(I)$. By Lemma~\ref{lem:lemmaP}, each window $[u,u+h(u)]$,
  $h(u)=u^{\theta/2}$, holds at most $h(u)L_u^{-4}$ absorbed elements; covering $I$ by $\lceil|I|/h(M/2)\rceil$ such windows,
  $\#\Abs(I)\le0.01|I|+(2M)^{\theta/2}$. As $K_{\mathrm{loc}}(P)+K_{\mathrm{loc}}(2P)+2\le M^{1/8}$ and $\theta<\frac12$, take $K^*(M):=2M^{1/4}$.
\item If $X_0/4\le M<2X_0$, apply the first case to $I\cap[1,\Xm)$ and the second, with $P:=X_0$, to
  $I\cap[X_0,\infty)\subset[X_0,4X_0]$; in between,
  $\NN\cap[\Xm,X_0)$ consists of integers and of at most $2H(X_0)$ jittered primes (\S\ref{sec:noise-notation}). This gives
  $K^*(M):=4X_0^{1/4}\le4(4M)^{1/4}$.
\end{itemize}
In all cases $K^*(M)/M\le6M^{-3/4}$.
\end{proof}

The second is a Rankin bound for smooth elements, valid at every height.

\begin{lemma}[Rankin bound]\label{lem:rankin-crude}
% v3: Lemma 3.14B, proof ("Rankin gives ..."); fix-Bprime §12.3 (B)
% Logged correction 73: on Good(<z') the hypothesis holds on [1,z] only for z below the horizon of z' (Remark rem:horizon).
Assume \eqref{eq:HBS} on $[1,z]$; for instance, work on $\Good(<z')$ for some $z'$ with $z\le\check z'$
(Remark~\ref{rem:horizon}; cf.\ Lemma~\ref{lem:cm}). Let $z\ge e^{10}$, $\ell\ge0$, and $v(\ell,z):=\sum_{n\in\NN\cap(e^\ell,2e^\ell],\,\Pplus(n)\le z}1/n$. Then $v(\ell,z)\le e^{14}\,(\log z)\,e^{-\ell/\log z}$.
\end{lemma}

\begin{proof}
Put $\beta:=1/\log z$. Since $\NN$ is the free semigroup on $\PP$,
$v(\ell,z)\le e^{-\beta\ell}\sum_{\Pplus(n)\le z}n^{\beta-1}=e^{-\beta\ell}\prod_{p\le z}(1-p^{\beta-1})^{-1}$. For $p\le z$,
$p^\beta\le e$, and $e^{\beta\log p}-1\le(e-1)\beta\log p$. So
\[
  \log\prod_{p\le z}(1-p^{\beta-1})^{-1}\le\sum_{p\le z}\frac1p+(e-1)\beta\sum_{p\le z}\frac{\log p}p
  +\sum_{p}\frac{p^{2\beta-2}}{1-p^{\beta-1}}.
\]
By Lemma~\ref{lem:cm}(iv) and (iii), the first two sums are $\le\log\log z+4.1$ and $\le(e-1)(1+0.43/\log z)\le1.8$; the
last is $\le1.3$: the primes $p\le7$ contribute $\le0.95$ (using $p^\beta\le p^{1/10}$), and those with $p\ge11$
contribute $\le1.33\,e^2\sum_{p\ge11}p^{-2}\le0.31$. Hence the product is $\le e^{7.2}\log z$, and the claim follows with
room.
\end{proof}

\paragraph{The three ranges.} Put
\begin{equation}\label{eq:jX0}
  j_{X_0}(x):=\sup_{y\in W_x}\log y\cdot\jz\Bigl(\frakB(y)+\frac{4.3\,C_NE_\theta}{\log y}\Bigr),
\end{equation}
so that $\log y\cdot|J_{X_0}(y)|\le j_{X_0}(x)$ for every $y\in W_x$, by Lemma~\ref{lem:xb}.

\begin{proposition}[The $X_0$ band in three ranges]\label{prop:xb-ranges}
% v3: Lemma 3.14A, "The three ranges"
Let $\Lambda\ge10^8$, $x\ge2X_0$ and $N:=x/X_0$. On $\Good(<x)$:
\begin{enumerate}[label=(\roman*)]
\item (small) if $N\le e^{0.01\Lambda}$, then $j_{X_0}(x)\le0.4209$ (computer-assisted; see Appendix~\ref{app:cert-x0band});
\item (middle) if $e^{0.01\Lambda}\le N\le e^{600\Lambda^4}$, then $\frakB(y)\le200\delta+400N^{-1/10}$ for $y\in W_x$, and
  $j_{X_0}(x)\le e^{-0.0009\Lambda}$;
\item (huge) if $N\ge e^{600\Lambda^4}$, then $\frakB(y)\le(400\delta+e^{-60})\,\tau_0/\log N$ for $y\in W_x$, and
  $j_{X_0}(x)\le e^{-59}$.
\end{enumerate}
The three ranges cover $x\ge2X_0$ for every $\Lambda\ge10^8$.
\end{proposition}

Throughout the proof, $y\in W_x$, and we write $N$ also for $y/X_0$, which differs from $x/X_0$ by a factor
$1+O(H(x)/x)$. We use the following consequences of Theorem~\ref{thm:kr-random} for the nominal kernels (with
$H_e/H_d\le(d/e)^\eps$ from \eqref{eq:RH}, and $|c_e|\le\mu^2(e)$):
\begin{equation}\label{eq:TVD-bounds}
\begin{gathered}
  \tfrac12\TV(D_d)\le|c_d|+\tfrac{\lambda_0}4\sum_{m\mid d,\,1<m<d}\Hord(m)m^{\eps-1}|c_{d/m}|,\\
  \frac{\lambda_0\KRn{D_d}}{wH_n}\le\frac{\lambda_0^2w^{\eps-1}}4\sum_{m\mid d,\,m<d}\Hord(m)m^{\eps-1}|c_{d/m}|\quad(n=dw).
\end{gathered}
\end{equation}
For the first, $\TV(D_d)\le|c_d|+\TV(C_d^{\nom})$; for the second, $H_e/H_n\le(n/e)^\eps=(wm)^\eps$ with $m=d/e$. Since
$\lambda_0^2\le1+7\delta$, the factors $\lambda_0$, $\lambda_0^2$ change the numerical constants below by a relative amount
$\le7\delta\le e^{-9\cdot10^7}$, and we absorb them. Consequently, for every set $I$ of levels,
\begin{equation}\label{eq:TVD-sums}
\begin{gathered}
  \sum_{d\in I}\frac{\TV(D_d)}{2d}\le\sum_{d\in I}\frac{|c_d|}d+\frac{\lambda_0}4\sum_{m\ge2}\Hord(m)m^{\eps-2}\sum_{e\in I/m}\frac{|c_e|}e,\\
  \sum_{d\in I}\frac{\lambda_0\KRn{D_d}}{d\,wH_{dw}}\le\frac{\lambda_0^2w^{\eps-1}}4\sum_{m\ge1}\Hord(m)m^{\eps-2}\sum_{e\in I/m}\frac{|c_e|}e .
\end{gathered}
\end{equation}

\begin{proof}[Proof of (i): the small range]
\emph{A single level.} Let $N\le e^{0.01\Lambda}$. If $n\in\mathbb B_0$ then $|n-N|\le nH_n/(2X_0)\le0.51\,(1.01N)\delta$.
If $(d,w)\in\mathbb B_1\cup\mathbb B_2$ then $|y/n-X_0|\le H_n/2+R_d/w+(X_0-\Xm)\le1.6H(X_0)$, so $|n-N|\le2N\delta$. Since
$N\delta\le e^{0.01\Lambda+8(\log\Lambda)^2-\Lambda}<10^{-10}$, every level occurring in $\frakB(y)$ is the integer $n_0$
nearest to $N$, and $n_0<X_0$. Below $\Xm$ the semigroup is $\N$, so all divisors $d$ of $n_0$ are rational integers, all
cofactors are rigid, and no step below level $n_0$ is $\tau$-near or $X_0$-near; in particular the $\tau$-term of
\eqref{eq:xb} vanishes.

\emph{The bound $B(n_0)$.} For $2\le n$, put
\begin{gather*}
  B(n):=\frac1n\Bigl[\tfrac12\mu^2(n)+\tfrac12\sum_{d\mid n,\,1<d<n}\Bigl(2\mu^2(d)+\tfrac12\lambda\frac{X(d)}d\Bigr)\Bigr],\\
  X(d):=\sum_{e\mid d,\,1<e<d}\mu^2(e)\,e\,\Hord(d/e),\qquad\lambda:=1.0001 .
\end{gather*}
Each pair $(d,n_0/d)$ lies in at most one of $\mathbb B_1,\mathbb B_2$. By \eqref{eq:TVD-bounds} with $\eps_*$ in place of
$\eps$ and $\lambda_0(d/e)^{\eps_*}\le\lambda$ for $d\le10^6$, the $\mathbb B_1$-weight of $d$ is $\le\mu^2(d)+\frac14\lambda X(d)/d$;
the $\mathbb B_2$-weight is $\le\frac18\lambda(\mu^2(d)+X(d)/d)$, which is smaller. Hence $\frakB(y)\le B(n_0)$ for
$n_0\le10^6$.

\emph{Computation, $n\le10^6$.} The certificate of Appendix~\ref{app:cert-x0band} computes $\mu^2$, $\Hord$, $X$ and the divisor
sums exactly in integers, and evaluates $B(n)$ as a sum of at most $\tau(n)$ nonnegative terms, inflated by $1+10^{-9}$.
The maximum over $2\le n\le10^6$ is $B(6)=(\frac12+\frac12(2+2))/6=5/12$; the next values are $B(12)=0.2778$ and
$B(2)=B(4)=B(10)=\frac14$, and $\sup_{n\ge100}B(n)=0.0881$.

\emph{Analytic tail, $10^6<n\le X_0^{0.01}$.} Here we use $H_e/H_d\le(d/e)^{\eps_*}$ directly instead of $\lambda$ (all heights are $\le X_0^{1.01}$). Write
$d=em$ with $e>1$ squarefree, $m\ge2$. Then $\sum_{d\mid n}\frac1n\sum_{e\mid d}\mu^2(e)\Hord(m)m^{\eps_*-1}$ is a sum over
factorisations $n=mk$ ($k\ge2$) of $\Hord(m)m^{\eps_*-2}\cdot2^{\omega(k)}k^{-1}$, which for $\sigma:=0.2$ is at most
\[
  n^{-\sigma}\sum_{m\ge2}\Hord(m)m^{\sigma+\eps_*-2}\cdot\sup_{k\ge2}2^{\omega(k)}k^{\sigma-1}
  =n^{-0.2}\,A(1.8-\eps_*)\cdot2^{0.2},
\]
where $A(s):=1/(2-\zeta(s))-1$ and the supremum is attained at $k=2$ (since $2p^{-0.8}<1$ for $p\ge3$). In ball arithmetic,
$A(1.8-3\cdot10^{-6})\le7.4915$ (the value is $7.49142$; Appendix~\ref{app:cert-consts}). With $2^{\omega(n)}\le1.64\sqrt n$ (the maximum of $\prod_{p\mid n}2p^{-1/2}$ is at $n=6$),
\[
  B(n)\le\frac{\frac12+1.64\sqrt n}{n}+\frac14\cdot7.4915\cdot2^{0.2}\,n^{-0.2}\le\frac{\frac12+1.64\sqrt n}{n}+2.152\,n^{-0.2}
  \le0.14\qquad(n>10^6).
\]

\emph{Conclusion.} $\log y\le\Lambda+\log N+10^{-10}$, so $\log y\cdot\jz\le1.0001\,(1.01+10^{-18})\le1.010102$, and
$4.3C_NE_\theta\le e^{-0.99\Lambda}$ by \eqref{eq:tau-term-bound}. So $j_{X_0}(x)\le1.010102\cdot\frac5{12}+e^{-0.99\Lambda}\le0.42088\le0.4209$.
\end{proof}

\begin{proof}[Proof of (ii): the middle range]
Let $e^{0.01\Lambda}\le N\le e^{600\Lambda^4}$. Every set below is an interval of relative half-width $\le0.01$, so
Lemma~\ref{lem:short-recip} applies to it, and we use $|c_e|\le1$.

\emph{($\alpha$) The levels $\mathbb B_0$.} Since $H_n\le H(1.01X_0)\le1.01H(X_0)$ there, $\mathbb B_0\subset
N[1-0.51\delta,1+0.51\delta]$, so this part of $\frakB(y)$ is $\le\frac12(28\cdot0.51\delta+7N^{-3/4})\le7.2\delta+3.5N^{-3/4}$.

\emph{($\beta_1$) $X_0$-near pairs with $w\le\sqrt N$.} For fixed $w$ the condition $|y/d-w\Xm|\le R_d\le1.02H(y/d)$ and
\eqref{eq:RH} give $d\in I_w:=M_w[1-1.01\delta_w,1+1.01\delta_w]$, with $M_w:=y/(w\Xm)$ and $\delta_w:=1.03w^{\eps-1}\delta$.
By \eqref{eq:TVD-sums} and Lemma~\ref{lem:short-recip},
\[
  \sum_{d\in I_w}\frac{\TV(D_d)}{2d}\le28.3\delta_w\Bigl(1+\frac{A_0}4\Bigr)+7M_w^{-3/4}
  +\frac74M_w^{-3/4}\sum_{m\le2M_w}\Hord(m)m^{\eps-1.25}\le43\,\delta_w+41\,M_w^{-0.2},
\]
using $A_0=A_0(\eps)\le2.07$ and $\sum_{m\le X}\Hord(m)m^{\eps-1.25}\le X^{0.55}A(1.8-\eps)\le13.1X^{0.55}$, where
$A(1.77)=13.08$ (the abscissa of $\sum\Hord(m)m^{-s}$ is $1.7286<1.77$). Weighting by $1/w$ and summing over
$2\le w\le\sqrt N$, with $M_w\ge N/w$ and $\sum_{w\ge2}w^{\eps-2}\le\zeta(1.97)-1=0.674$, gives
$\le43\cdot1.03\cdot0.674\,\delta+41N^{-0.2}\cdot5N^{0.1}\le30\,\delta+205\,N^{-0.1}$.

\emph{($\beta_2$) Pairs in $\mathbb B_2$ with $w\le\sqrt N$.} Here $|y/d-wX_0|\le wH_n/2+R_d\le1.04\,w H(X_0)$, so $d$
ranges over an interval of relative half-width $\le1.05\delta$ around $y/(wX_0)$. By \eqref{eq:TVD-sums} and
Lemma~\ref{lem:short-recip}, for fixed $w$ this part is
$\le\frac14w^{\eps-1}\bigl(29.4\delta(1+A_0)+142(w/N)^{0.2}\bigr)$, and summing over $w$ gives
$\le\frac14\cdot29.4\cdot3.07\cdot0.674\,\delta+36N^{-0.2}\sum_{w\ge2}w^{-1.77}\le15.3\,\delta+34\,N^{-0.2}$.

\emph{($\gamma$) Pairs with $w>\sqrt N$.} Then $d<1.01\sqrt N$; fix $d$ and count $w$. For $X_0$-near pairs the
$w$-window $\{w:|y/d-w\Xm|\le R_d\}$ has relative half-width $\le1.02\delta(d/N)^{1-\eps}$ at height
$M':=y/(d\Xm)\ge0.98\sqrt N$, so $\sum_w1/w\le29\delta(d/N)^{1-\eps}+7M'^{-3/4}$. For pairs in $\mathbb B_2$ the window has
relative half-width $\le1.05\delta$, and by \eqref{eq:TVD-bounds} the weight carries $\frac14w^{\eps-1}\le\frac14N^{-0.485}$.
By Theorem~\ref{thm:kr-random}, $\sum_{d\le D}\TV(D_d)/(2d)\le1.52(\log D+1)$ and
$\sum_{d\le D}d^{-1}\sum_{m\mid d}\Hord(m)m^{\eps-1}|c_{d/m}|\le3.07(\log D+1)$. Hence this part is
$\le(\log N+1)\bigl(1.52(29\delta N^{-0.485}+7.2N^{-3/8})+0.77N^{-0.485}(30\delta+7.2N^{-3/8})\bigr)\le N^{-1/4}$.

\emph{Total.} $\frakB(y)\le53\,\delta+3.5N^{-3/4}+205N^{-0.1}+34N^{-0.2}+N^{-1/4}\le200\,\delta+400\,N^{-1/10}$. Finally
$\log y\le601\Lambda^4$, so $\log y\cdot\jz\le1.0001\cdot601\Lambda^3$, and with $N\ge e^{0.01\Lambda}$ and
$4.3C_NE_\theta\le e^{-0.99\Lambda}$ by \eqref{eq:tau-term-bound},
\[
  j_{X_0}(x)\le1.0001\cdot601\Lambda^3\bigl(200\,e^{8(\log\Lambda)^2-\Lambda}+400\,e^{-0.001\Lambda}\bigr)+e^{-0.99\Lambda}\le e^{-0.0009\Lambda},
\]
because $0.0001\Lambda\ge12.4+3\log\Lambda$ for $\Lambda\ge10^8$.
\end{proof}

\begin{proof}[Proof of (iii): the huge range]
Let $J:=\log N\ge600\Lambda^4$ and $W_1:=e^{100}$. The main pieces use Lemma~\ref{lem:si} below, with $J_1=300\Lambda^4$.
Recall from Lemma~\ref{lem:lemmaM} that $c_e(y)\ne0$ forces $e$ to be squarefree and $4v$-dense whenever $y/e\in[v,2v)$,
including the condition $p_1\le4v$ on its smallest prime.
% Logged correction 40: t-dense includes j = 0 (p_1 <= t)
Applying this with $v$ just above $y/(2e)$, so that $4v\le2.03\,sX_0$, an element $e$ with $y/e\le1.01\,sX_0$ is $t$-dense
with $t=4.01\,sX_0$.

\emph{($\alpha$) The levels $\mathbb B_0$.} $\mathbb B_0\subset[N(1-0.51\delta),N]\cup[N,N(1+0.51\delta)]$. Lemma~\ref{lem:si}
with $t=4.01X_0$ and $\eta=0.52\delta\in[e^{-2\Lambda},\frac12]$ gives
$\frac12\cdot2\cdot61.1(1+2.1\tau_0/J)\cdot0.52\delta\,\tau_0/(J-1)\le32\,\delta\tau_0/J$.

\emph{($\beta_1$) $X_0$-near pairs with $w\le W_1$.} As in (ii), $d\in I_w=M_w[1-1.01\delta_w,1+1.01\delta_w]$ with
$\delta_w=1.03w^{\eps-1}\delta\ge\delta e^{-100}\ge e^{-2\Lambda}$ and $\log M_w\ge J-101$. In \eqref{eq:TVD-sums}, the
element $d$ is $t_w$-dense with $t_w:=4.01wX_0$, and the element $e\in I_w/m$ is $t_{wm}$-dense with $t_{wm}:=4.01wmX_0$.
\begin{itemize}
\item For $m\le N^{1/100}/w$: the height of $I_w/m$ is $\ge N^{0.99}/1.01$, so its logarithm $J'$ is $\ge0.99J-1\ge J_1$;
  and $t_{wm}\le4.01N^{1/100}X_0\le(N^{0.99})^{1/10}$ because $X_0\le N^{1/(600\Lambda^3)}$. So Lemma~\ref{lem:si} applies, with
  $\tau_{wm}/J'\le0.0103$ and factor $1+2.1\tau_{wm}/J'\le1.022$:
  $\sum_{e\in I_w/m}|c_e|/e\le2\cdot61.1\cdot1.022\cdot1.02\,\delta_w\,\tau_{wm}/J'$, with $\tau_{wm}=\tau_0+\log w+\log m$.
\item For $m>N^{1/100}/w$: Lemma~\ref{lem:short-recip} gives $\sum_{e\in I_w/m}1/e\le29\delta_w+7$, and
  \[
    \sum_{m>N^{1/100}/w}\Hord(m)m^{\eps-2}\le(N^{1/100}/W_1)^{-0.2}A(1.77)\le e^{22.6}N^{-0.002}.
  \]
\end{itemize}
With $A_0\le2.07$ and $A_1=A_1(\eps)\le9.42$, and $J'\ge0.99J-1$,
\[
  \sum_{d\in I_w}\frac{\TV(D_d)}{2d}\le\frac{127.4\,\delta_w}{0.99J-1}\Bigl(\tau_w\bigl(1+\tfrac{A_0}4\bigr)+\tfrac{A_1}4\Bigr)
  +e^{26}N^{-0.002},\qquad \tau_w:=\tau_0+\log w .
\]
Weighting by $1/w$ and summing over $2\le w\le W_1$, with
\[
  \sum_{w\ge2}w^{-1.97}=0.674\quad\text{and}\quad\sum_{w\ge2}w^{-1.97}\log w=-\zeta'(1.97)\le1.0001
\]
(ball arithmetic; Appendix~\ref{app:cert-consts}), gives
\[
  \le\frac{132.5\,\delta}{0.99J-1}\bigl(0.674(1.518\tau_0+2.36)+1.518\cdot1.0001\bigr)+e^{31}N^{-0.002}\le138\,\delta\tau_0/J.
\]

\emph{($\beta_2$) Pairs in $\mathbb B_2$ with $w\le W_1$.} Now $d$ ranges over an interval of relative half-width
$\le1.05\delta$ around $y/(wX_0)$, and by \eqref{eq:TVD-sums} with the factor $\frac14w^{\eps-1}$ and the same use of
Lemma~\ref{lem:si} (now including $m=1$),
this part is $\le\frac14\sum_{2\le w\le W_1}w^{\eps-2}\cdot\frac{2\cdot61.1\cdot1.022\cdot1.06\,\delta}{0.99J-1}
\bigl(\tau_w(1+A_0)+A_1\bigr)+e^{31}N^{-0.002}\le72\,\delta\tau_0/J$.

\emph{($\gamma'$) Pairs with $w>W_1$.} Fix $d$; then $y/d>W_1X_0/2$. Group $d$ by the dyadic block $y/d\in[v,2v)$,
$v\ge W_1X_0/2$. For $X_0$-near pairs, the $w$-window has relative half-width $\le1.02H(v)/v\le1.02\delta(X_0/v)^{1-\eps}$
at height $\ge v/X_0$, so $\sum_w1/w\le29\delta(X_0/v)^{0.97}+7(X_0/v)^{3/4}$ by Lemma~\ref{lem:short-recip}. For pairs in
$\mathbb B_2$ the window has relative half-width $\le1.05\delta$, and the weight carries
$\frac14w^{\eps-1}\le\frac14(X_0/v)^{0.97}$. By Proposition~\ref{prop:true-kernels} and Corollary~\ref{cor:saias} (and
Corollary~\ref{cor:kr-weights} for the $m$-sums), the blocks satisfy
\begin{gather*}
  \sum_{y/d\in[v,2v)}\TV(D_d)/(2d)\le(C_N+C_S)\log(2v)/\log y\\
  \text{and}\quad\sum_{y/d\in[v,2v)}d^{-1}\sum_{m\mid d}\Hord(m)m^{\eps-1}|c_{d/m}|\le C_S\bigl((1+A_0)\log(2v)+A_1\bigr)/\log y.
\end{gather*}
Summing
the geometric series over $v=2^jW_1X_0/2$,
\begin{multline*}
  (\gamma')\le\frac{2.5\,(C_N+C_S)(\tau_0+102)}{\log y}\Bigl(29\delta\bigl(\tfrac2{W_1}\bigr)^{0.97}+7\bigl(\tfrac2{W_1}\bigr)^{3/4}\Bigr)\\
  +\frac{2.5\cdot0.77\,C_S(\tau_0+102)}{\log y}\bigl(\tfrac2{W_1}\bigr)^{0.97}(30\delta+7)\le e^{-60}\,\frac{\tau_0}J .
\end{multline*}

\emph{Total.} $\frakB(y)\le(32+138+72)\,\delta\tau_0/J+e^{-60}\tau_0/J\le(400\delta+e^{-60})\,\tau_0/\log N$. Since
$\log y\le J+\Lambda+1$ and $\tau_0\le1.0001\Lambda$, $\log y\cdot\jz\cdot\tau_0/J\le1.0003$, so with
$4.3C_NE_\theta\le e^{-0.99\Lambda}$ by \eqref{eq:tau-term-bound}, $j_{X_0}(x)\le1.0003(400\delta+e^{-60})+e^{-0.99\Lambda}\le e^{-59}$.
\end{proof}

\paragraph{Contiguity.} The three ranges meet at $N=e^{0.01\Lambda}$ and $N=e^{600\Lambda^4}$ and so cover $x\ge2X_0$ for
every $\Lambda\ge10^8$. No step uses a short-interval count at relative width below $e^{-2\Lambda}$, and no sum of the form
$\sum_w1/(w\log(N/w))$ arises: the pairs with $w>W_1$ are counted by fixing $d$.

\paragraph{A short-interval Saias bound.} The huge range rests on the following lemma. Saias' asymptotic for integers with
$t$-dense divisors has relative error $O(1/\log t)$, which is far larger than the relative widths $\eta\approx\delta$ needed
here, so it cannot be used directly. Instead we combine the joint ``$t$-dense and smooth'' profile of
Lemma~\ref{lem:bs-profile} with the short-window prime counts of Theorem~\ref{thm:lp}\eqref{eq:lpplus}.

\begin{lemma}[Short-interval Saias bound]\label{lem:si}
% v3: Lemma 3.14B (Lemma SI(Λ))
Put $\Lambda_1:=100\Lambda^3$ and $J_1:=300\Lambda^4$. Let $J:=\log N\ge J_1$, $t\in[X_0,N^{1/10}]$, $\tau:=\log t$ and
$\eta\in[e^{-2\Lambda},\frac12]$. Then, on $\Good(<y)$ for every $y\ge N$,
\begin{equation}\label{eq:si}
  S_t\bigl([N,N(1+\eta)]\bigr):=\sum_{\substack{n\in\NN\cap[N,N(1+\eta)]\\ n\ \text{squarefree, $t$-dense}}}\frac1n
  \le61.1\Bigl(1+\frac{2.1\,\tau}J\Bigr)\frac{\eta\,\tau}J .
\end{equation}
\end{lemma}

The trivial bound is about $27\eta$; \eqref{eq:si} is the Saias density $\tau/J$ inside an interval of relative width
$\eta$, with no loss in $\eta$.

\begin{proof}
Let $I:=[N,N(1+\eta)]$ and let $n\in I$ be squarefree and $t$-dense. Since $n>t$, $n$ is not prime. Put $p:=\Pplus(n)$ and
$m:=n/p$. Then $m$ is squarefree and $t$-dense with $\Pplus(m)<p\le tm$ (Theorem~\ref{thm:bs}, Step 1 of its proof); in
particular $p^2\le t\,mp=tn$, so $m=n/p\ge\sqrt{n/t}\ge\sqrt{N/t}$ (this uses the lower end $n\ge N$ of $I$). Put
$P_1:=e^{\Lambda_1}$ and split according to $p<P_1$ or $p\ge P_1$.

\emph{(B) $p<P_1$.} Then $n$ is $P_1$-smooth, and $I\subset(N/2,N]\cup(N,2N]$. By Lemma~\ref{lem:rankin-crude}, this part is
$\le2e^{14}\Lambda_1e^{-(J-1)/\Lambda_1}$. This is $\le10^{-6}\eta\tau/J$, because $\eta\ge e^{-2\Lambda}$, $\tau\ge\Lambda$, and
$(J-1)/\Lambda_1\ge3\Lambda-10^{-20}\ge2\Lambda+\log J+3\log\Lambda+34$ at $J=J_1$, while the left side grows linearly and the
right logarithmically in $J$.

\emph{(A) $p\ge P_1$: the prime sum.} For fixed $m$, $p$ runs over primes in $I/m=[X,X(1+\eta)]$ with $X:=N/m$ and
$\Lambda_p:=\log X\ge\Lambda_1-1$; in particular $X\ge X_{\mathrm{LP}}$, since $\Lambda_1-1=100\Lambda^3-1\ge10^5\Lambda^2\log^2\Lambda=\log X_{\mathrm{LP}}$.
% Logged correction 70 (CORRECTIONS-vs-v3.md): X >= X_LP made explicit (was "log X_LP = O(Lambda^2 log^2 Lambda)").
% Logged correction 45: new L_1 (log k + 3.08), the extra term Lambda_p^{-2} in Rel, the hypothesis v(1+2 eta) <= y, and the
% explicit window of Theorem lp(iii)
% Logged correction 73 (CORRECTIONS-vs-v3.md): v3 applies Theorem lp with its parameter y equal to the cutoff of Good(<y),
% where (LP-a) is not certified up to y; it is now applied with y' = 2N^{0.55}, far below the horizon.
Since $m\ge\sqrt{N/t}\ge N^{0.45}$ and $\eta\le\frac12$, also $X(1+2\eta)\le2N^{0.55}=:y'$. We apply Theorem~\ref{thm:lp}(iii)
with $y'$ in place of its parameter $y$ and $v=X$, so that its hypothesis $v(1+2\eta)\le y'$ holds; (LP-a) holds up to $y'$ on
$\Good(<y)$ by Remark~\ref{rem:horizon}, as $y'\le0.9N\le0.9y$. Theorem~\ref{thm:lp}(iii) is stated for $\Lambda_p\ge76\Lambda^3$; here
$\Lambda_p\ge\Lambda_1-1\ge76\Lambda^3$. There $\mathrm{Rel}\le0.05+\Lambda_p^{-2}\le0.1$ as soon as the exponent in
\eqref{eq:lpplus} is at least $\log(360\,G(L_1)\Lambda_p^3\eta^{-1})$. Put $\ell_p:=\log\Lambda_p$, so that
$k\le\log(1/\eta)+20\ell_p+1\le2\Lambda+20\ell_p+1$ and
$L_1=2\ell_p+\log(1/\eta)+\log k+3.08\le2\Lambda+2\ell_p+\log(2\Lambda+20\ell_p+1)+3.08$, and write
$D:=L_1+442+2\log\Lambda+\log L_1$ for the denominator of the exponent. Since $c_{\mathrm{zf}}\ge0.088/\Lambda$, the exponent
is at least $0.0792\,\Lambda_p/(\Lambda D)$.
\begin{itemize}
\item \emph{Case $\ell_p\le e^{\Lambda/2}/3$.} Then $L_1\le e^{\Lambda/2}$, so $\log G(L_1)\le\Lambda$ by Theorem~\ref{thm:lp}(iii).
  Write $\Lambda_p=76\Lambda^3r$ with $r\ge1$, and $\ell_0:=\log(76\Lambda^3)\le0.001\Lambda$. Then $\ell_p=\ell_0+\log r$; since
  $k$ and $L_1$ are at most $3\Lambda(1+\log r)$ and $1+\log r\le r$, both $\log k$ and $\log L_1$ are at most
  $\log(3\Lambda)+\log r$. So
  \[
    D\le2.003\Lambda+4\log r,\qquad \log\bigl(360\,G(L_1)\Lambda_p^3\eta^{-1}\bigr)\le3\Lambda+3\ell_p+5.9\le3.002\Lambda+3\log r,
  \]
  since $2\ell_0+446+2\log\Lambda+2\log(3\Lambda)\le0.003\Lambda$ and $3\ell_0+5.9\le0.002\Lambda$. Using $\log r\le r$ and
  $\log^2r\le r$,
  \[
  \begin{aligned}
    \Lambda D\,(3.002\Lambda+3\log r)&\le\Lambda\bigl(6.0131\Lambda^2+18.02\Lambda\log r+12\log^2r\bigr)\le6.0132\,\Lambda^3r\\
    &<0.0792\cdot76\,\Lambda^3r=0.0792\,\Lambda_p ,
  \end{aligned}
  \]
  which is the required inequality.
\item \emph{Case $\ell_p>e^{\Lambda/2}/3$.} Then $\ell_p\ge\Lambda^2$, so $k\le21\ell_p$ and
  $L_1\le2\ell_p+2\Lambda+\log(21\ell_p)+3.08$. Since $2\Lambda+457+2\log\Lambda+2\log(21\ell_p)\le\ell_p$ (with huge room, as
  $\ell_p\ge\Lambda^2\ge10^{16}$), both $L_1+450+2\log\Lambda\le3\ell_p$ and $D=L_1+442+2\log\Lambda+\log L_1\le3\ell_p$. So
  $G(L_1)\le400\cdot3\ell_p\le\Lambda_p$, hence $\log(360\,G\Lambda_p^3\eta^{-1})\le4\ell_p+2\Lambda+5.9\le5\ell_p$, and the
  required inequality $0.0792\,\Lambda_p\ge\Lambda\cdot3\ell_p\cdot5\ell_p$ holds because
  $190\Lambda\ell_p^2\le\ell_p^3\le e^{\ell_p}=\Lambda_p$.
\end{itemize}
So $\mathrm{Rel}\le0.1$ for every $\Lambda_p\ge\Lambda_1-1$. The window $\varphi$ of Theorem~\ref{thm:lp}(iii) majorises
$\mathbf 1_{[X,X(1+\eta)]}$, takes values in $[0,1]$ and is supported in $[X(1-w),X(1+\eta+w)]$ with $w=\eta\Lambda_p^{-2}$. As
$d\Pi_0\le du/\log u$, this gives
$\Pi([X,X(1+\eta)])\le\int\varphi\,d\Pi\le1.1\int\varphi\,d\Pi_0\le1.1\,(1+3\Lambda_p^{-2})\,\eta X/\log X$: the support has
length $\eta(1+2\Lambda_p^{-2})X$, and $\log u\ge\log X-2w$ on it. As
$\Lambda_p\ge10^{25}$, $1.1(1+3\Lambda_p^{-2})\le1.1001$. Hence
\[
  \sum_{p\in I/m}\frac1p\le\frac1X\,\Pi\bigl([X,X(1+\eta)]\bigr)\le1.1001\,\frac{\eta}{\log X}=\frac{1.1001\,\eta}{\log(N/m)} .
\]

\emph{(A) $p\ge P_1$: the $m$-sum.} Group $m$ into dyadic blocks $(M_k,2M_k]$, $M_k:=2^kM_0$ with $M_0:=\frac12\sqrt{N/t}$,
up to $M_k\le N(1+\eta)/P_1$. In a block, $\Pplus(m)<p\le N(1+\eta)/m<2N/M_k$ and $\log(N/m)\ge\sigma_k:=\log(N/2M_k)$, so
\[
  S_A\le1.1001\,\eta\sum_k\frac{W_t(M_k;\,2N/M_k)}{\sigma_k},\qquad
  W_t(M;z):=\sum_{\substack{m\in\NN\cap(M,2M]\\ m\ \text{sqf, $t$-dense},\ \Pplus(m)<z}}\frac1m .
\]
Put $\ell_k:=\log M_k$ and $u_k:=\ell_k/\log(2N/M_k)$. Since $\tau\le J/10$, $\ell_k\ge(J-\tau)/2-2\log2\ge\tau$. Recall from
Lemma~\ref{lem:bs-profile} the profile $U(\ell)=\ell^{1/4}\bigl(1+\sqrt{2\log(\ell/\tau)+4}\bigr)$,
$Q(\ell,u)=A+a$ with $A=11(1+u)^3\log(2+u)/\ell$ and $a=(u-\ell^{1/4})_+^2/(2\sqrt\ell)$, and $\Phi_\infty\le38.3$.

\emph{Region 1: $u_k\le U(\ell_k)$.} By Lemma~\ref{lem:bs-profile},
$W_t(M_k;2N/M_k)\le(\tau/\ell_k)\,\Phi_\infty\,\rho(u_k-1)\,e^{Q(\ell_k,u_k)}$, with $\rho(v):=1$ for $v\le1$.
\begin{itemize}
\item For $u\le\ell^{1/4}$: $a=0$ and $A\le11\cdot1.01\,\ell^{-1/4}(\frac14\log\ell+0.01)\le7\cdot10^{-7}$, because
  $\ell\ge0.449J\ge10^{34}$ (as $J\ge300\Lambda^4\ge3\cdot10^{34}$) and $\ell^{-1/4}\log\ell$ is decreasing. For $\ell^{1/4}<u\le U(\ell)$: $e^{Q}\le e^{2}\ell/\tau$ and
  $\rho(u-1)\le1/\Gamma(\ell^{1/4})$, so these $k$ contribute less than $10^{-6}\tau/J$ in total.
\item With $\tilde J:=J+\log2$ we have $u_k=\ell_k/(\tilde J-\ell_k)$ and $\sigma_k=\tilde J-\ell_k-2\log2$. The function
  $f(\ell):=\rho(u(\ell)-1)/(\ell(\tilde J-\ell))$ satisfies
  \[
    \int f(\ell)\,d\ell=\frac1{\tilde J}\int\rho(u-1)\frac{du}u\qquad(\ell=\tilde Ju/(1+u)).
  \]
  The sum over $\ell_k\in\ell_0+(\log2)\Z$ differs from $(\log2)^{-1}\int f$ by at most the total variation of $f$ on
  Region~1, which is $\le4\sup f\le10^{-6}/J$, since there $\tilde J-\ell\ge\ell/U(\ell)$.
\item The lower limit is $u_{\min}\ge(J-\tau-6)/(J+\tau+8)$, and by $\int_a^\infty\rho(w-1)\,dw/w=\rho(a)$ ($a\ge1$),
  \[
    \int_{u_{\min}}^\infty\rho(u-1)\frac{du}u=\log\frac1{u_{\min}}+\rho(1)\le1+\log\frac{J+\tau+8}{J-\tau-6}\le(1+10^{-6})\Bigl(1+\frac{2.1\tau}J\Bigr),
  \]
  using $\log\frac{1+x}{1-x}=2x(1+\frac{x^2}3+\frac{x^4}5+\cdots)\le2.1x$ for $0\le x\le0.1$.
\end{itemize}
Hence Region~1 contributes at most
$1.1001\cdot\frac{38.3}{\log2}\,(1+3\cdot10^{-6})\bigl(1+\frac{2.1\tau}J\bigr)\frac{\eta\tau}J
\le60.81\bigl(1+\frac{2.1\tau}J\bigr)\frac{\eta\tau}J$.

\emph{Region 2: $u_k>U(\ell_k)$.} Drop the density condition: by Lemma~\ref{lem:rankin-crude},
$W_t(M_k;2N/M_k)\le e^{14}\log(2N/M_k)\,e^{-u_k}$, and $\log(2N/M_k)\le1.01\sigma_k$. Since $U(\ell)\ge3\ell^{1/4}\ge3(0.449J)^{1/4}$
and there are at most $J/\log2+1$ blocks, Region~2 contributes $\le1.1001\,\eta\,e^{14}\cdot1.01\,(J/\log2+1)\,e^{-3(0.449J)^{1/4}}
\le10^{-6}\eta\tau/J$.

Adding (B), Region~1 and Region~2 gives \eqref{eq:si}, since $60.81+2\cdot10^{-6}\le61.1$.
\end{proof}

\begin{remark}
% v3: §3.6, after Lemma 3.14B
No cancellation of the M\"obius function and no short-interval asymptotic for smooth numbers is used. The smooth part
needs no density saving: its Rankin factor $e^{-J/\Lambda_1}$ beats $\eta/J$ outright. The thresholds $\Lambda_1$ and $J_1$
depend only polynomially on $\Lambda$; if the zero-free constant of Theorem~\ref{thm:lp} were only $c_{\mathrm{zf}}\ge\Lambda^{-A}$,
it would suffice to take $\Lambda_1=100\Lambda^{2+A}$ and $J_1=300\Lambda^{3+A}$, and the middle range of
Proposition~\ref{prop:xb-ranges} would still close.
\end{remark}


\paragraph{Oscillation of the target.}\label{par:osc} All bounds are for $\Tc$, and none uses a local count at the height where
the oscillation is read.
% v3: (Osc)
\begin{equation}\label{eq:osc}
  \sum_u\log(2u)\,o_u\le e^{-83},\qquad\sum_v(\log3v)^2\,o_v\le e^{-41},\qquad\sup_v\log v\cdot o_v\le e^{-120},
  \qquad E_{\mathrm{BR}}\le e^{-76}.
\end{equation}
\begin{proof}
\begin{itemize}
\item Near $u=X_0$, $\Tc$ is continuous: it is $0$ below $X_0$ and $1-\rhodet-\jz$ above, which vanishes at $X_0$.
\item $o_{\mathrm{det}}\le(2.04H(u)+2)X_0^{5/12}/u$, from the linear window of length $\kappa u$ in \eqref{eq:rhodet} and the
  unit spacing of the rigid set.
\item On $\Good(<u)$, the $\Abs$-part of $\Tc$ is $\sigma_{\Abs}=K_{\Abs}\rhot-\rho_{\Abs}$ (as $\bar\rho_{\Abs}=\rho_{\Abs}$,
  Remark~\ref{rem:stopped}). So $o_{\Abs}\le2\sup|\sigma_{\Abs}|\le4L_u^{-4}$, by Lemma~\ref{lem:rhoA}(c) (from parts (a), (b)
  and Lemma~\ref{lem:lemmaP}). Lemma~\ref{lem:lemmaP} holds on $\Good(<t)$ at heights up to $t+1.01H(t)$, using only
  $E_{\mathrm{cl}}$ at lower heights and Chebyshev below $2y$. So reading $o_{\Abs}$ at the current height $t$ costs no event
  later than $t$. Moreover $o_{\Abs}$ vanishes unless
  $\log\log u>2\log\Lambda+5.97$ ($=42.8$ at $\Lambda=10^8$), since $\Abs$ is empty below that height.
\item $o_R\le\sum_{k:\,x_k\le u}\gamma_k\kappa_kH(u)u^{-1}(x_k/u)^{1-\theta}\le H(u)u^{-\theta}$, since
  $\gamma_k\kappa_k\le x_k^{1-\theta}$ and $u\ge x_1\ge\exp(X_0^2)$.
\item Hence $\sum_u\log(2u)\,o_u\le4.2H(X_0)X_0^{-7/12}\log X_0+5e^{-2\cdot42.8}\le e^{-83}$, and
  $E_{\mathrm{BR}}\le883e^{-83}+8.4C_NE_\theta\le e^{-76}$. Also $\sum_v(\log3v)^2o_v\le5.8e^{-42.8}\le e^{-41}$ and
  $\sup_v\log v\cdot o_v\le4e^{-3\cdot42.8}\le e^{-120}$. All of these improve with $\Lambda$.
\end{itemize}
No triangular windows and no Lipschitz bound are needed.
\end{proof}

\begin{lemma}[the noise terms, in predictable form]\label{lem:predictable}
% v3: Lemma 3.15
With \eqref{eq:decomp},
\begin{equation}\label{eq:predictable}
  M(x)-(B\Rres T')(x)=(\vartheta/H)(x)-(B\Kpt B\Rres F_r)(x)-(B\Rres G)(x),
\end{equation}
and, on $\Good(<x)$,
\begin{multline*}
  \log x\cdot|M-B\Rres T'|(x)\le\eps_N:=5.1C_N\sum_u\log(2u)\bar F_r(u)+15.3C_NC_S'\sum_v(\log3v)^2F_r^+(v)\\
  +\sup_x\frac{6\log x}{H(x)}\le10^{-13},
\end{multline*}
where $\bar F_r(v):=o_v+6/H(v)+\bar\nu(v)+54\eta(v)$ bounds $|F_r|$ on $[v,2v)$, and $F_r^+:=2\bar F_r$ bounds $|(1-B)F_c|$.
Every sum is over dyadic scales $\ge X_0/2$; there is no dependence on $x$, and no event at the current scale is used.
\end{lemma}

\begin{proof}
\emph{The identity.} By \eqref{eq:pushthrough} for $F_c$ and $\Kpt B\Rres=1-\Rres$, $\Kpt\rhot_c=F_c-\Rres F_c+\Rres G$, where
$\rhot_c:=\rhot-(1+B\Kpt)^{-1}BJ$. So $M=B\Rres J+B\Tc+\vartheta/H-BF_c+B\Rres F_c-B\Rres G$. Insert $F_c=\Tc+F_r$ and
$\Rres\Tc+\Rres J=\Rres T'$. In particular, the current-scale terms $-B(\nu+\NT)(x)$ of $BF_c$ cancel exactly against those of
$B\Rres F_c$ (its $n=1$ term). Every term on the right of \eqref{eq:predictable} involves $\nu$, $\NT$ and $\rhot$ only at scales
$\le x/2+2H$.

\emph{The bounds.}
\begin{itemize}
\item $\log x\cdot|(B\Kpt B\Rres F_r)(x)|\le5.1C_N\sum_{u\le x/2+2H}\log(2u)\bar F_r(u)$. The kernels of $\Kpt B\Rres=1-\Rres$
  are the $C_n$ with $n\ge2$, so Proposition~\ref{prop:true-kernels} applies ($5.1C_N$ bounds both kinds of blocks).
\item $\log x\cdot|(B\Rres G)(x)|\le15.3C_NC_S'\sum_v(\log3v)^2F_r^+(v)$, since
  $|G(y)|\le3\sum_v\omega_y(v)\sup_{[v,3v]}F_r^+$ (Lemma~\ref{lem:S}(iii)).
\item $|\vartheta/H|\le6/H(x)$, and $\sup\rhot\le1.01$ surely (\S\ref{sec:noise-perscale}), so $|\NT|\le2.02\eta\le54\eta$ by
  \eqref{eq:NT}.
\end{itemize}
By \eqref{eq:sw-sums} (note that $\sum_u\log(2u)\bar\nu(u)$ and $\sum_v(\log 3v)^2\bar\nu(v)$ are dominated by the sums
there), Lemma~\ref{lem:tau-near} ($E_\eta^0$ and $E_\eta^{(2)}$) and \eqref{eq:osc}, $\eps_N\le10^{-13}$. The largest part is
$30.6C_NC_S'e^{-41}$, from the oscillation of $\sigma_{\Abs}$.
\end{proof}

\paragraph{Design facts used} (\S\ref{sec:construction}).
\begin{itemize}
\item $R=\sum_kR_k$, with $R_k\equiv0$ on $[1,x_k+2H(x_k)]$, $x_1\ge\exp(X_0^2)$, $x_{k+1}\ge\exp(x_k)$, and
  $\gamma_k\in[\Gamma_k,2\Gamma_k]$ with $\Gamma_k=\kappa_0x_k^{1-\theta}/(2\log^3x_k)$; so $\log\gamma_k\asymp(1-\theta)\log x_k$.
\item \eqref{eq:Rkprime}: $|R_k'(v)|\le\kappa_k(v/x_k)^{\theta-1}$ with $\kappa_k=2\kappa_k'/\log x_k$, $\kappa_k'=\kappa_0/\log^2x_k$,
  $\kappa_0\le\frac14$.
\end{itemize}

\paragraph{A squarefree average.} $F_\theta:=\sup_yy^{\theta-1}\sum_{n\in\NN,\,n\le y}\mu^2(n)n^{-\theta}\le1.33$ for $\theta\le\frac12$
(in particular for $\theta\in[\frac14,\theta_1]$):
\begin{itemize}
\item $F_\theta\le F_{1/2}$, since $y^{\theta-1}n^{-\theta}=y^{-1}(y/n)^\theta$ increases with $\theta$;
% Logged correction 54: 1.32964 (an upper bound) replaces 1.3296; the maximum is 1.3296349 (independent check R8, C2).
\item a direct computation of $F_{1/2}$ on $\N$ for $y\le10^6$ gives the maximum $1.32964$, at $y=7$ (see
  Appendix~\ref{app:cert}, \texttt{finite-n});
\item for $y\ge10^6$, the elementary bound $|Q(u)-6u/\pi^2|\le3\sqrt u$ for the squarefree count $Q$ gives
  $y^{\theta-1}\sum_{n\le y}\mu^2(n)n^{-\theta}\le(6/\pi^2)/(1-\theta)+3y^{-1/2}+1.5y^{-1/2}\log y\le1.25$; and
  \textup{(H$_{\mathrm{BS}}$)} changes the count for $\NN$ by a relative $O(\eta_1)$.
\end{itemize}
% Logged correction 9 (CORRECTIONS-vs-v3.md): explicit bound |Q(u) - 6u/pi^2| <= 3 sqrt u replaces v3's (1+o(1)).

\begin{proposition}[positivity of the sharp reference]\label{prop:g1}
% v3: Proposition 3.16
On $\Good(<x)$,
% Logged correction 47 (CORRECTIONS-vs-v3.md): v3 states 0.97 - o(1) for all x >= X_0 + H(X_0); the edge term rho_edge of
% Lemma margin(d) lowers it to 0.82 on [X_0 + H, X_0^{1.01}].
\[
  \rhosh(x)\cdot\log x\ge
  \begin{cases}
    0.82-o(1), & X_0+H(X_0)\le x\le X_0^{1.01},\\
    0.97-o(1), & x\ge X_0^{1.01}/2.
  \end{cases}
\]
\end{proposition}

The cutoff $\Xm$ in $\Kpt$ does not affect $\rhosh$: $c_n(y)=\mu(n)\chi^t_y(n)\mathbf 1[y/n\ge\Xm]$ (Lemma~\ref{lem:lemmaM}),
and $T'=0$ on $[\Xm,X_0)$. The proof has two ranges. Put $\log X_{\mathrm{sw}}:=X_0^{1/100}$. For every $\Lambda\ge10^8$,
$X_0^{1.01}\le X_{\mathrm{sw}}$, $\log X_{\mathrm{sw}}\le(0.0048/\kappa)^{1/5}$ (the limit for $\mu_0$ below, which is
$\ge X_0^{0.079}$), and $\log X_{\mathrm{sw}}=e^{\Lambda/100}\ge10^5\Lambda^2\log^2\Lambda=\log X_{\mathrm{LP}}$.
% Logged correction 69 (CORRECTIONS-vs-v3.md): X_sw >= X_LP made explicit (was ">>").

\begin{proof}[Proof for range \textup{(A)}: $X_0+H\le x\le X_{\mathrm{sw}}$]
No designed oscillation is active, since $x_1\ge\exp(X_0^2)>X_{\mathrm{sw}}$. So $T'=1-\rhodet+\sigma_{\Abs}$ there.
Lemma~\ref{lem:margin} (\eqref{eq:ratid}, \eqref{eq:shdecomp}) gives
$(1+\Kpt)^{-1}(1-\rhodet)=\mathrm{MA}[\mathbb P]+\rho_{\mathrm{edge}}+(1+\Kpt)^{-1}(\mu_0+\mu_1)$, where:
% Logged correction 47: support convention for MA[P] (primes >= X_0 only), the edge term moved out of mu_0 into its own
% item (d), and the conclusion stated piecewise.
\begingroup\emergencystretch=3em
\begin{enumerate}[label=(\alph*)]
\item $\mathrm{MA}[\mathbb P](q):=\mathrm{MA}_{\kappa q}[\text{rational primes}\ge X_0](q)$ is the density of the new primes (the
  rational primes $\ge X_0$) in the linear window at $q$ (\S\ref{sec:baserange}); it vanishes for $(1+\kappa)q\le X_0$. By
  \eqref{eq:huxley} and Lemma~\ref{lem:margin}(a) (Huxley and Heath-Brown, $X_0\ge X_{\mathrm{Hux}}$),
  $\mathrm{MA}[\mathbb P](x)\log x\ge0.99$ for every $x\ge X_0$;
\item $|\mu_0(u)|\le2\kappa(\log u)^3$ for every $u\ge X_0$, in particular on range (A)
  (Lemma~\ref{lem:margin}(b)). This smooth part contains the drift $+(\kappa/2)\Psi$, the partial windows, the decision-order
  mismatch (bounded by a direct count, $\le2.02(H(X_0)+1)(\log u+1)X_0^{-1}\sup|f_0|$; ``the rest'' in the proof of
  Lemma~\ref{lem:margin}(b)), and the $O(1/(\kappa x))$ term of \eqref{eq:ratid};
% Logged correction 70: the mismatch is described as in correction 53 (was "the order mismatch H(q)/(kappa q)"; pass 20, N40).
\item $\sum_{u\text{ dyadic}\ge X_0}\log(2u)\sup_{[u,2u]}|\mu_1|\le7\cdot10^{-6}$ \eqref{eq:mu1sum};
\item $\rho_{\mathrm{edge}}:=(1+\Kpt)^{-1}\mu_{\mathrm{edge}}$ is the edge at $X_0$, where $\mu_{\mathrm{edge}}$ is the part of
  $K_{\mathbb Q}\mathrm{MA}[\mathbb P]$ that lies below $X_0$ \eqref{eq:shdecomp}. By Lemma~\ref{lem:margin}(d),
  $\rho_{\mathrm{edge}}(x)\log x\ge-0.167$ for $X_0\le x\le X_0^{1.01}$, and
  $|\rho_{\mathrm{edge}}(x)|\log x\le3X_0^{-0.01}+X_0^{-0.34}$ for $X_0^{1.01}/2\le x\le X_{\mathrm{sw}}$.
\end{enumerate}
\endgroup
Now bound the responses.
\begin{itemize}
\item $\mu_0$: by Lemma~\ref{lem:lemmaM} ($|c_n|\le1$ and $\sum|c_n|/n\le\log x+B_1$, with $B_1=1$),
  \[
    \log x\cdot|(1+\Kpt)^{-1}\mu_0(x)|\le\log x\,(\log x+1)\sup_{u\le x}|\mu_0(u)|\le2.1\,\kappa(\log x)^5\le X_0^{-0.34}
  \]
  for $\log x\le X_0^{1/100}$, since $\kappa=X_0^{-5/12+2\eps_\kappa}=X_0^{-0.39666\ldots}\le X_0^{-0.3966}$.
% Logged correction 69 (R13, C2): "kappa <= X_0^{-0.3967}" was false as written; 2.1 X_0^{-0.3966} X_0^{0.05} <= X_0^{-0.34} still.
\item $\mu_1$: by Corollary~\ref{cor:saias}, $\log x\cdot|(1+\Kpt)^{-1}\mu_1(x)|\le C_S\cdot7\cdot10^{-6}\le6\cdot10^{-4}$.
\item $\sigma_{\Abs}$: likewise, at most $C_S\sum_v\log(2v)\cdot2L_v^{-4}\mathbf 1[\log\log v>2\log\Lambda+5.97]\le e^{-80}$.
\end{itemize}
So, pointwise, for $X_0+H\le x\le X_0^{1.01}$,
\[
  \rhosh(x)\log x\ge0.99-0.167-6\cdot10^{-4}-X_0^{-0.34}-e^{-80}\ge0.822,
\]
and for $X_0^{1.01}/2\le x\le X_{\mathrm{sw}}$,
\[
  \rhosh(x)\log x\ge0.99-6\cdot10^{-4}-2X_0^{-0.34}-3X_0^{-0.01}-e^{-80}\ge0.989\ge0.97,
\]
for every $X_0\ge\max(e^{10^8},X_{\mathrm{Hux}})$; no constant here depends on $\Lambda$. By Lemma~\ref{lem:margin}(d)(i), the loss
$0.167$ occurs only on the windows $N\cdot(X_0/(1+\kappa),X_0)$ with $\mu(N)=1$, each of relative length $\kappa$.
\end{proof}

\begin{proof}[Proof for range \textup{(B)}: $x\ge X_{\mathrm{sw}}$]
Call $\gamma_j$ \emph{resolvable at $x$} if $\gamma_j\le\Gamma(x):=e^{c_{\mathrm{LP}}\sqrt{\log x}/4}$, and let $\gamma_*$ be the largest
resolvable $\gamma_j$ ($\gamma_*:=1$ if there is none). Put
\[
  \delta:=\min\big(\gamma_*^{-2}\log^{-4}x,\ X_0^{-1/2}\big),\qquad J:=[x(1-2\delta),\,x(1-\delta)].
\]
% Logged correction 69: the two comparisons below made explicit (were "delta x >> H" and "c_LP sqrt(x_j) >> log x_j").
The window $J$ ends $\delta x\ge x^{0.98}>H(x)$ before $x$ (as $1/\delta\le x^{0.02}$, Step~3), so the proof uses only the realised
data before $x-H$, on $\Good(<x)$. In particular Theorem~\ref{thm:lp} is applied with $y:=x(1-\delta)\le x-H(x)$, below the
certified horizon of $\Good(<x)$, where (LP-a) holds (Remark~\ref{rem:horizon}).
% Logged correction 73: the local prime number theorem is applied below the certified horizon of Good(<x).
% Logged correction 84 (R21, 84-1): k is defined here, before its first use, and only for x >= x_1.
For $x\ge x_1$ let $k:=\max\{j:x_j\le x\}\ge1$; for $x<x_1$ no $R_j$ is active on $[1,x]$ (Step~4). All old oscillations $j<k$ are resolvable: $\log\gamma_j\le(1-\theta)\log x_j\le\log x_j$ (as
$\gamma_j\le2\Gamma_j\le x_j^{1-\theta}$), while $\log x\ge x_{k-1}\ge x_j$ and $x_j\ge\exp(X_0^2)$, so
$c_{\mathrm{LP}}\sqrt{\log x}/4\ge(0.0044/\Lambda)\sqrt{x_j}\ge\log x_j$.

\emph{Step 1: the window average of $\rhot$.}
\begin{itemize}
\item $\int_J\rhot=\#\{q:s_q\in J\}\pm(H(x)+2)$, the error coming from the partial boxes at the two ends of $J$; and slots lie
  within $H/2$ of values.
% Logged correction 31 (CORRECTIONS-vs-v3.md): the same partial-box effect; v3 has +-2. Absorbed by the +-27H(x) below.
  By the local count, and since the prime powers
  $p^k\in J$ with $k\ge2$ number at most $2\sqrt x$, $\int_J\rhot=\Pi(J)\pm(27H(x)+2K_{\mathrm{loc}}+6+2\sqrt x)$.
\item By Theorem~\ref{thm:lp}(ii), $\Pi(J)=\delta x\,\pi_0(x)(1+\eps_J)+O(Kxe^{-c_{\mathrm{LP}}\sqrt{\log x}})$ with
  $0\le\eps_J\le2.01\,\delta/\log x$, where $\pi_0(x):=(1-1/x)/\log x$ (which decreases, and $J\subset[x(1-2\delta),x]$).
\item Since $1/\delta\le\Gamma^2\log^4x+X_0^{1/2}$ with $\Gamma^2\le e^{c_{\mathrm{LP}}\sqrt{\log x}/2}$, the error term is at most
  \[
    \pi_0(x)\bigl(2K\log^5x\,e^{-c_{\mathrm{LP}}\sqrt{\log x}/2}+2K\log x\,X_0^{1/2}e^{-c_{\mathrm{LP}}\sqrt{\log x}}\bigr);
  \]
  and
  $c_{\mathrm{LP}}\sqrt{\log x}\ge2(\log K+\Lambda+10\log\sqrt{\log x})+10$, as $\log K\le98+5.9\log\Lambda\le\Lambda$ and
  $(0.0176/\Lambda)\sqrt{\log x}\ge4\Lambda+20\log\sqrt{\log x}+10$ for $\sqrt{\log x}\ge X_0^{0.005}$. So the error is at most
  $\pi_0(x)\,e^{-\Lambda}$, and $\mathrm{avg}_J\rhot=\pi_0(x)+o(1/\log x)$.
% Logged correction 70 (check R13, C4): the condition is stated against both error terms (c_LP sqrt(log x)/2 must beat
% log(2K) + 5 log log x); the earlier form covered only log log x <= 0.6 Lambda.
% Logged correction 70 (CORRECTIONS-vs-v3.md): the comparison with log K made explicit (was "recall log K = O(Lambda)").
\end{itemize}

\emph{Step 2: $\rhosh$ against $\rhot$ on average.}
\begin{itemize}
\item $\rhot-\rhosh=[\rhot+\nu+\NT-B\Rres T']-\nu-\NT+B[\Rres T'-\rhosh]+(B\rhosh-\rhosh)$.
\item Lemma~\ref{lem:predictable} and Proposition~\ref{prop:br} bound the first and fourth terms by
  $(\eps_N+E_{\mathrm{BR}}+j_{X_0}(x))/\log x$. In range (B), $\log(x/X_0)\ge X_0^{1/100}-\Lambda\ge600\Lambda^4$, so $x$ lies in
  the huge range of Proposition~\ref{prop:xb-ranges}, where $j_{X_0}\le e^{-59}$. Also $|\nu|\le\bar\nu$ and $|\NT|\le54\eta$.
\item Averaged over $J$, $\int_J(Bf-f)$ is a boundary term $\le2H(x)\sup|f|$, and $\sup|\rhosh|\le2(\log x+B_1)$.
\item Hence $\mathrm{avg}_J\rhosh\ge\pi_0(x)-(\eps_N+E_{\mathrm{BR}}+e^{-59}+o(1))/\log x$.
\end{itemize}

\emph{Step 3: pointwise against average, for all components except the current oscillation.} For $y\in J$,
$\rhosh(x)-\rhosh(y)=\sum_n[c_n(x)T'(x/n)-c_n(y)T'(y/n)]/n$.
\begin{itemize}
\item \emph{Threshold changes:} at most $2\cdot$(Lemma~\ref{lem:threshold-mass} at $2\delta$) $=o(1/\log x)$, since $|T'|\le2$,
  $\delta\le\log^{-4}x$ and $\log H(x)=8(\log\log x)^2$.
% Logged correction 59: the hypothesis 6H(x)x^{-1/3} <= 2 delta <= 1/4 of Lemma threshold-mass, checked here.
  The lemma applies at $2\delta$: $2\delta\le2X_0^{-1/2}\le\frac14$; and $1/\delta\le\Gamma(x)^2\log^4x+X_0^{1/2}\le x^{0.02}$ (as
  $\gamma_*\le\Gamma(x)$ and $\log x\ge X_0^{1/100}$), while $6H(x)x^{-1/3}\le x^{-0.3}$ by Theorem~\ref{thm:caplw}(c), so
  $2\delta\ge6H(x)x^{-1/3}$.
\item \emph{The jump of $T'$ at $X_0$:} only $n\in[y/X_0,x/X_0]$ see it. Put $N_*:=x/X_0$, $I:=[N_*(1-2\delta),N_*]$ and
  $L_I:=\log(N_*/2)$. By the lag bound,
% Logged correction 66: renamed from L_*, which clashes with L_* = 10^7 of eq:absorbed.
$|N-T|\le H(N_*)$ at both ends of $I$, and under the target \eqref{eq:target},
  $T(I)\le2.0001|I|+1+\#\Abs(I)$ (as in Corollary~\ref{cor:localcount}(i)). Cover $I$ by $\lceil|I|/h(N_*(1-2\delta))\rceil$ consecutive windows $[u,u+h(u)]$,
  $h(u)=u^{\theta/2}$; by Lemma~\ref{lem:lemmaP} each holds at most $h(u)L_u^{-4}$ absorbed elements, so
  $\#\Abs(I)\le(2.01\,\delta N_*+N_*^{\theta/2})L_I^{-4}$. Hence $\#(\NN\cap I)\le4.01\,\delta N_*+2H(N_*)+2+N_*^{\theta/2}$, and
  this part is at most $\jz\sum_{n\in\NN\cap I}1/n\le\jz\big(4.1\,\delta+3H(N_*)/N_*+2N_*^{\theta/2-1}\big)=o(1/\log x)$.
% Logged correction 14 (CORRECTIONS-vs-v3.md): v3 omits #A(I) here (same omission as correction 4); with T(I) <= 2|I|+2+#A(I)
% the constant 2.1 delta becomes 4.1 delta (noted under #14).
% Logged correction 93 (CORRECTIONS-vs-v3.md): for y in J the level points x/n and y/n differ by up to 2 delta x/n, not
% delta x/n; the smooth-change bounds for rho_det and for the resolvable R'_j are restated at displacement 2 delta.
\item \emph{$\rhodet$, smooth change:} for $y\in J$ and $x/n\ge X_0$ put $v:=x/n$ and $v':=\max(y/n,X_0)$, so that $|v-v'|\le2\delta v$
  (if $y/n<X_0$, then $T'(x/n)-T'(y/n)=\jz+\rhodet(X_0)-\rhodet(x/n)$, as $R'$ and $\sigma_{\Abs}$ vanish near $X_0$, and $\jz$ is
  the previous item); then
  $|\rhodet(v)-\rhodet(v')|\le\min\big(2(2\delta(1+\kappa)v+1)/(\kappa v),\ 3\Psi(2v,X_0)/(\kappa v)\big)$. Indeed the windows
  $[v,(1+\kappa)v]$ and $[v',(1+\kappa)v']$ overlap (as $2\delta<\kappa$) and differ by two intervals of length $\le2\delta(1+\kappa)v$,
  which enter the difference of the counts with opposite signs; so the counts differ by at most
  $2\delta(1+\kappa)v+1$, and the normalisations $\kappa v$, $\kappa v'$ change the quotient by at most $2.01\delta(1+1/(\kappa v))$,
  which is smaller; and $\rhodet(v)+\rhodet(v')\le\Psi(2v,X_0)/(\kappa v)+\Psi(2v,X_0)/(\kappa v')\le3\Psi(2v,X_0)/(\kappa v)$, as
  $(1+\kappa)v'\le2v$. Put $u:=\log v/\Lambda$
  and $U_0:=\Lambda$. Use the first bound for $u<U_0$ and the second beyond; for $u\ge U_0$, Rankin's bound (exponent
  $1-1/(2\Lambda)$) gives $\Psi(2v,X_0)\le10v\Lambda^{1.65}e^{-u/2}$. With the weights of Corollary~\ref{cor:saias}
  ($\sum_{v<X_0^{U_0}}\log(2v)\le1.5(U_0\Lambda)^2$) and $1/\kappa\le X_0^{5/12}$, $\log x$ times this part of $\rhosh(x)-\rhosh(y)$
  is at most
  \begin{multline*}
    C_S\Big[1.5\Lambda^4(4.0001\,\delta+2/X_0)/\kappa+\sum_{u\ge\Lambda}\log(2v)\cdot30\Lambda^{1.65}e^{-u/2}/\kappa\Big]\\
    \le C_S\big[7\Lambda^4X_0^{-1/12}+90\Lambda^{4.7}e^{-\Lambda/12}\big]\le e^{-10^6},
  \end{multline*}
  using $\delta\le X_0^{-1/2}$. Both terms decrease with $\Lambda$, so this holds for every $\Lambda\ge10^8$; at $\Lambda=10^8$ (with
  $C_S\le77$) the logarithm of the bound is $\le-8.33\cdot10^6$.
\item \emph{Resolvable $R_j'$, smooth change:} by \eqref{eq:Rk} and \eqref{eq:ampl},
  $|R_j''(w)|\le\alpha_j\gamma_j^2w^{\theta-2}(1+O(1/\gamma_j))$, so at a level with $|v-v'|\le2\delta v$,
  $|R_j'(v)-R_j'(v')|\le1.01\,\delta\gamma_j\kappa_j(v/x_j)^{\theta-1}$. Summed with the weights $\mu^2(n)/n$ over $n\le x/x_j$, this is at
  most $1.01F_\theta\,\delta\gamma_j\kappa_j\le2\delta\gamma_j\kappa_j/(1-\theta)\le2\kappa_j/((1-\theta)\gamma_j\log^4x)$, as $F_\theta\le1.33$
  and $\delta\gamma_j^2\le\log^{-4}x$ for resolvable $\gamma_j\le\gamma_*$. This is summable over $j$.
\item \emph{$\sigma_{\Abs}$:} $|\sigma_{\Abs}|\le2L_v^{-4}$, and it vanishes unless $\log\log v>2\log\Lambda+5.97$. So the
  weighted sum is at most $4C_S\sum_v\log(2v)L_v^{-4}/\log x\le C_Se^{-80}/\log x$.
\end{itemize}

% Logged correction 84 (CORRECTIONS-vs-v3.md): v3 defines k := max{j : x_j <= x}, which is undefined for x < x_1, a range
% the proposition covers ([X_sw, x_1) is non-empty); there R' = 0 on [1,x] and the step is void. The case k = 1 (gamma_* = 1) is stated.
\emph{Step 4: the current oscillation, if it is not resolvable.} If $x<x_1$, every $R_j$ vanishes on $[1,x]$ (as $R_j=0$ on
$[1,x_j^+]$), so $R'=0$ there, the oscillation contributes nothing, and this step is void; this happens on the whole range
$[X_{\mathrm{sw}},x_1)$, since $\log X_{\mathrm{sw}}=X_0^{0.01}<X_0^2\le\log x_1$. Otherwise $k\ge1$ is as above; $R_j$ vanishes on
$[1,x]$ for $j>k$. Put $\rho^{(k)}_{\mathrm{osc}}(y):=\sum_nc_n(y)R_k'(y/n)/n$.
\begin{itemize}
\item \emph{Separation.} If $\gamma_k$ is not resolvable, then $\gamma_k>\Gamma(x)$. Every $j<k$ is resolvable (and no $\gamma_j$ with $j>k$ is, as
  $\log\gamma_j\ge\log\Gamma_j\ge(1-\theta)\log x_j-\Lambda-3\log\log x_j-1>c_{\mathrm{LP}}\sqrt{\log x}/4$ for $x_j>x$), and
  $\gamma_*\le\max(1,\gamma_{k-1})\le\log x$ (for $k\ge2$, $\gamma_{k-1}\le x_{k-1}\le\log x$ as $x_k\ge\exp(x_{k-1})$; for $k=1$ there is
  no resolvable $\gamma_j$, and $\gamma_*=1$). So
  $\gamma_k\delta\ge\Gamma(x)\min(\log^{-6}x,X_0^{-1/2})\ge e^{c_{\mathrm{LP}}\sqrt{\log x}/5}$, since, for
  $\sqrt{\log x}\ge X_0^{0.005}=e^{\Lambda/200}$,
  \[
    c_{\mathrm{LP}}\sqrt{\log x}/20\ge(8.8\cdot10^{-4}/\Lambda)\sqrt{\log x}\ge\max(6\log\log x,\Lambda/2).
  \]
% Logged correction 70: the comparison of correction 69 displayed (it overflowed the line).
% Logged correction 69: the comparison made explicit (was "c_LP sqrt(log x) >> log X_0 + log log x").
\item \emph{Average.} Where $c_n(y)$ is constant on $J$, $\int_JR_k'(y/n)\,dy/n=R_k(x(1-\delta)/n)-R_k(x(1-2\delta)/n)$, and
  $|R_k(v)|\le2\kappa_kv(v/x_k)^{\theta-1}/((1-\theta)\gamma_k)$ (integrate the oscillation at frequency $\gamma_k$). So, with
  Lemma~\ref{lem:lemmaM}, the $J$-average is at most $4\kappa_k(\log x+B_1)/((1-\theta)\gamma_k\delta)$, plus the threshold
  part $\le\kappa_k\cdot R(2\delta)$, where $R(\cdot)$ is the bound of Lemma~\ref{lem:threshold-mass}. Both are $o(1/\log x)$.
\item \emph{Pointwise.} $|c_n|\le\mu^2(n)$, and $R_k'(x/n)\ne0$ forces $n\le x/x_k$. So
  \[
    |\rho^{(k)}_{\mathrm{osc}}(x)|\le\kappa_k(x/x_k)^{\theta-1}\sum_{n\le x/x_k}\mu^2(n)n^{-\theta}\le\kappa_kF_\theta.
  \]
  \begin{itemize}
  \item For $x\le x_k^2$: $\kappa_kF_\theta\log x\le4\kappa_k'F_\theta=4\kappa_0F_\theta/\log^2x_k=o(1)$.
  \item For $x_k^2<x<\exp(C'\log^2x_k)$, with $C':=(4(1-\theta)/c_{\mathrm{LP}})^2$:
    \[
      \kappa_kF_\theta\log x\le2\kappa_k'F_\theta C'\log x_k=2\kappa_0F_\theta C'/\log x_k=o(1),
    \]
    since $\log x_k\ge X_0^2$ and $C'\asymp\Lambda^2$.
  \end{itemize}
\item For $x\ge\exp(C'\log^2x_k)$, $\gamma_k\le x_k^{1-\theta}$ is resolvable, and Step~3 covers it.
\end{itemize}

\emph{Step 5.} Combining Steps 1--4,
$\rhosh(x)\log x\ge\pi_0(x)\log x-\eps_N-E_{\mathrm{BR}}-e^{-59}-o(1)\ge1-10^{-12}-o(1)\ge0.97$.
% Logged correction 69 (CORRECTIONS-vs-v3.md): the terms behind o(1) listed with explicit bounds; their sum is
% <= 1.6e-3 Lambda^{-4}, inside the convention o(1) <= Lambda^{-3} of Section sec:noise-notation.
Here $o(1)$ is the sum of the following contributions to $\rhosh(x)\log x$; we use $\log x\ge X_0^{1/100}=e^{\Lambda/100}$,
$1/\delta\le x^{0.02}$, $\delta\le\log^{-4}x$, $\log(1/\delta)\le0.6\sqrt{\log x}$ and $H(x)\le x^{0.01}$.
\begin{itemize}
\item Step~1: $\le x^{-0.4}$ from the partial boxes and the local count; $\le e^{-0.4c_{\mathrm{LP}}\sqrt{\log x}}$ from
  Theorem~\ref{thm:lp}(ii) (as $\log K\le98+5.9\log\Lambda$); and $\le3\delta/\log x$ from the variation of $\pi_0$ over $J$.
\item Step~2: $\le129\log^{-5}x$ from $\bar\nu$ (\S\ref{sec:noise-perscale}); $\le55e^{-0.99\Lambda}$ from $\NT$
  (Lemma~\ref{lem:tau-near}); and $\le x^{-0.9}$ from the boundary term.
% Logged correction 93: the rho_det and R'_j items are at displacement 2 delta; the bounds X_0^{-0.08} and log^{-2} x still hold.
\item Step~3: $\le160\log^{-1.5}x$ from the threshold changes; $\le\log^{-3}x$ from the jump at $X_0$; $\le X_0^{-0.08}$ from
  $\rhodet$ (the bound of Step~3, at displacement $2\delta$, is $\le C_S[7\Lambda^4X_0^{-1/12}+90\Lambda^{4.7}e^{-\Lambda/12}]\le X_0^{-0.08}$);
  $\le\log^{-2}x$ from the resolvable $R_j'$ (namely $\le4\sum_j\kappa_j\log^{-3}x$, as $1-\theta\ge\frac12$); and $\le4C_S\sum_{L_v>391\Lambda^2}\log(2v)L_v^{-4}\le1.5\cdot10^{-3}\Lambda^{-4}$
  from $\sigma_{\Abs}$.
\item Step~4: $\le e^{-c_{\mathrm{LP}}\sqrt{\log x}/6}+40\log^{-1.5}x$ from the average, and $\le2\cdot10^4\Lambda^2X_0^{-2}$ from the
  pointwise bounds (with $\kappa_0\le\frac14$, $F_\theta\le1.33$ and $C'\le(171\Lambda)^2$).
\end{itemize}
In total $o(1)\le1.5\cdot10^{-3}\Lambda^{-4}+e^{-0.01\Lambda}\le1.6\cdot10^{-3}\Lambda^{-4}$, for every $\Lambda\ge10^8$.
\end{proof}

\begin{remark}
No twisted M\"obius estimate is used: old oscillations are \emph{resolved} by the window of the local prime number theorem,
not cancelled. Range (A) uses no local prime number theorem, and range (B) uses no Huxley-type input beyond
$\jz\le1.0001/\Lambda$. The ranges fit: $\log X_{\mathrm{LP}}=10^5\Lambda^2\log^2\Lambda\le\log X_{\mathrm{sw}}=X_0^{1/100}\le X_0^{0.079}
\le(0.0048/\kappa)^{1/5}$.
% Logged correction 69: the ranges made explicit (were "<<").
\end{remark}

% 04g-margin.tex — v3 §3.7: Theorem 3.17, Proposition 3.18, Corollary 3.19. Author: gap-localpnt.
\subsection{The margin: window average, pointwise, and R1\texorpdfstring{$'$}{'}}\label{sec:noise-margin}
% v3: §3.7

\begin{theorem}[the margin on $H$-windows, uniformly in $x$]\label{thm:g2}
% v3: Theorem 3.17
On $\Good(<x)$, for every $x\ge2X_0$, the compensator margin \eqref{eq:M} satisfies
\[
  \log x\cdot M(x)=\log x\cdot(\tilde T'-B\Kpt\rhot)(x)\ge0.97-o(1)-\eps_N-E_{\mathrm{BR}}-j_{X_0}(x),\qquad
  \eps_N+E_{\mathrm{BR}}\le10^{-13},
\]
where $j_{X_0}$ is defined in \eqref{eq:jX0}; $j_{X_0}\le0.4209$ for $x\le X_0e^{0.01\Lambda}$ and $j_{X_0}\le e^{-0.0009\Lambda}$ beyond (Proposition~\ref{prop:xb-ranges}).
\end{theorem}

\begin{proof}
\begin{itemize}
\item By \eqref{eq:predictable} and Lemma~\ref{lem:predictable}, $\log x\cdot M(x)\ge\log x\cdot(B\Rres T')(x)-\eps_N$.
\item By Proposition~\ref{prop:br}, Lemma~\ref{lem:xb} and \eqref{eq:jX0} at each point $y$ of $W_x$,
  $\Rres T'(y)\ge\rhosh(y)-(E_{\mathrm{BR}}+j_{X_0}(x))/\log y$. (For $y<2\Xm$, $\Rres T'=\rhosh=T'$.)
\item By Proposition~\ref{prop:g1} at each point of $W_x$, $B\rhosh(x)\cdot\log x\ge0.97-o(1)$.
\item Every noise input lives at scales $\le x/2+2H$. Proposition~\ref{prop:g1} at points of $W_x$ uses only data up to
  $x-H/2$ (its windows $J$ end before $x$).\qedhere
\end{itemize}
\end{proof}

This is the place where the divergence of the sup-norm noise response is avoided. A bound of the form
$|Ke|\le(1+\log x)^{2K+1}\sup|\nu|$ diverges, because $\nu$ is of size $\asymp H(X_0)^{-1/2}$ at $u\approx X_0$. Here the noise
enters only through \eqref{eq:sw-sums}, whose weights $\log(2u)$ and $(\log3v)^2$ come from Corollary~\ref{cor:saias} and
Lemma~\ref{lem:S}. The exponent of $1/\log x$ in Theorem~\ref{thm:bs} is exactly~$1$, which is what makes these sums
independent of $x$.

\begin{proposition}[the pointwise margin]\label{prop:pointwise}
% v3: Proposition 3.18
Let $m(t)=T'(t)-(\Kpt\rhot)(t)-\NT_{\mathrm{pt}}(t)$ be the pointwise margin density (\S\ref{sec:noise-identity}), so that
$M=Bm+B\,\NT_{\mathrm{pt}}+\vartheta/H$. Put $v(t):=\sup_{|h|\le1.01H(t)}|m(t+h)-m(t)|$, $N:=t/X_0$, and
\begin{gather*}
  E_S:=\sup_v\log v\cdot o_v+2\sup_y\log y\cdot\eta(y)+27\sum_{v\text{ dyadic}}\log(3v)\,\hat\nu(v)\qquad(\le10^{-28}),\\
  V_{X_0}(t):=1.02\,(1+\log N/\Lambda)\sum_{M\in I'_t}\Hord(M)\,M^{\eps'-2},\qquad I'_t:=\{M\in\NN:|t/M-X_0|\le5.5H(X_0)\},
\end{gather*}
with $\eps':=\eps+0.015\le0.045$ and $\Hord(M)$ the number of ordered factorisations of $M$ in $\NN$. On $\Good(<t)$:
\begin{enumerate}[label=\textup{(\alph*)}]
\item for $t\ge X_0^{1.01}$: $v(t)\log t\le6.5E_S+V_{X_0}(t)$ with $V_{X_0}(t)\le e^{-0.0015\Lambda}$, and
  \[
    m(t)\log t\ge0.97-o(1)-\eps_N-E_{\mathrm{BR}}-54\sup_y\log y\cdot\eta(y)-6.5E_S-j_{X_0}(t)-V_{X_0}(t)\ge0.96-o(1);
  \]
\item for $2X_0\le t\le X_0^{1.01}$ the jump is paid once:
  $m(t)\log t\ge\rhosh(t)\log t-j_{X_0}(t)-10^{-12}\ge0.97-0.4209-10^{-12}\ge0.549$;
\item for $X_0\le t<2X_0$: $m(t)\log t\ge0.73$, surely (no event is used). Below $X_0$ there are no slots.
\end{enumerate}
\end{proposition}

\begin{proof}
\emph{Step 1: short windows.} For a slot window $I\subset[u,2u]$ with $|I|\le H(u)$, Lemma~\ref{lem:noise} gives
$\sigma(I)=\int_Im+E(I)$, with $|E(I)|\le\hat\nu(u)H(u)$ on $G_u^{\mathrm{sw}}$. Write $\eps_s(z):=\sup|E(I)|$ over windows with
$|I|\le s$ at $z$.

\emph{Step 2: the sub-$H$ increments of $\rhot$ are strip differences of $m$.} For $0\le s\le1.01H(z)$, with
$I_\pm:=[z\pm H/2,\,z\pm H/2+s)$,
\[
  \rhot(z+s)-\rhot(z)=\frac{\sigma(I_+)-\sigma(I_-)}{H(z)}=\frac1{H(z)}\int_0^s\big[m(z+H/2+r)-m(z-H/2+r)\big]dr+\frac{E(I_+)-E(I_-)}{H(z)} .
\]
This uses the centred boxes, and it holds also near $X_0$, where $\sigma=0$ and $m=T'=0$ below $X_0$. So
$|\Delta_s\rhot(z)|\le(s/H(z))\sup v+2\eps_s(z)/H(z)$.

\emph{Step 3: the recursion.} Write $\kappa_w:=\mathbf 1[\tau_w<t/w]$. Then
$\Delta_hm(t)=\Delta_hT'(t)-\sum_w(\kappa_w/w)\Delta_{h/w}\rhot(t/w)+(\text{filter switching})+\Delta_h\NT_{\mathrm{pt}}$.
\begin{itemize}
\item Filter switching and $\NT$ occur only for $w$ $\tau$-near at $t$, and cost at most $2\eta(t)$. The cutoff $\Xm$ switches
  nothing on, since $\rhot(\Xm)=0$.
\item $h/w\le1.01H(t/w)w^{\eps-1}$ by \eqref{eq:RH}. So
  \[
    v(t)\le S(t)+\sum_{w\ge2}1.01\,\kappa_w\,w^{\eps-2}\sup\{v(t'):|t'-t/w|\le2.1H(t/w)\},\qquad
    S(t):=o_t+2\eta(t)+2\sum_w\frac{\kappa_w}w\cdot\frac{\eps_{H(t)/w}(t/w)}{H(t/w)} .
  \]
\item Base case $t<2\Xm$: there $\Kpt\rhot=\NT_{\mathrm{pt}}=0$, so $m=T'=\Tc+J$ and
  $v(t)\le o_t+\jz\mathbf 1[|t-X_0|\le1.01H(t)]$.
\item $v(t)$ depends on $v$ and on the noise only at points $\le t/2+2.1H(t/2)<t$. So the recursion is a well-founded
  induction on $t$, and each short-window event is used only at strictly earlier times. This is the pointwise form of the
  event ordering in Definition~\ref{def:good} (event E3$'$).
\end{itemize}

\emph{Step 4: $\log t\cdot S(t)\le E_S$.} By Lemma~\ref{lem:noise}, $\eps_{H(t)/w}(t/w)/H(t/w)\le2\hat\nu(t/w)$, so the last
term of $S$ is at most $4\sum_{v'}\omega_t(v')\hat\nu(v')$.
\begin{itemize}
\item For $v'\ge\sqrt t/2$ use $\omega\le1.2$ and $\log t\le2\log(2v')$.
\item For $v'<\sqrt t/2$ use Lemma~\ref{lem:S}(ii): $\log t\cdot\omega_t(v')\le2.1\max_u(1+u)\rho(u)\cdot\log3v'\le4.3\log3v'$.
\item With $\log t\cdot o_t\le\sup_v\log v\cdot o_v$ and $\log t\cdot2\eta(t)\le2\sup_y\log y\cdot\eta(y)$, this gives
  $\log t\cdot S(t)\le E_S$.
\end{itemize}

\emph{Step 5: solving the recursion.} Iterating Step~3, a chain with product $M$ has weight at most $\Hord(M)M^{\eps'-2}$, and
it ends at a base point within $\sum_i2.1H(t_i)(t_k/t_i)\le4.4H(X_0)$ of $t/M$. So
$v(t)\le\sum_M\Hord(M)M^{\eps'-2}E_S/\log(t/M)+v_{X_0}(t)$, where $v_{X_0}$ collects the base terms
$\jz\mathbf 1[\cdots]$.
\begin{itemize}
\item Split at $M=\sqrt t$. The part $M\le\sqrt t$ is at most $2(1+A_0(0.045))E_S/\log t=6.44E_S/\log t$, since
  $A_0(\eps'):=\sum_{M\ge2}\Hord(M)M^{\eps'-2}$ has $A_0(0.045)=2.218$. The tail is at most
  $27t^{-0.1}(\log t/\Lambda)E_S/\log t\le X_0^{-0.09}E_S/\log t$, by Rankin's trick at $s=1.75$ (as $1/(2-\zeta(1.75))=26.5$).
\item \emph{The jump.} The base term $\jz\mathbf 1[\cdots]$ is reached only by chains with $M\in I'_t$. So
  $v_{X_0}(t)\le\jz\sum_{M\in I'_t}\Hord(M)M^{\eps'-2}$, and $\jz\log t\le1.01(1+\log N/\Lambda)$. For $N\ge e^{0.01\Lambda}$, every
  $M\in I'_t$ has $M\ge N/2$, so
  $\sum_{M\in I'_t}\Hord(M)M^{\eps'-2}\le(N/2)^{-0.155}\sum_M\Hord(M)M^{-1.8}\le8.5(N/2)^{-0.155}$, using
  $1/(2-\zeta_{\NN}(1.8))-1\le7.5(1+O(\eta_1))$. This gives $V_{X_0}(t)\le e^{-0.0015\Lambda}$.
\end{itemize}

\emph{Step 6(a).} $m\ge Bm-v$ pointwise, and $Bm=M-B\,\NT_{\mathrm{pt}}-\vartheta/H\ge M-54\eta-6/H$. Insert
Theorem~\ref{thm:g2}.

\emph{Step 6(b): the small range; the jump is paid once.}
\begin{itemize}
\item For $y\in W_t$, every $n$ with $c_n(y)\ne0$ or $C_n^y\ne0$ has $n\le y/\Xm<X_0$, so $n$ is a rigid integer. Hence
  $\tau_w=\Pplus(w)<y/w$ always: the $\tau$-filter of $\Kpt$ is void, $\NT=\NT_{\mathrm{pt}}=0$, and there are no $\tau$-near steps.
\item \emph{Split.} $m=m_c+m_j$ and $\rhot=\rhot_c+B\Rres J$, with $m_j:=\Rres J$ exactly (since the $J$-part of $\rhot$ is
  $B\Rres J$, and $1-\Kpt B\Rres=\Rres$).
\item \emph{The jump part.} By Lemma~\ref{lem:xb}, pointwise, $m_j(t)\ge\jz\rho^{\mathrm{sh}}_1(t)-\jz\mathfrak B(t)$; no step is $\tau$-near here, so the $\tau$-term of \eqref{eq:xb} is absent. Here
  $\rho^{\mathrm{sh}}_1=(1+\Kpt)^{-1}\mathbf 1_{[X_0,\infty)}$ and $\mathfrak B$ is defined there.
\item \emph{The smooth part.} Steps 1--5 apply verbatim to $(m_c,\rhot_c)$, because $\Delta_s(Bf)$ is exactly the strip
  difference of $f$. The base case is $m_c=\Tc$, which has no jump, so $v_c(t)\log t\le6.5E_S$. By \eqref{eq:predictable} and
  Proposition~\ref{prop:br}, $M_c:=M-B\Rres J$ satisfies $\log t\cdot M_c\ge\log t\cdot B\rhosh_c-\eps_N-E_{\mathrm{BR}}$, where
  $\rhosh_c:=\sum_nc_n\Tc(\cdot/n)/n$; and $Bm_c=M_c-\vartheta/H$.
\item \emph{$\rhosh_c$ is flat on $W_t$.} Here $c_n(y)=\mu(n)\mathbf 1[y/n\ge\Xm]$, and $|\Tc|\le o_{X_0}$ on
  $[\Xm,X_0+H]$. So $|\rhosh_c(y)-\rhosh_c(t)|\le\sum_uo_u+o_{X_0}\le11H(X_0)X_0^{-7/12}$, which is at most
  $e^{-10^7}/\log t$. (Here $o_R=0$, and $o_{\Abs}=0$ because $\Abs$ is empty while $\log\log n\le2\log\Lambda+5.97$.)
\item \emph{Sum.} $m(t)\ge\rhosh_c(t)+\jz\rho^{\mathrm{sh}}_1(t)-\jz\mathfrak B(t)-(\eps_N+E_{\mathrm{BR}}+6.5E_S+e^{-10^7})/\log t-6/H$,
  and $\rhosh_c+\jz\rho^{\mathrm{sh}}_1=\rhosh$. Then use $\rhosh(t)\log t\ge0.97$ (Proposition~\ref{prop:g1}(A); the range
  $2X_0\le t\le X_0^{1.01}$ lies in range (A) for every $\Lambda$).
\end{itemize}

\emph{Step 6(c): below $2X_0$.} Here $\Abs=\emptyset$, $R=0$, and $\Kpt\rhot(t)$ can be nonzero only through $w$ with
$t/w\ge\Xm$, i.e.\ $w=2$.
\begin{itemize}
\item For $X_0\le t<2\Xm$: $\Kpt\rhot(t)=\NT_{\mathrm{pt}}(t)=0$, so $m(t)=T'(t)=1-\rhodet(t)$. The integers in $[t,(1+\kappa)t]$
  that are not $X_0$-smooth include every prime there, so
  $1-\rhodet(t)\ge\mathrm{MA}[\mathbb P](t)-1/(\kappa t)\ge(0.99-10^{-9})/\log t$, by \eqref{eq:huxley} and
  Lemma~\ref{lem:margin}(a).
\item For $2\Xm\le t<2X_0$: $m(t)=T'(t)-\frac12\rhot(t/2)$, with $t/2\in[\Xm,X_0)$ (the filter $\tau_2=2<t/2$ is void, and
  $\NT_{\mathrm{pt}}=0$). Since slots start at $X_0$ and boxes are centred, $\rhot(t/2)$ counts the primes decided at slots in
  $[X_0,X_0+H/2]$, at most one per continuous slot. So
  $\rhot(t/2)\le(\int_{X_0}^{X_0+H/2}T'_{\mathrm{cont}}+3)/H(X_0)\le\frac12(\jz+o_{X_0})+3/H\le0.5001/\Lambda$, since
  $T'_{\mathrm{cont}}\le\jz+o_{X_0}$ there and $\jz\le1.0001/\Lambda$. With $T'(t)\ge(0.99-10^{-9})/\log t$ as above and
  $\log t\le\Lambda+0.7$,
  \[
    m(t)\log t\ge0.99-0.2501(1+0.7/\Lambda)-10^{-9}\ge0.739 .
  \]
\item These bounds use no event at all, only \eqref{eq:huxley} and the slot rule.\qedhere
\end{itemize}
\end{proof}

\begin{remark}
The contraction in Step~5 comes from compressing a shift $h$ at scale $t$ to a shift $h/w$ at scale $t/w$: a relative shift
$w^{\eps-1}$, times the weight $1/w$. In the small range the jump is compared pointwise with its sharp value
(Lemma~\ref{lem:xb}) and never enters a variation; bounding it as an additive variation term as well would pay for it twice.
\end{remark}

\begin{corollary}[R1$'$ on every window]\label{cor:r1}
% v3: Corollary 3.19
Let $I$ be any window: a slot window, or a physical interval $[a,b)$ with arbitrary (for instance integer, deterministic)
endpoints, of any length, anywhere in $[1,\infty)$. It may start below $X_0$ and straddle $2X_0$ or $X_0^{1.01}$. Let
$\ell_{\mathrm{jit}}(I):=\#\{\text{continuous slots in }I\}\le1.01|I\cap[X_0,\infty)|+1$, and let
$L\ge\log\max(I\cup\{X_0\})$. On $\Good(<\max I)$,
\[
  \ell_{\mathrm{jit}}(I)-\Comp(I)=\int_Im\pm2\ge0.549\,\frac{|I\cap[X_0,\infty)|}L-2\ge\frac{c'}L\,\ell_{\mathrm{jit}}(I)-K_d,
  \qquad\Comp^{(2)}(I)\le s_2\,|I|,
\]
with $c'=0.15$, $K_d=3$, and $s_2$ the supremum of $s_2(u)$ \eqref{eq:s2} over the dyadic scales meeting $I$.
\end{corollary}

\begin{proof}
\begin{itemize}
\item $\int_Im=\ell_{\mathrm{jit}}(I)-\Comp(I)\pm2$ is the bookkeeping of \S\ref{sec:noise-identity} without the $B$-average:
  $\#\{\text{continuous slots in }I\}=\int_IT'_{\mathrm{cont}}\pm2$ and $\Comp(I)=\int_I(\Kpt\rhot-K_{\Abs}\rhot+\NT_{\mathrm{pt}})$
  exactly, for any interval $I$. Below $X_0$ there are no slots, and below $2\Xm$ no jittered composites, so $m=0$ on
  $[1,X_0)$ and that part contributes nothing.
\item $\Comp$ is the compensator $\Comp_{\mathrm{jit}}$ of \S\ref{sec:queue} (latest prime dividing once). The $q^2$-part is
  the separate source $\Comp^{(2)}$, bounded by \eqref{eq:s2}.
\item $m\log t\ge0.549$ pointwise on all of $[X_0,\infty)$, by Proposition~\ref{prop:pointwise}(a)--(c): it is $\ge0.73$ on
  $[X_0,2X_0)$, $\ge0.549$ up to $X_0^{1.01}$, and $\ge0.959$ beyond. Since $0.549\ge1.01c'$, the $+1$ in the slot count only
  enters $K_d$.
\end{itemize}
In the form used in Theorem~\ref{thm:queue}: $\Comp(I)\le\ell_{\mathrm{jit}}(I)(1-c'/L)+K_d$ for every such $I$ with
$L\ge\log\max I$. Above $X_0^{1.01}$ the hole rate is at least $0.959/\log x$, independently of the block length $L_0$ and of
the tolerance $M_1\le H/L^{14}$.
\end{proof}

